% 本文件由 LaTeX 排版 Agent 根据有效 translation.json 自动生成。
% author=Augustine of Hippo; work_id=1202; title=论基督教教义

\chapter{\WorkTitle{论基督教教义}}

% structure: p0001 source=h1 target=omitted reason=page-title-already-rendered-by-parent

\Emph{卷册：前言 | I | II | III | IV}\par

\Emph{前言} 论述《论基督教教义》一书的效用。\par

\Emph{卷 I} 作者将他的著作分为两部分，一部分关乎圣经真义的发现，另一部分关乎圣经真义的表达。他表明，要发现意义，我们必须同时关注事物与符号，因为有必要知道我们应该向基督徒教导什么事物，以及这些事物的符号，即这些事物的知识应从何处寻求。在第一卷中，他论述事物，将其分为三类——应享受的事物、应使用的事物，以及使用并享受的事物。唯一应被享受的对象是\OriginalTerm{天主圣三}{Trinity}，他是我们的至高善和真正的幸福。我们的罪使我们无法享受天主；为使我们的罪被除去，\BibleQuote{圣言成了血肉}，我们的主受苦、死亡、复活，升天，以教会作为自己的新娘，我们在其中获得罪的赦免。如果我们的罪被赦免，我们的灵魂因\Emph{恩宠}而更新，我们可以怀着希望期待身体复活进入永恒的光荣；若不然，我们将被复活而遭受永恒的惩罚。这些关乎\Emph{信德}的事既已阐明，作者继而表明，除天主外，一切事物都是为了使用；因为虽然其中有些可以被爱，但我们的爱不应停留在它们身上，而应指向天主。我们自己并非天主享受的对象；他使用我们，却是为了我们自己的益处。作者继而表明，爱——为天主自身之爱天主，并为天主之爱邻人——是全部圣经的完成与目的。在补充了几句关于\Emph{望德}的话后，作者最后表明，\Emph{信德}、\Emph{望德}和\Emph{爱德}是那些想要正确理解并解释\Emph{圣经}的人所本质上必需的恩宠。\par

\Emph{卷 II} 作者在完成了对事物的阐述后，现在进而讨论符号的主题。他首先定义什么是符号，并表明符号有两类，自然的与约定的。在约定的符号（此处仅讨论此类）中，词语是最繁多和最重要的，也是圣经诠释者主要关注的符号。圣经的困难与晦涩主要源于两个来源，未知的符号和含糊的符号。本卷仅处理未知的符号，语言的含糊性则留待下一卷处理。因无知符号而产生的困难，可以通过学习圣经所用的希腊语和希伯来语，通过比较各种译本，以及通过关注上下文来消除。在解释比喻性表达时，对事物的知识与对词语的知识同样必要；异教的各种科学与技艺，只要它们真实且有用，可以被用来消除我们对符号的无知，无论这些符号是直接的还是比喻的。在揭露许多异教迷信和做法的愚昧与徒劳的同时，作者指出，他们科学与哲学中一切健全和有用的东西，都可以被用于基督徒的目的。最后，他表明我们应该以何种精神来着手研究和解释\Emph{圣经}。\par

\Emph{卷 III} 作者在前一卷中讨论了处理未知符号的方法，现在在第三卷中进而论述含糊的符号。此类符号可以是直接的或比喻的。在直接符号的情况下，含糊性可能源于标点、发音或词语的疑义，可以通过关注上下文、比较译本或参考原语言来解决。在比喻符号的情况下，我们需要防范两种错误：一、将直接表达解释为比喻；二、将比喻表达解释为直接。作者制定了一些规则，我们可以借此判断一个表达是直接的还是比喻的；一般规则是，任何在直接意义上与纯洁生活或正确教义不一致的，都必须被理解为比喻。作者继而制定了解释已被证明是比喻的表达的规则；一般原则是，任何不促进天主之人和人之爱的解释都不能是真实的。作者继而阐述并说明了多纳徒派提科尼乌斯的七条规则，他推荐给\Emph{圣经}的学生注意。\par

\WorkTitle{论基督教教义} \newline{} \Emph{第四部} \newline{} 转入其著作的第二部分，即关于表达的部分，作者首先表明，他无意撰写一篇关于修辞学法则的论文。这些法则可以在其他地方学到，且不应被忽视，因为它们对基督徒教师确实特别必要，他们应当在口才和言辞能力上出类拔萃。在极其细致地详述了演说家的各种品质后，他推荐圣经作者作为雄辩的最佳典范，在雄辩与智慧的结合上远胜于所有其他作者。他指出，明晰是文风最本质的品质，教师应当特别注重培养，因为它是教导的主要要求，尽管其他品质对于使听者愉悦和信服也是必要的。所有这些恩宠都应当从天主那里热切祈求，尽管我们不应忘记在学习上要热忱勤勉。他展示了有三种文风：\OriginalTerm{ subdued}{subdued}、\OriginalTerm{ elegant}{elegant} 和 \OriginalTerm{ majestic}{majestic}；第一种用于教导，第二种用于赞美，第三种用于劝诫：对于每一种，他都提供了从圣经和教会早期教师\Person{西彼廉}（Cyprian）和\Person{安博罗修}（Ambrose）那里选取的例证。他表明这些不同的文风可以混合，以及它们何时为何种目的而混合；并且它们都指向同一个目标，即把真理带给听者，使他能够理解它，愉快地聆听它，并在生活中实践它。最后，他劝勉基督徒教师本人，指出他所担任的职位的尊严与责任，要过一种与自己教导相和谐的生活，并为所有人树立良好的榜样。\par

\SourceRecord{Translated by James Shaw.；Nicene and Post-Nicene Fathers, First Series；Vol. 2.；Edited by Philip Schaff.；Buffalo, NY: Christian Literature Publishing Co.,；1887.；<http://www.newadvent.org/fathers/1202.htm>.}

% structure: p0001 source=h1 target=section reason=page-title

\section{论基督徒教义（序言）}

教导解释圣经的规则并非多余的任务\par

1. 有一些解释圣经的规则，我认为若教给热衷研读圣经的人，会大有益处，使他们不仅从他人揭示圣经奥秘的著作中获益，而且也能亲自将这些奥秘传授给别人。我计划将这些规则教导给那些有能力且愿意学习的人，只要上主天主在我写作时不离弃我，赐予我祂向来在默想这主题时所赐予我的思想。但在开始这项工作之前，我认为应先回应那些可能会反对这著作的人，或者那些若我不事先安抚他们就会反对的人。即使后来还有人提出反对，至少他们无法说服他人（他们本来可能影响那些人，若那些人未受预防而受其攻击），使他们放弃有益的研读，陷入无知中的迟钝懒惰。\par

2. 因此，有些人可能会反对我这著作，因为他们未能理解这里列出的规则。另一些人则会认为我白费了力气，因为他们虽然理解了规则，但在尝试应用这些规则来解释圣经时，却未能澄清他们所想澄清的问题；这些人，由于自己从这著作中未得到任何帮助，便会断言它对任何人都没有用处。还有第三类反对者，他们要么确实精通圣经，要么自以为精通，并且因为他们知道（或想象）自己在未阅读任何我将在此列出的这类指引的情况下，就已获得了某种解释圣书的能力，便会叫嚷说这样的规则对任何人都没有必要，而是认为，凡是为澄清圣经疑难所做的正确之事，最好完全依靠天主的恩宠来完成。\par

3. 简略回应这一切。对于那些不理解我所写内容的人，我的回答是：他们的不理解不能归咎于我。好比他们急于想看新月或残月，或某颗非常暗淡的星，而我用手指向它；如果他们的视力连我的手指都看不见，他们当然无权因此向我发怒。至于那些虽然知道并理解我的指示，却未能参透圣经中难解段落之意义的人，他们好比在上述情况下，只能看见我的手指，却看不见我所指的方向。因此，这两类人最好停止责备我，转而祈求天主赐予他们眼睛的视力。因为虽然我能移动手指指向某物，但我无法开启人的眼睛，使他们看到我是在指，或者我所指的对象。\par

4. 但是，至于那些自夸天主恩宠的人，吹嘘没有我现在拟订的这种指引也能理解并解释圣经，因而认为我所要写的完全是多余的人。我希望这样的人能冷静下来，记住：不论他们多么有理由因天主的大恩而喜乐，他们自己最初也是从人间的老师那里学会阅读的。他们几乎不会认为，因此他们就应当被\Place{埃及}隐修士\Person{安当}——一位正义而圣洁的人，他自己不识字，据说通过听别人诵读而将圣经记在心中，并靠着智慧的默想通达了圣经——所鄙视，或被我最近从非常可敬而可信的见证人那里听说的那位外邦奴隶\Person{基利斯当}所鄙视，他未经任何人的教导，仅靠祈祷求天主启示他，就完全掌握了阅读技艺；经过三天的祈求，他如愿以偿，当场就能读完整本被惊呆的旁观者递给他的书。\par

5. 但如果有人认为这些故事是假的，我并不坚持。因为，我是在与那些宣称不用从人那里得到任何指引就能理解圣经的基督徒打交道（如果事实如此，他们确实夸耀一个真正的、非同寻常的优点），他们想必会承认，我们每个人都是从孩提时代不断听人说话而学会了自己的母语，而我们学会的任何其他语言——希腊语、希伯来语或其他语言——要么是通过听人说话，要么是从人间的老师那里学来的。那么，假设我们建议所有弟兄不要教子女这些事，因为圣神降临之际，宗徒们立即开始说各种语言；并且警告每一个没有类似经验的人，他不必认为自己是个基督徒，或至少可以怀疑自己是否已领受了圣神？不，不；宁可让我们放下虚假的骄傲，学习一切能从人那里学到的东西；让教导他人的人，以无傲慢、无嫉妒的心，传递他所领受的。并且不要试探我们所信的主，免得陷入仇敌的这等诡计和我们自己的执拗之中，以致我们甚至拒绝去教堂聆听福音本身，或拒绝读一本书，或拒绝听别人诵读或讲道，而指望自己会被提到第三层天上，如宗徒所说，\Quote{或在身内，或在身外，}\BibleRef{格后 12:2}，在那里听到不可言说的话语，是人不可说的，或看见主耶稣基督，并从祂亲口而非从人口中听到福音。\par

6. 让我们警惕这种骄傲的危险诱惑，反而要思想：宗徒保禄本人虽被来自天上的天主声音击倒并受训诫，却仍被差遣到一个人那里去领受圣事并加入教会（\BibleRef{宗 9:3}）；而百夫长科尔乃略，虽有天使向他宣告他的祈祷已蒙垂听、他的施舍已蒙纪念，却仍被交给伯多禄去受教导，不仅从宗徒手中领受圣事，而且还受他关于信德、望德和爱德的正当对象的教导。（\EditorialNote{宗徒大事录第十章}）无疑，借着天使的媒介完成一切是可能的，但假若天主不愿意使用人作为祂圣言向同胞传递的工具，我们人类的情况就会更加卑下。因为如果天主不从祂的人性圣殿发出神谕，而是事事通过天上的声音或天使的服务来教导人，那么所写的\BibleQuote{天主的宫殿是圣的，这宫殿就是你们}（\BibleRef{格前 3:17}）这句话怎能成为事实？再者，若世人从不从同胞那里学到什么，那将人结合在团结纽带中的爱德本身，也无法将灵魂注入灵魂并使之彼此相融。\par

7. 我们知道，那位宣读依撒意亚先知书却不明白所读之物的太监，并没有被宗徒差遣到天使那里去，也没有天使向他解释他不明白的，更没有在天主恩宠中不经人中介而受内在光照；相反，因天主的提示，懂得那先知书的斐理伯来到他那里，与他同坐，用人的言语和人的口舌为他开启了圣经。（\BibleRef{宗 8:26}）天主不是曾与梅瑟交谈吗？然而他却以极大的智慧、毫无嫉妒的骄傲，接纳了他那异族岳父的计划，来治理和管理托付给他的伟大民族的事务。（\BibleRef{出 18:13}）因为梅瑟知道，一个明智的计划，无论它出自何人的心智，都应归功于那位创造它的主，即那不变的真理、天主。\par

8. 最后，凡自夸借着天主光照理解圣经疑难、而未受任何诠释规则指教的人，同时相信（且正确地相信）这种能力不是出于自己（即不是源于自身），而是天主的恩赐。因为他寻求的是天主的荣耀，而非自己的荣耀。然而，既然他无需任何人间诠释者的帮助就能阅读和理解，为什么他自己却要为别人诠释？为什么他不直接把他们送到天主那里，使他们也可以在没有人的帮助、借着圣神的内在教导而学习？事实是：他害怕招致这样的指责：\BibleQuote{你这可恶又懒惰的仆人，你该把我的钱交给兑换银钱的人。}（\BibleRef{玛 25:26}）那么，既然这些人或通过言论或通过写作把自己所理解的教给别人，他们当然不能责备我，如果我也教导他们所理解的以及他们所遵循的诠释规则。因为除了虚假，无人应把自己的东西视为己有。一切真理都属于那说\BibleQuote{我是真理}者（\BibleRef{若 14:6}）。因为\BibleQuote{你们有什么不是领受的呢？既然是领受的，为什么还夸耀，好像不是领受的呢？}（\BibleRef{格前 4:7}）\par

9. 对听众宣读的人，大声读出他眼前所见的话语；教授诵读的人，则是为了使别人能够自己阅读。然而，每个人都将自己所学的传授给别人。同样，向听众解释他所理解的圣经段落的人，就像大声宣读面前文字的人。另一方面，制定诠释规则的人，就像教授诵读的人，即向他人展示如何为自己阅读。因此，正如知道如何阅读的人，在找到一本书时，不依赖别人告诉他里面写的是什么；同样，拥有我在此尝试制定的规则的人，若在阅读的书中遇到难懂段落，将不需要一位诠释者向他揭示奥秘，而是坚守某些规则，遵循某些迹象，就能达到隐藏的意义而无错误，或至少不陷入任何显著的荒谬。所以，尽管在本书的过程中会充分表明，没有人能正当地反对我这项除了服务别无他图的工作，但既然在开头答复可能提出初步异议的人似乎是合宜的，这就是我对即将在本书中走过的道路所选择的开端。\par

\SourceRecord{Translated by James Shaw.；Nicene and Post-Nicene Fathers, First Series；Vol. 2.；Edited by Philip Schaff.；Buffalo, NY: Christian Literature Publishing Co.,；1887.；<http://www.newadvent.org/fathers/12020.htm>.}

% structure: p0001 source=h1 target=section reason=page-title

\section{论基督圣道（卷一）}

包含圣经所论主题之总览\par

% structure: p0003 source=h2 target=subsection reason=relative-heading-rank; base=h2

\subsection{第一章——圣经的解释取决于意义的发现与表达，且须依赖天主的助佑而进行。}

1. 一切圣经的解释都依赖于两件事：确定正当意义的方式，以及当意义确定后将之表达出来的方式。我们将先讨论确定意义的方式，其次讨论表达意义的方式——这是一项伟大而艰巨的工作，如果难以实行，我担心冒昧着手。如果我是依靠自己的力量，那无疑是冒昧的；但由于完成这工作的希望寄托于那已经在这主题上赐给我许多思想的天主，我不害怕祂会在我开始使用祂已赐予的东西之后，继续供应所欠缺的。因为一份与他人分享却不减少的财富，如果拥有而不分享，就还不是应当拥有那样拥有。主说：\Quote{凡有的，还要给他。}\BibleRef{玛 13:12} 因此，祂将赐给那有的；也就是说，如果他们自由而乐意地使用所领受的，祂将增添并圆满祂的恩赐。奇迹中的饼，在门徒开始分给饥饿的人群之前，只有五个和七个。但当他们开始分发时，虽然满足了成千上万人的需要，却装满了剩下的碎块篮子。正如那饼在掰开的过程中增加一样，主已惠赐给我用于着手这项工作的思想，一旦我开始将它们传授给他人，也必因祂的恩宠而倍增，以致在这分施的工作中，我非但不致遭受损失和匮乏，反而要因财富的奇妙增加而欢欣。\par

% structure: p0005 source=h2 target=subsection reason=relative-heading-rank; base=h2

\subsection{第二章——什么是物，什么是记号。}

2. 一切教导都或是关于事物的，或是关于记号的；但事物是通过记号来认识的。我现在严格使用\Quote{物}这个词，指那些从不被用作其他事物记号的东西：例如，木头、石头、牲畜以及其他这类东西。然而，并非我们读到梅瑟投到苦水中使之变甜的木头（\BibleRef{出 15:25}），也不是雅各伯用作枕头的石头（\BibleRef{创 28:11}），也不是亚巴郎代替自己儿子献祭的公羊（\BibleRef{创 22:13}）；因为这些虽然是物，但也同时是其他事物的记号。还有另一种记号，它们从来只被用作记号：例如，词语。没有人使用词语而不作为其他事物的记号；因此可以理解我所谓的记号：即那些用来指示其他事物的东西。所以，每个记号同时也是一事物；因为不是事物的东西根本不存在。然而，并非每一事物同时也是记号。因此，关于事物和记号的这种区分，当我谈论事物时，我会以这样的方式谈论，即使其中一些也可能被用作记号，这也不会干扰我按照先讨论事物后讨论记号的划分。但我们必须小心记住，我们现在要考虑关于事物的是它们本身是什么，而不是它们是什么事物的记号。\par

% structure: p0007 source=h2 target=subsection reason=relative-heading-rank; base=h2

\subsection{第三章——有些事物是供使用的，有些是供享受的。}

3. 因此，有些事物是应享受的，有些是应使用的，还有一些既是享受者也是使用者。那些作为享受对象的事物使我们幸福。那些作为使用对象的事物协助并（可以说）支持我们追求幸福，以便我们能达到那些使我们幸福并安息于其中的事物。而我们自己，又享受和使用这些事物，被置于两类对象之间，如果我们执意享受那些本应使用的事物，就会在旅程中受阻，有时甚至被引入歧途；以致陷入对低级满足的爱恋之中，懈怠于、甚至完全背离了追求真正而恰当的享受对象。\par

% structure: p0009 source=h2 target=subsection reason=relative-heading-rank; base=h2

\subsection{第四章——使用与享受的区别。}

4. 因为享受一件事物，就是为它本身而满意地安息于其中。而使用，则是运用可用的任何手段来获得所渴望的东西，如果这渴望是正当的对象；因为不正当的使用更应称为滥用。那么，假如我们是在异乡的流浪者，无法离开祖国而幸福生活，感到流浪的痛苦，希望结束苦难，决心返回家乡。然而我们发现，必须借助某种运输方式，或陆路或水路，才能到达那个将开始享受的家乡。但我们沿途经过的风景之美，以及运动本身的乐趣，迷住了我们的心，把这些本应使用的事物当作享受的对象，就不愿意加快旅程的终点；并且沉溺于虚假的愉悦中，我们的思想就背离了那个其欢乐能使我们真正幸福的家乡。这就是我们在今世可朽生命中的状况。我们已经远离天主；如果我们愿意返回父家，就必须使用这个世界，而不是享受它，以便藉着所造之物，可以清楚地看见那看不见的天主的事（\BibleRef{罗 1:20}）——也就是，藉着物质和暂时的事物，去把握那属灵和永恒的事物。\par

% structure: p0011 source=h2 target=subsection reason=relative-heading-rank; base=h2

\subsection{第五章——天主圣三是享受的真正对象。}

5. 因此，真正的享受对象是圣父、圣子和圣神，他们同时也是天主圣三，同一本体，超越万有，并共同与所有享受祂的人同在——如果祂是一个对象，而不是一切对象的根源，甚或即便祂是万有的根源。因为不易找到一个名称能恰当地表达如此伟大的卓越，除非这样说：天主圣三，唯一的天主，万有都出于祂，藉着祂，并在祂内。\BibleRef{罗 11:36} 因此，圣父、圣子和圣神，以及每一位单独而言，都是天主，而同时他们全体是唯一的天主；每一位单独而言都是一个完整的实体，然而他们全体是同一实体。圣父不是圣子，也不是圣神；圣子不是圣父，也不是圣神；圣神不是圣父，也不是圣子：但圣父只是圣父，圣子只是圣子，圣神只是圣神。三者拥有同一永恒、同一不变性、同一威严、同一权能。在圣父内是统一，在圣子内是相等，在圣神内是统一与相等的和谐；而这三者因圣父而归于统一，因圣子而归于相等，因圣神而归于和谐。\par

% structure: p0013 source=h2 target=subsection reason=relative-heading-rank; base=h2

\subsection{第六章  天主在何种意义上不可言喻。}

6. 我曾以任何相称的方式论及天主，或颂扬祂吗？不，我感到我只不过渴望言说；而若我说了什么，那并非我所渴望说的。我如何知道这一点，除非因为天主是不可言喻的？但我所说的，若是不可言喻，就不可能被说出。因此，天主甚至不可被称为\Quote{不可言喻}，因为这样说也是在谈论祂。于是就产生了一种奇怪的词语矛盾，因为若不可言喻是无法被言说的，那么它若能被叫做不可言喻，它就不是不可言喻的。这种词语的对立更宜以沉默规避，而非以言语解释。然而，虽然对于天主的伟大，没有什么相称之言可说，祂却屈尊接受人嘴唇的敬拜，并愿意我们藉着我们自己的言语在赞美祂时喜乐。正是基于这一原则，祂被称为\OriginalTerm{天主}{Deus}。因为这两个音节的声音本身并不传达有关祂本性的真知识；但凡是通晓拉丁文的人，当那声音传到他们耳中时，都会被引向思想一种在卓越上至高、在存在上永恒的本性。\par

% structure: p0015 source=h2 target=subsection reason=relative-heading-rank; base=h2

\subsection{第七章  所有人对\Quote{天主}一词的理解。}

7. 因为当人们想到那至高的唯一的天主（诸神之神）——甚至被那些相信有其他神明并称它们为神、敬拜它们为神的人——时，他们的思想努力达到一种概念，即没有什么比这更卓越或更崇高。由于人受不同种类的快乐所驱使，一部分与身体感官相关，一部分与理智和灵魂相关，那些受感官束缚的人认为，要么是天空，要么是天空中最耀眼的东西，要么是宇宙本身，是神之神；或者如果他们试图超越宇宙，他们便向自己描绘某种耀眼明亮的东西，模糊地将其视为无限的，或具最可想象的最美形态；或者他们以人体的形式来表现它，如果他们认为那胜过其他一切。或者如果他们以为没有一个至高的天主高于其余，而是有许多甚至无数的等级相同的神明，他们仍然认为这些神明具有形状和形式，各人按照自己认为的卓越模式。另一方面，那些努力以理智的努力达到对天主的理解的人，将祂置于一切可见的和有形体的本性之上，甚至置于一切可变化的理智和精神本性之上。然而，所有人都争先恐后地推崇天主的卓越：也找不到任何人相信任何存在者如果有比它更优越的就是天主。因此，所有人都一致相信，天主就是那在尊贵上超越一切其他对象的存在。\par

% structure: p0017 source=h2 target=subsection reason=relative-heading-rank; base=h2

\subsection{第八章  天主应被尊崇超越一切，因为祂是不变的智慧。}

8. 既然所有思考天主的人都认为祂是活的，那么只有那些将祂视为生命本身的人才能形成对祂的任何不荒谬且相称的观念；并且，无论他们心中联想到什么形体形式，他们都认识到，是生命使之活或不活，并且偏爱活的东西胜过死的东西；他们明白，活的形体本身，无论它在光彩上如何超越其他，在尺寸上如何凌驾，在美丽上如何卓越，都与使生命活跃的生命完全不同；并且他们认为生命在尊贵和价值上无比优越于它所活跃和赋予生气的物质。然后，当他们继续考察生命本身的本性时，如果他们发现它仅仅是营养性的生命，没有感性，如植物的生命，他们认为它低于有感性的生命，如牲畜的生命；而在此之上，他们又将理智性的生命，如人的生命，置于更高地位。而察觉到即使是这也处于变化中，他们被迫再次将不变的生命置于其上，这不变的生命不是一时愚拙、一时智慧，而是相反，它是智慧本身。因为一个明智的理智，即已获得智慧的理智，在获得智慧之前是不明智的。但智慧本身从未是不明智的，也永远不会成为不明智的。如果人们从未瞥见这智慧，他们绝不可能完全确信地偏爱不变智慧的生命胜过可变化的生命。这一点将是明显的，如果我们考虑他们用以断言不变生命更为卓越的真理规则本身是不变的：而他们无法找到这样的规则，除非超越他们自己的本性；因为他们在自身内找不到任何不是变化的东西。\par

% structure: p0019 source=h2 target=subsection reason=relative-heading-rank; base=h2

\subsection{第九章  所有人都承认不变智慧优于可变智慧。}

9. 现在，没有一个人会如此极端愚蠢，以至于问：\Quote{你怎么知道不变智慧的生命比可变智慧的生命更可取？} 因为那关于他所问的\Quote{我怎么知道}的真理本身，是不变地刻在所有人的心中，并呈现在他们共同的默观中。而那看不见这真理的人，就像阳光下的盲人，那光芒的灿烂如此清晰、如此临近，倾泻入他的眼球中，却对他毫无益处。反之，那看见却退缩于这真理的人，是因长久居于肉身的阴影中而心智视力薄弱。因此，人们被邪恶习惯的逆风从故乡驱回，追求那些较低等、较无价值的对象，而非他们自己承认更优越、更值得的对象。\par

% structure: p0021 source=h2 target=subsection reason=relative-heading-rank; base=h2

\subsection{第十章 为见天主，灵魂必须净化。}

10. 因此，既然我们有责任完全享受那不变而生活的真理，而且圣三的天主按这真理为祂所造之物筹谋，那么灵魂必须被净化，才能有能力感知那光，并在感知后安息于其中。让我们把这净化看作是一种旅程或航行，返回我们的故乡。因为不是通过改变地点，我们才能接近那位无所不在者，而是通过培养纯洁的愿望和德行习惯。\par

% structure: p0023 source=h2 target=subsection reason=relative-heading-rank; base=h2

\subsection{第十一章 智慧降生成人，成为我们净化的楷模。}

11. 但对此我们本来完全无能为力，除非智慧屈尊俯就，适应我们的软弱，并以我们人性的形态向我们展示圣洁生活的楷模。然而，当我们接近祂时是明智的，但当祂来到我们这里时，却被骄傲的人视为极其愚拙。当我们接近祂时变得强壮，但当祂来到我们这里时，却被视为软弱。但是，\BibleQuote{天主的愚拙比人更智；天主的软弱比人更强。} \BibleRef{格前 1:25} 因此，虽然智慧本身就是我们的家园，祂也使祂自己成为我们到达家园的道路。\par

% structure: p0025 source=h2 target=subsection reason=relative-heading-rank; base=h2

\subsection{第十二章 天主的智慧以何种意义来到我们当中。}

虽然祂处处临在于内眼之中，当内眼健全清晰时，祂却屈尊俯就，使祂自己向那些内眼微弱模糊的人的外眼显现。\BibleQuote{因为，在天主的智慧里，世人凭智慧不认识天主，天主就乐意用人所当作愚拙的宣讲来拯救那些相信的人。} \BibleRef{格前 1:21}\par

12. 那么，不是在穿越空间的意味上，而是因为祂以可朽血肉的形态向凡人显现，因此祂被说成是来到我们这里。因为祂来到一个祂一直所在的地方，虽然如此，\BibleQuote{祂在世界，世界原是藉祂造成的。} 但是，因为人们——他们迫切地享受受造物而非造物主，因而变得像这世界，所以被最恰当地称为\Quote{世界}——没有认出祂，因此福音作者说，\BibleQuote{世界却不认识祂。} \BibleRef{若 1:10} 因此，在天主的智慧里，世人凭智慧不认识天主。那么，既然祂已经在这里，为什么祂还要来呢？除非是祂乐意用愚拙的宣讲来拯救那些相信的人？\par

% structure: p0028 source=h2 target=subsection reason=relative-heading-rank; base=h2

\subsection{第十三章 圣言成了血肉。}

祂来除了这样，\BibleQuote{圣言成了血肉，居住在我们中间}？\BibleRef{若 1:14} 正如当我们说话时，为使我们心中所想能通过耳朵进入听者的心，我们心里的言词成为外在的声音，称为言语；然而我们的思想并不消失在声音里，而是保持完整，并呈现为言语的形式，但其本性并不因这一变化而有所改变：同样，圣言虽然本性毫无改变，却成了血肉，为能居住在我们中间。\par

% structure: p0030 source=h2 target=subsection reason=relative-heading-rank; base=h2

\subsection{第十四章 天主的智慧如何治愈了人。}

13. 此外，正如使用药物是通往健康的道路，同样，这药物接纳了罪人，以治愈和恢复他们。又正如外科医生包扎伤口时，并非草率行事，而是仔细地，使包扎除了实用之外，也有一定程度的整洁；同样，我们的药物——智慧——藉祂摄取人性而适应我们的伤口，以相反者治愈某些伤口，以类似者治愈另一些伤口。正如处理身体创伤的人，有时施用相反者，如冷对热、湿对干等等；在其他情况下则施用类似者，如圆布对圆伤口，或长布对长伤口，并且不给所有肢体配上同样的绷带，而是以类似对类似；同样，天主的智慧在治愈人时，将祂自己应用于人的治疗，祂自己同时是治疗者和药物。因此，既然人因骄傲而堕落，祂便藉谦卑恢复了人。我们被蛇的智慧所诱捕：我们被天主的愚拙所释放。且正如前者被称为智慧，实为那些蔑视天主之人的愚拙，后者被称为愚拙，但在那些战胜魔鬼的人身上，却是真智慧。我们如此滥用我们的不朽，以致招致死亡的惩罚；基督如此善用祂的可死性，以致恢复我们的生命。疾病是通过一个女人的腐败灵魂引入的；解药是通过一个女人的童贞身体而来。同样属于相反疗法的范畴，我们的恶习因祂的德行榜样而治愈。另一方面，以下则像是与所缠裹的肢体的伤口形状相同的绷带：祂由女人所生，为拯救我们这些因女人而堕落的人；祂作为人而来，为拯救我们这些身为人的；作为可死者，为拯救我们这些可死的；藉着死亡，为拯救我们这些已死的。而那些能够更完整地探究此事的，未被执行某项既定任务的必要性所催促的人，在思考基督宗教的医学中使用的药物（无论是相反者还是类似者）时，将会发现许多其他教导之处。\par

% structure: p0032 source=h2 target=subsection reason=relative-heading-rank; base=h2

\subsection{第十五章 信德因基督的复活与升天而加强，并因祂的审判来临而受激励。}

14. 我主从死者中复活并升天的信仰，藉着增添望德的巨大支柱，坚固了我们的信德。因为这事清楚地显示，当祂有能力如此再次取回生命时，祂是多么自由地为我们舍去了生命。那么，当信徒们思考，当他们仍在不信中时，那为他们忍受了如此伟大事物的是何等伟大的主时，他们的望德是何等地受到鼓舞啊！当人们期盼祂从天降临，作为审判生者死者的审判者时，这在粗心大意的人心中引起极大的恐惧，促使他们投身于勤勉的准备，并学习藉着圣洁的生活来渴望祂的来临，而不是因自己的恶行而战栗。当我们想到，为在这尘世的旅途中安慰我们，祂如此慷慨地将祂的圣神赐予我们，使我们在今生的逆境中仍能保持对那尚未看见之主的信赖与爱慕，并且祂还赐予每人适当的神恩，用以建立祂的教会，使我们能遵行祂所指示当做的事，不仅毫无怨言，甚至满心喜悦时，有什么口舌能述说，或什么想象能构思祂在末后要赐予的赏报呢？\par

% structure: p0034 source=h2 target=subsection reason=relative-heading-rank; base=h2

\subsection{第十六章——基督以医治性的苦难炼净祂的教会。}

15. 因为教会是祂的身体，正如宗徒的教导向我们显示的；它甚至被称为祂的新娘。所以，祂将那个有许多肢体、且各具不同功能的身体，在合一与爱德的联系中结合一起，这就是它真正的健康。此外，祂在现今时期锻炼它，并用许多有益的苦难炼净它，以便当祂将它从这个世界移植到永恒世界时，祂能把它当作新娘，毫无玷污或皱纹，或任何这类东西，带到自己面前。\par

% structure: p0036 source=h2 target=subsection reason=relative-heading-rank; base=h2

\subsection{第十七章——基督藉着赦免我们的罪，打开了通往我们家园的道路。}

16. 此外，当我们行在路上，且那路并非穿越空间，而是藉着情感的转变，并且我们过去之罪的过犯如同荆棘篱笆阻挡我们时，那愿意将自己作为我们应回归之路而躺下的祂，除了赦免我们所有的罪，并为我们被钉十字架，以除去那阻挡我们回归之门的严厉法令外，还能做什么更仁爱、更慈悲的事呢？\par

% structure: p0038 source=h2 target=subsection reason=relative-heading-rank; base=h2

\subsection{第十八章——钥匙交给了教会。}

17. 因此，祂将钥匙交给了祂的教会，以便凡在地上捆绑的，在天上也捆绑；凡在地上释放的，在天上也释放；这就是说，凡在教会中不相信自己的罪被赦免的人，他的罪就不得赦免；但凡相信并悔改，且离开其罪的人，将因同样的信德和悔改而得救，并基于此被收进教会的怀抱。因为那不相信自己的罪能被赦免的人，便陷入绝望，并变得更为恶劣，仿佛除了作恶之外，再没有更大的好处留给他，因为他不再对自己的悔改结果持守信德。\par

% structure: p0040 source=h2 target=subsection reason=relative-heading-rank; base=h2

\subsection{第十九章——肉身与灵魂的死亡与复活。}

18. 此外，正如有一种灵魂的死亡，它在于摒弃旧习和旧的生活方式，并且是通过悔改而来的，同样地，肉身的死亡也在于先前生命原则的解散。并且，正如灵魂在藉着悔改摒弃并摧毁旧习之后，按照一个更好的模式被重新创造，同样地，我们必须希望并相信，肉身在经历那因罪所欠的死亡之后，将在复活时改变成更美好的形态；——并非血肉之躯能继承天主的国（因为那是不可能的），而是这必朽坏的要穿上不朽坏，这必死的要穿上不死。\BibleRef{格前 15:50} 因此，身体，因不再感到匮乏而成为无忧之源，将由一个完全纯洁幸福的灵魂所驱策，并享有不间断的平安。\par

% structure: p0042 source=h2 target=subsection reason=relative-heading-rank; base=h2

\subsection{第二十章——复活的永罚。}

19. 如今，那灵魂不死于这个世界、且在此世未开始肖似真理的人，肉身死亡时，便落入更可怕的死亡，并且他将复活，不是要将他地上的居所换成天上的居所，而是要承受他罪恶的刑罚。\par

% structure: p0044 source=h2 target=subsection reason=relative-heading-rank; base=h2

\subsection{第二十一章——死亡时肉身与灵魂均不消灭。}

因此，信德紧握着这确信，我们必须相信事实确是如此：人的灵魂和肉身都不会遭受完全的消灭，而是恶人复活承受不可思议的惩罚，善人复活获得永生。\par

% structure: p0046 source=h2 target=subsection reason=relative-heading-rank; base=h2

\subsection{第二十二章——唯独天主应被享受。}

20. 那么，在这一切事物中，只有那些我们所说永恒不变的事物才是真正的享受对象。其余的都是为使用，使我们能够达到对前者的圆满享受。然而，我们这些享受和使用其他事物的人，我们自己也是事物。因为人实在是一件伟大的受造物，是按天主的肖像和相似而造的，不是就其披戴的必死肉身而言，而是就其理性灵魂而言，这灵魂使人尊荣超乎禽兽之上。因此，一个重要的问题是：人是否应当享受自己，或使用自己，或两者都做？因为我们被命令要彼此相爱：但问题是，一个人是因其自身而被另一个人爱，还是为了别的原因？若是为了其自身，我们便享受他；若是为了别的，我们便使用他。那么，在我看来，他是为了别的原因而被爱的。因为若一个事物是为了其自身而被爱，那么幸福的生活就在于享受它，至少在现世，这幸福的希望——即使不是现实——是我们的安慰。但寄望于人的人，是被诅咒的。\BibleRef{耶 17:5}\par

21. 细加审视，任何人也不当以自己为乐，因为没有人应当为自身的原因而爱自己，甚至爱自己，也应为祂而爱，祂才是真正的享乐对象。因为一个人的最佳状态，莫过于他整个生命都在向那不变的生命迈进，并且他的情感完全专注于那生命。然而，如果他为了自身的缘故爱自己，他就不是从与天主的关系来看自己，而是将心神转向自我，因此他并没有专注于任何不变的事物。这样，他就没有以最好的方式享受自己，因为当他的心完全专注于不变之善，情感完全投入其中时，他比自己从那里转向去享受自身时更好。因此，如果你甚至不应为自身而爱自己，而应为祂而爱，你的爱在祂那里找到最值得的对象，那么若你也为天主的缘故而爱他人，没有人有权生气。因为这是神圣权威所制定的爱的法律：\Quote{你当爱邻人如同自己}；但是，\Quote{你当全心、全灵、全意爱天主}；这样，你就应当将你的一切思想、整个生命和全部理智都集中于祂，你从祂那里领受你所带来的一切。因为当祂说\Quote{全心、全灵、全意}时，意思是我们的生命中没有任何部分闲置，不为享受其他对象的欲望留下空间，而是任何其他似乎值得爱的事物，都应当被导入我们全部情感所流的同一渠道。因此，凡是正确爱邻人的人，也应当敦促邻人全心、全灵、全意爱天主。因为这样，人爱邻人如同自己，就将对自己和邻人的全部爱之流都导入爱天主的渠道，这渠道不容许任何支流从中分走，以致其自身的流量减少。\par

% structure: p0049 source=h2 target=subsection reason=relative-heading-rank; base=h2

\subsection{第二十三章 人无需受命爱自己和自己的身体}

22. 然而，并非所有用物都当爱，只有那些要么与我们在共同与天主的关系中结合在一起的，如一个人或一位天使，要么如此与我们相关，以致需要经由我们而获得天主的善的，如身体。因为确然，殉道者并不爱迫害者的恶，虽然他们利用这恶以获得天主的恩宠。因此，当爱之物有四类：第一，在我们之上的；第二，我们自己；第三，与我们平等的；第四，在我们之下的。关于第二类和第四类，无需颁布任何诫命。因为无论人离真理多远，他仍然继续爱自己，并爱自己的身体。那远离不变之光、万有主宰的灵魂，之所以如此行，是为了统治自己和自己的身体；因此它不能不既爱自己又爱自己的身体。\par

23. 此外，它以为若能使同伴，即其他人，服从自己，便获得了非常伟大的成就。因为对有罪的灵魂而言，最渴望并自认为当得的是唯独天主才当得的。然而，这种对自己的爱更正确地被称为恨。因为欲望之下的东西服从自己，而自己却不服从自己的在上者，这是不义的；因此至公义地说：\Quote{喜爱不义的人，恨自己的灵魂。}于是灵魂变得软弱，并因可朽的身体而承受许多痛苦。当然，它必须爱身体，并为身体的腐朽而悲伤；而身体的不朽和不腐性出自灵魂的健康。灵魂的健康在于坚定地依附于更好的那一位，即不变的天主。但当它想甚至统治那些本性上与它平等的，即同伴时，这是完全无法容忍的傲慢。\par

% structure: p0052 source=h2 target=subsection reason=relative-heading-rank; base=h2

\subsection{第二十四章 没有人憎恨自己的肉身，连那些滥用它的人也不例外}

24. 没有人恨自己。关于这点，确实任何学派都从未提出过疑问。但也没有人恨自己的身体。因为宗徒确实说：\Quote{没有人曾经恨过自己的肉身。}\BibleRef{弗 5:29}当有些人说他们宁愿完全没有身体时，他们完全欺骗了自己。因为他们恨的不是他们的身体，而是身体的腐朽和沉重。因此，他们想要的是没有身体，而是不朽而且非常轻盈的身体。但他们以为那样的身体根本就不是身体，因为他们认为这种东西必然是一种神灵。至于他们似乎以某种方式通过禁食和劳苦来鞭笞身体，那些以正意这样做的人，并非为了摆脱他们的身体，而是为了使身体服从并随时准备做一切必要的工作。因为他们藉着身体的某种劳苦训练，根除那些有害于身体的欲望，即导致享受不值得之物的灵魂的习惯和情感。他们不是在毁灭自己，而是在照顾自己的健康。\par

25. 反之，那些以乖戾精神行事的人，向自己的身体作战，仿佛它是天生的仇敌。在这事上，他们因误解所读的经文而走入歧途：\BibleQuote{肉性相反圣神，圣神相反肉性，二者彼此相敌。}（\BibleRef{迦 5:17}）这是指尚未被制服的肉性习染，圣神渴望反对它，并非要毁灭身体，而是要铲除身体的贪欲——\Emph{即}其邪恶习惯——从而使身体服从圣神，这正是自然秩序的要求。因为正如复活后，身体完全服从圣神，将在永恒的完全平安中生活；同样，在今世我们也必须以改善肉性习染为目标，使其不当的情欲不再攻击灵魂。直到这事实现，\BibleQuote{肉性相反圣神，圣神相反肉性}；圣神奋斗，不是出于仇恨，而是为了主宰，因为它渴望它所爱的应服从更高原则；肉性奋斗，不是出于仇恨，而是因为它从祖传而来的习惯束缚，这束缚藉自然法则滋长并变成根深蒂固。因此，圣神在制服肉性时，犹如在摧毁一种邪恶习惯的虚假平安，并带来一种源于良好习惯的真正平安。然而，即使那些因错误观念而仇恨自己身体的人，也不愿牺牲一只眼睛，即使他们可以毫无痛苦地这样做，并且一只眼睛的视力与从前两只眼睛一样，除非为了达到某种超过损失的目标。这一点和其他类似迹象足以向真诚寻求真理者表明，宗徒所说的：\BibleQuote{从来没有人恨恶自己的肉身}是多么有根据。他又加上：\BibleQuote{反而养育它，珍爱它，如同主对待教会一样。}（\BibleRef{弗 5:29}）\par

% structure: p0055 source=h2 target=subsection reason=relative-heading-rank; base=h2

\subsection{第二十五章——人可能爱某物胜过自己的肉身，但并不因此恨自己的肉身。}

26. 因此，人应当受到教导关于爱的适当尺度，即他应以何种尺度爱自己，从而对自己有益。因为没有人（除非愚人）会怀疑他爱自己并愿意对自己行善。同样，他也应受教导以何种尺度爱自己的身体，以便明智地在适当限度内照料它。因为同样明显，他也爱自己的身体，并愿意保持它安全健康。然而，人可能有比身体的安全和健康更爱的东西。因为许多人曾自愿忍受痛苦和切除某些肢体，以获取他们更珍视的其他对象。但没有人被告知不应渴望身体的安全和健康，因为他有更渴望的东西。因为吝啬者虽然爱钱，却为自己买面包——也就是说，他付出他很喜欢并渴望积攒的钱——但这是因为更珍视面包所维持的身体健康。再详细讨论如此明显之点是多余的，但这正是恶人的错误常常迫使我们做的。\par

% structure: p0057 source=h2 target=subsection reason=relative-heading-rank; base=h2

\subsection{第二十六章——爱天主和爱近人的诫命包含了爱自己的诫命。}

27. 既然不需要命令人爱自己和自己的身体——也就是说，我们通过从未被违反的自然法则爱自己以及我们之下但与我们相连的事物（因为连野兽也爱自己和自己的身体），那么只需对我们关于我们之上的天主和我们旁边的近人颁布训令即可。他说：\BibleQuote{你应全心、全灵、全意爱上主你的天主，并应爱你的近人如你自己。全部法律和先知都系于这两条诫命。}（\BibleRef{玛 22:37}）这样，诫命的终向就是爱，而且是双重的爱，即爱天主和爱近人。现在，如果你以你的整体——即灵魂和身体一起——以及你的近人的整体——灵魂和身体一起（因为人由灵魂和身体构成），你会发现这两条诫命并未忽略任何应爱之物的类别。因为尽管爱天主在先，并且我们爱他的尺度这样被规定，以至于显然其他一切都要以他为中心，但关于爱自己似乎没有说什么；然而当说\Quote{你应爱你的近人如你自己}时，立刻就显明我们对自身的爱没有被忽略。\par

% structure: p0059 source=h2 target=subsection reason=relative-heading-rank; base=h2

\subsection{第二十七章——爱的秩序。}

28. 一个公义圣洁的人，他对事物作出无偏见的判断，并严格约束自己的情感，以致他不爱他本不该爱的，也不缺乏他本应爱的，爱不当爱之深者不过分，爱不当爱之浅者不浅，爱当平等者不偏不倚。没有一个罪人应被当作罪人来爱；每一个人都应为天主之故被当作人来爱；但天主应为他自身之故被爱。如果天主应被爱超过任何人，那么每个人爱天主应超过爱自己。同样，我们爱另一个人应胜过爱自己的肉身，因为所有事物都应参照天主来爱，而另一个人能与我们共享天主之福乐，我们的肉身却不能；因为肉身只靠灵魂生活，而我们是藉着灵魂享受天主。\par

% structure: p0061 source=h2 target=subsection reason=relative-heading-rank; base=h2

\subsection{第二十八章——如何决定该帮助谁。}

29. 再者，所有的人应受同等的爱。但既然你无法对所有人都行善，就应特别留意那些因时间、地点或环境的机缘而与你更为亲近的人。例如，假设你拥有大量某种物品，并感到有责任将它送给某个一无所有的人，而它又不能给予多人；若两人同时出现，其中任何一个既无需要也无亲属关系比另一个有更大的权利要求你，那么你所能做的最公平之事，莫过于用抽签来决定将不能给两人的物品给其中一人。同样地，在众人中间：既然你不能为所有人的益处着想，就必须将此事视为已由某种抽签为你决定，即根据每个人在当时与你更亲近的程度来对待。\par

% structure: p0063 source=h2 target=subsection reason=relative-heading-rank; base=h2

\subsection{第29章——我们应渴望并努力使所有人爱天主。}

30. 在一切能与我们一起享有天主的人中，我们爱一部分人是出于向他们提供服务，一部分人是出于他们向我们提供服务，一部分人是出于他们既帮助我们于需要之中，又反过来受我们帮助，一部分人是出于我们既未给予他们利益，也不期望从他们获得任何利益。然而，我们应当渴望他们全都与我们一同爱天主，而我们给予他们或从他们接受的一切帮助，都应趋向于那同一个目的。因为在剧场中（尽管它们是邪恶的巢穴），若一个人喜爱某位演员，并将他的技艺视为伟大甚至至善，他就会喜爱所有与他一同赞赏其偶像的人——不是为了他们本身，而是为了他们所共同崇拜的那一位；他越是热切地赞赏，就越会尽力设法为他的偶像获得新的崇拜者，也越急于将偶像展示给别人；若他发现某人相对冷漠，他就会极力赞扬偶像的功绩以激发其兴趣；然而，若他遇见反对者，他便因偶像受轻视而极为不悦，并竭尽全力消除这种轻视。如果连在剧场中也是如此，那么，我们这些生活在天主的爱之共融中的人——享见天主是人生的真正福乐，所有爱祂的人都因祂而获得存在和所怀的爱，我们无需担心任何认识祂的人会因祂而失望，祂渴望我们的爱并非为了自身的利益，而是为使爱祂的人获得永恒的赏报，即他们所爱的天主本身——我们应当做什么呢？正因如此，我们甚至爱我们的仇敌。因为我们并不惧怕他们，因为他们不能夺走我们所爱的；相反，我们更怜悯他们，因为他们越是恨我们，就越是远离我们所爱的那一位。因为他们若转向祂，就必定会爱祂为至高之善，也会爱我们为同享如此宏恩的同伴。\par

% structure: p0065 source=h2 target=subsection reason=relative-heading-rank; base=h2

\subsection{第30章——天使是否应被视为我们的近人}

31. 在此，有关天使的问题又进一步出现。因为他们正幸福地享见我们所渴望享见的天主；而我们在今生如同在镜中观看，模糊不清，越是分享享见，就越能忍受我们的流徙，也越热切渴望其终结。但问及在这两条诫命中是否也包括对天使的爱，并非不合理。因为那命令我们爱近人的主，就人而言并未作任何例外，这由我们的主在福音中及宗徒保禄所证明。当主将这两条诫命交付给他、并对他说全法律和先知都系于这两条诫命的那人问祂：\Quote{谁是我的近人？}祂便告诉他有一个人从耶路撒冷下到耶里哥，落在强盗手中，被他们打得半死，剥去衣服丢下。祂向他表明，除了那怜悯他、上前照顾他的那人以外，无人是他的近人。而提问的那人在被反问时也承认了这点。主对他说：\Quote{你去照样做吧。}以此教导我们：凡是我们在其需要中有义务帮助的人，或是如果他有需要我们就应帮助的人，就是我们的近人。由此可推，反过来也有义务帮助我们的人，就是我们的近人。因为\Quote{近人}一名是相对的，无人能成为近人，除非相对于另一位近人。再者，当主将诫命甚至延伸到我们的仇敌时，谁看不出没有一个人被排除在应得仁慈之事之外呢？\Quote{爱你们的仇敌，善待恨你们的人。}\BibleRef{玛 5:44}\par

32. 宗徒保禄也是如此教训的，他说：\Quote{因为\InnerQuote{不可奸淫，不可杀人，不可偷盗，不可作假见证，不可贪婪}，以及别的任何诫命，都包含在这句话里：\InnerQuote{你当爱你的近人如你自己。}爱不加害于近人。}\BibleRef{罗 13:9} 所以，凡认为宗徒未在这条诫命中包括每一个人的人，就不得不承认（这既极荒谬又极有害）宗徒认为若一个人不是基督徒或是仇敌，与他的妻子通奸、或杀他、或贪图他的财物都无罪。既然除了愚人无人会这样说，那么显然每个人都应被视为我们的近人，因为我们不可加害于任何人。\par

33. 然而，如果每个我们应当向其显示或应当向我们显示慈悲服务的人，按权利被称为近人，那么显然，爱近人的诫命也包括圣天使，因为他们在我们身上施行了如此伟大的慈悲服务，这一点可以很容易地通过留意圣经中的许多段落来证明。基于此，连天主自己、我们的主，也愿意被称为我们的近人。因为我们的主耶稣基督以那救助被强盗打伤、半死半活丢在路边之人的比喻，指向祂自己。圣咏作者在祈祷中说：\Quote{我待他们有如朋友或兄弟。}但由于天主性比我们的本性更高贵、远远超越，爱天主的诫命与爱近人的诫命是不同的。因为祂出于自身的良善而怜悯我们，而我们却是因祂而彼此怜悯——也就是说，祂怜悯我们，使得我们可以完全享受祂自己；我们彼此怜悯，使得我们能够完全享受祂。\par

% structure: p0069 source=h2 target=subsection reason=relative-heading-rank; base=h2

\subsection{第三十一章——天主使用我们而非享受我们。}

34. 基于此，当我们说我们只享受我们为其自身而爱的事物，并且除了那使我们幸福的事物之外，没有真正的享受对象，而所有其他事物都是供使用——似乎仍有需要说明之处。因为天主爱我们，圣经也常向我们展示祂对我们的爱。那么祂是如何爱我们的呢？是作为使用对象还是作为享受对象？如果祂享受我们，祂就必然从我们这里需要好处，没有理智的人会这样说；因为我们所享受的一切好处，要么是祂自己，要么来自于祂。没有人会不知道或怀疑，光本身并不需要它自己所照亮之物的照耀。圣咏作者非常明确地说：\Quote{我向上主说：祢是我的天主，因为祢不需要我的善。}因此，祂并不享受我们，而是使用我们。因为如果祂既不享受也不利用我们，我不知道祂如何能爱我们。\par

% structure: p0071 source=h2 target=subsection reason=relative-heading-rank; base=h2

\subsection{第三十二章——天主如何利用人。}

35. 但祂的使用也并非我们那种使用方式。因为当我们使用对象时，目的是为了完全享受天主的良善。然而，天主在使用我们时，则是为了祂自身的良善。因为祂是善的，我们才存在；而只要我们真实地存在，我们就是善的。此外，因为祂也是公义的，所以我们不能作恶而不受惩罚；而我们越是恶，存在就越不完全。祂是首要的、至高的存在，完全不变，能以最充分的意义说：\BibleQuote{我是自有者，}并且\BibleQuote{你要对以色列子民说：那自有者派遣我到你们这里来。}\BibleRef{出 3:14}因此，所有其他存在的事物，其存在完全依赖于祂，并且它们之所以善，仅仅是因为祂赐予它们如此存在。因此，所谓天主使用我们，并非指向祂自己的利益，而仅仅指向我们的利益；就祂而言，仅仅指向祂的良善。当我们怜悯一个人并照顾他时，我们这样做是为他的益处；但不知何故，我们自己的益处也通过某种自然的结果而随之而来，因为天主不会让我们向需要者所施的怜悯得不到赏报。这就是我们最高的赏报：完全享受祂，并且所有享受祂的人都在祂内彼此享受。\par

% structure: p0073 source=h2 target=subsection reason=relative-heading-rank; base=h2

\subsection{第三十三章——人应当如何被享受。}

36. 因为如果我们在一人身上找到完全的幸福，我们就会停在这条路上，将我们幸福的希望寄托于人或天使。傲慢的人和傲慢的天使将这种希望据为己有，并且乐于看到他人依赖他们。但相反，圣善的人和圣善的天使，即使当我们疲惫并渴望与他们同住、安息在他们内时，也会用他们从天主那里为他人或为自己所领受的供应来恢复我们的力量；然后催促我们这些恢复了精力的人继续前行，走向祂，在享受祂时我们找到共同的幸福。因为宗徒大声说：\Quote{保禄曾为你们被钉死在十字架上吗？或者你们曾因保禄的名字受洗吗？}\BibleRef{格前 1:13}又说：\Quote{所以，栽种的算不了什么，浇灌的也算不了什么，只在乎那使之生长的天主。}\BibleRef{格前 3:7}天使也告诫那准备朝拜他的人，应当朝拜他的主人，而他本人与他同是仆人。\BibleRef{默 19:10}\par

37. 但当你因天主而喜悦一人时，你享受的是天主，而不是那人。因为你享受那使你幸福的那一位，并且你因来到祂面前而欢喜，在祂面前放置你喜乐的希望。因此，保禄对费肋孟说：\Quote{是的，弟兄，让我在主内从你得到喜乐。}因为他若是没有加上\InnerQuote{主内}，而只说\InnerQuote{让我从你得到喜乐，}他就会暗示他把幸福的希望寄托在那人身上；尽管即使在上下文中，\Quote{享受}一词是以\Quote{愉快地使用}的意义来使用的。因为当我们所爱的事物临近时，它自然应带来喜悦。如果你超越这喜悦，并使之成为你最终安息的那一目标的途径，那么你是在使用它，而你若说你享受它，则是言语上的滥用。但如果你紧抓住它并安息在其中，在其中找到完全的幸福，那么你可以真正地、恰当地说你享受它。我们决不应这样做，除非是在谈论有福的天主圣三，祂是至高不变之善。\par

% structure: p0076 source=h2 target=subsection reason=relative-heading-rank; base=h2

\subsection{第三十四章——基督是通往天主的首要道路。}

38. 请注意，即使祂本人就是真理和圣言，万物都是藉着祂造成的，却成了血肉，好住在我们中间，宗徒仍然说：\BibleQuote{纵使我们曾按肉性认识基督，但如今不再这样认识祂了。}\BibleRef{格后 5:16} 因为基督不但愿意把产业赐给那些已经走完旅程的人，而且祂自己愿意成为那些刚出发的人的道路，所以祂决定取一个血肉的身体。因此也有那句话说：\BibleQuote{上主在祂道路的起初就造了我}\BibleRef{箴 8:22}，就是说，那些愿意来的人可以在祂内开始旅程。因此，宗徒虽然仍在路上，追随那召叫他承受天上赏报的天主，但忘记背后，努力向前\BibleRef{斐 3:13}，他已经跨越了道路的起点，如今不再需要它了。然而，凡渴望达到真理并在永生中安息的人，都必须在祂开始旅程。因为祂说：\BibleQuote{我是道路、真理、生命}\BibleRef{若 14:6}；就是说，人们藉着我来到，向着我来，在我内安息。因为当我们来到祂面前时，我们也来到圣父面前，因为藉着平等者认识平等者；圣神连系我们，并好像印了我们，使我们能够永恒地安息在那至高的不变之善中。因此我们可以学习，即使我们的主祂自己，就祂屈尊成为我们的道路而言，也不愿留住我们，反而希望我们继续前进，那么还有什么比不被任何事情耽搁在路上更重要的呢？我们不要软弱地紧抓暂时之物，即使这些东西是祂为了我们的救恩而取用和穿戴的，而要迅速越过它们，努力达到祂自己那里，祂已使我们的本性从暂时之物的束缚中解放出来，并把它安置在祂圣父的右边。\par

% structure: p0078 source=h2 target=subsection reason=relative-heading-rank; base=h2

\subsection{第三十五章　圣经的满全与目标是爱天主及爱近人。}

39. 那么，自从我们开始讨论事物以来，所说的一切可以总结为：我们应当清楚明白，法律和全部圣经的满全与目标是爱那当享受的对象，并爱那能与我们一同享受那对象的对象。因为没有人需要命令爱自己。因此，为我们的救恩而设的整个暂时性的安排，是由天主眷顾所制定的，好让我们认识这个真理并能够付诸实行；我们应当使用这个安排，不是带着一种如安息之善的爱与喜悦，而应带着一种暂时性的情感，如同我们对待道路、车辆或其他仅是工具的东西。或许可以找到其他更合适的比喻来表达这个观念：我们爱那承载我们的东西，只是为了我们所趋向的那个目标。\par

% structure: p0080 source=h2 target=subsection reason=relative-heading-rank; base=h2

\subsection{第三十六章　那在爱德中建立我们的圣经解释，即使有错，也不是有害的欺骗或谎言；然而解释者应当被纠正。}

40. 所以，谁若自以为了解圣经或其任何部分，却加上一种不倾向于建立这双重的爱——爱天主和爱近人的解释，他就尚未按应该的方式了解它们。反之，若有人从圣经中得出一种可用于建立爱的意义，即使他并未恰好碰上作者在该处想要表达的精确意义，他的错误并非有害，他也完全免于欺骗的指责。因为欺骗涉及说假话的意图；我们发现很多人意图欺骗，但没有人愿意被欺骗。既然知道的人行骗，而不知道的人被骗，那么显而易见，在任何特定情况下，被骗的人比行骗的人更好，因为受不义比行不义更好。每一个说谎的人就行不义；若有人认为谎言有时有益，他必定认为不义有时有益。没有说谎者在他说谎的事上守信。他当然希望他所欺骗的人信任他，但他却因说谎而辜负了这种信任。每一个背信的人就是不义的。因此，要么不义有时是有益的（这是不可能的），要么谎言从来不是有益的。\par

41. 谁若从圣经取出作者原意以外的其他意义，他就走错了路，但这并非由于圣经有虚假。然而，正如我即将说的，若他的错误解释倾向于建立爱——这爱的命令的满全，他就如同一个人错误地离开了大路，却穿过田野到达了同一地点。然而他应当被纠正，并被告知不离开直路要好得多，免得他养成走岔路的习惯，有时会走上歧途，甚至完全走错方向。\par

% structure: p0083 source=h2 target=subsection reason=relative-heading-rank; base=h2

\subsection{第三十七章　错误解释的危险。}

因为如果他轻率地采纳一种他所读的作者并未意指的含义，他常常会遇到其他他无法与此含义协调的陈述。而如果他承认这些陈述是真实且确定的，那么，他先前赋予前一段落的含义就不可能是真实的；于是，几乎难以说清，出于对自己意见的偏爱，他开始对圣经比对自己更加恼怒。而一旦他容许这种邪恶潜入，它必将彻底毁灭他。\BibleQuote{因为我们凭信德行走，而非凭目睹。}\BibleRef{格后 5:7}如今，如果圣经的权威开始动摇，信德也将动摇。而如果信德动摇，爱德本身也会冷却。因为如果一个人从信德堕落，他也必然从爱德堕落；因为他不能爱他认为不存在的东西。但如果他既相信又爱，那么，通过善工和对道德诫命的殷勤关注，他也将开始希望他终将获得他所爱的对象。因此，所有知识与所有预言都服务于这三样：信德、望德与爱德。\par

% structure: p0085 source=h2 target=subsection reason=relative-heading-rank; base=h2

\subsection{第三十八章．爱德永不止息。}

42. 但目睹将取代信德；望德将被吞没于我们将要到达的那完美福乐中；而爱德，当其他二者消失时，反而会增长。因为如果我们借着信德爱那尚未看见的，当我们开始看见时，我们将何等更爱它！如果我们借着望德爱那尚未达到的，当我们达到时，我们将何等更爱它！因为暂时之物与永恒之物之间有一个重大区别：暂时之物在拥有之前更被珍视，一旦获得便开始显得毫无价值，因为它不能满足灵魂——灵魂只有在永恒中才有真正而可靠的安息；相反，永恒之物在拥有时比在渴望时更被热烈地爱着，因为没有人能在渴望时对其赋予比其真正价值更高的评价，以至于当发现它并不像所想象的那样有价值时就认为它相对无价值；相反，无论一个人在走向拥有它的路上对其评价多高，当他实际拥有它时，会发现它的价值更高。\par

% structure: p0087 source=h2 target=subsection reason=relative-heading-rank; base=h2

\subsection{第三十九章．在信德、望德与爱德上成熟的人不再需要圣经。}

43. 因此，一个安息于信德、望德与爱德并紧握这三者的人，除了为教导他人之外，并不需要圣经。因此，许多人即使没有圣经抄本，甚至独处时，也因着这三样恩宠而生活。所以，我认为在他们身上，这句话已经应验：\BibleQuote{至于预言，它们终必消失；至于语言，它们终必停止；至于知识，它终必消逝。}\BibleRef{格前 13:8}然而，借着这些工具（姑且如此称呼），如此伟大的信德与爱德之殿宇在他们内被建立起来，以至于他们持守那完美的，并不寻求那仅在部分上完美的——当然，我是就今生所可能而言；因为与来世相比，今生没有一位义人或圣人的生命是完美的。因此，宗徒说：\BibleQuote{如今常存的有信德、望德、爱德，这三样；但其中最大的是爱德：}\BibleRef{格前 13:13}因为当人到达永恒世界时，虽然其他两种恩宠将消失，爱德将更伟大、更稳固。\par

% structure: p0089 source=h2 target=subsection reason=relative-heading-rank; base=h2

\subsection{第四十章．圣经要求怎样的读者。}

44. 因此，如果一个人完全理解\BibleQuote{诫命的终向是爱德，出自纯洁的心、良善的良心和无伪的信德}\BibleRef{弟前 1:5}，并决心使他一切对圣经的理解都服务于这三样恩宠，那么他就可以放心地去解释这些书卷。因为宗徒说到\Quote{爱德}时，又加上\Quote{出自纯洁的心}，以防止人爱上任何不应爱的事物。他又加上\Quote{良善的良心}，关乎望德；因为如果一个人背负着不良良心的重担，他会绝望于达到他所相信和所爱的事物。第三，他说：\Quote{和无伪的信德}。因为如果我们信德毫无虚伪，那么我们既会戒绝爱那不应爱的事物，又会因正直生活而能怀抱希望，不致使我们的望德落空。\par

为此，我一直渴望谈论信德的对象，就我按当前目的认为必要的而言；因为关于这一主题，他人或我自己已在其他书卷中说过很多。因此，就让这作为本书的结束。在下一卷中，我将尽天主所赐的光照，讨论记号的课题。\par

\SourceRecord{Translated by James Shaw.；Nicene and Post-Nicene Fathers, First Series；Vol. 2.；Edited by Philip Schaff.；Buffalo, NY: Christian Literature Publishing Co.,；1887.；<http://www.newadvent.org/fathers/12021.htm>.}

% structure: p0001 source=h1 target=section reason=page-title

\section{论基督圣道（卷二）}

% structure: p0002 source=h2 target=subsection reason=relative-heading-rank; base=h2

\subsection{第一章　记号：其性质与多样性}

1. 正如我在论述事物时，曾预先警告不要关注事物本身之外的东西，即使它们是某个其他事物的记号；同样，现在轮到讨论记号时，我定下这条原则：不要关注记号本身，而要注意它们是记号，即它们所标示的东西。因为记号是一种事物，除了在感官上产生的印象外，它本身还导致某物进入心中：例如，当我们看见足迹时，我们断定留下这个足迹的动物经过；当我们看见烟时，我们知道烟下有火；当我们听见活人的声音时，我们想到他心中的情感；当号角吹响时，兵士们知道该前进或后退，或做其它战场所需的事。\par

2. 有些记号是自然的，有些是约定的。自然的记号是那些无需任何使用记号作为记号的意图或愿望，却能引导认识其他事物的记号，例如烟指示火。因为烟并非有意成为记号，而是通过经验观察，我们即使在仅能看见烟时也知道下面有火。同样，动物经过的足迹也属此类记号。愤怒或忧伤之人的面容，无论其意愿如何，都揭示出心中的情感；同样，心灵的其他任何情绪都会通过传情达意的面容泄露出来，即使我们无意让人知道。不过，这类记号目前不在我的讨论计划之内。但由于它属于这一主题的划分，我不能完全忽略它。做到这一点就够了。\par

% structure: p0005 source=h2 target=subsection reason=relative-heading-rank; base=h2

\subsection{第二章　关于我们现在所关注的记号种类}

3. 另一方面，约定的记号是活物为了尽可能表现他们心中的情感、感知或思想而相互交换的。给出记号的唯一理由就是希望将记号发出者心中的东西引出并传达到另一个人心中。因此，我们希望就与人有关的这类记号进行考虑和讨论，因为即使是天主赐予我们、包含在圣经中的记号，也是通过人——即那些撰写圣经的人——传达给我们的。野兽之间也有某些记号，用以表明它们心中的愿望。例如，公鸡发现食物时，会用叫声信号让母鸡跑来；鸽子通过咕咕叫唤它的配偶，或被配偶呼唤；许多类似的记号为常见之事。现在，这些记号——如同人的表情或悲叹——是本能地、毫无目的地随心灵活动而产生，还是确实有目的地用于表意，是另一个问题，与当前主题无关。我排除这部分，因为对于当前目的来说不必要。\par

% structure: p0007 source=h2 target=subsection reason=relative-heading-rank; base=h2

\subsection{第三章　在记号中，言语占据首要地位}

4. 那么，在人用以相互交流思想的记号中，有些与视觉有关，有些与听觉有关，很少与其它感官有关。因为当我们点头时，我们仅对希望借此传达意愿之人的眼睛给出记号。手部的动作也能传达大量信息；演员通过全身的动作向懂行的人给出某些记号，可以说是向眼睛进行对话；军事旗帜和旗语通过眼睛传达指挥官的意志。所有这些记号都像是一种可见的言语。而诉诸听觉的记号，如我所说，更为众多，且大多数由言语构成。因为号角、长笛和竖琴不仅发出悦耳的声音，也常常发出有意义的声音，但所有这些记号与言语相比却极其稀少。因为在人类中，言语作为表达思想的方式，占据了绝对主要的地位。诚然，我们的主曾通过倒在祂脚上的香膏气味给出记号；在祂身体和血的圣事中，祂通过味觉显明祂的旨意；当妇人触摸祂衣服的边而痊愈时，这一行动也不乏表征意义。\BibleRef{玛 9:20} 但人表达思想的无数记号都由言语组成。因为我能够用言语表达我简要提到的各种记号，却无法用那些记号来表达言语。\par

% structure: p0009 source=h2 target=subsection reason=relative-heading-rank; base=h2

\subsection{第四章　文字的起源}

5. 然而，由于言语一旦撞击空气就消逝，其声音存在时间不长，人们便通过字母形成言语的记号。这样，声音通过某些记号变得肉眼可见，当然不是作为声音本身。然而，由于人类纷争之罪——源于每个人都试图为自己夺取首位——使这些记号无法成为所有民族共通。而那座通天的高塔就是这种傲慢精神的表现；参与其中的不虔敬之人理应受到惩罚，不仅他们的心灵，连他们的舌头也陷入混乱与不和。\EditorialNote{创世纪第十一章}\par

% structure: p0011 source=h2 target=subsection reason=relative-heading-rank; base=h2

\subsection{第五章　圣经被译为多种语言}

6. 因此，就连那为人类意志的可怕疾病带来医治的\WorkTitle{圣经}，起初以一种语言写成，以便在适当的时候传遍全世界，而后被翻译成各种语言，广传四方，从而为各民族的救恩所认识。人在阅读它时，所追求的不过是找出作者的心思和旨意，并由此找出天主的旨意，他们相信这些作者是凭着天主的旨意说话的。\par

% structure: p0013 source=h2 target=subsection reason=relative-heading-rank; base=h2

\subsection{第六章——论\WorkTitle{圣经}中因喻象语言而产生的晦涩之益处。}

7. 然而，轻率粗心的读者会被许多各样的晦涩与歧义引入歧途，用一个意义代替另一个；在某些地方，他们甚至无法得出合理的解释。有些表达如此晦涩，仿佛将含义笼罩在最浓密的黑暗中。我不怀疑，这一切都是天主上智的安排，为的是以辛劳制伏骄傲，并防止理智产生厌腻之感——理智通常会贬低不费力就发现的事物。因为，我问，为什么如果有人这样说：有些圣善正义的人，他们的生活和交往被基督的教会用来救赎那些从各种迷信中来到教会的人，并通过他们效法善人而使他们成为教会身体的肢体；这些人作为天主善良忠实的仆人，来到圣洗池，卸下世界的重担，从那里上升，因圣神的植入，结出双重的爱，即爱天主和爱近人的果子——我说，如果一个人这样说，为什么他不能让听众感到像他从\BibleBook{}{雅歌}中引用同一含义时那样愉快？那里论到教会，以一位美丽女子的形象受到赞美：\BibleQuote{你的牙齿像一群剪过毛的母羊，从洗濯中上来，个个都怀双胎，没有一只是不生育的。}\BibleRef{歌 4:2} 难道听众所学到的，比他用最直白的语言表达同一思想而不用这个形象时更多吗？然而，不知何故，当我看见圣人被视作教会的牙齿，把人们从错误中撕扯出来，带进教会的身体，他们的一切粗鲁被软化，就像被牙齿撕碎咀嚼过一样，我默观他们时感受到更大的喜悦。我也以最大的喜悦认出他们是在剪过毛的绵羊的形象下，卸下世界的重担如同羊毛，从洗濯（\Emph{即}圣洗）中上来，个个都怀双胎，\Emph{即}爱的双重诫命，没有一只在那圣洁的果实中是不生育的。\par

8. 但我为什么在那种形象下默观他们时，比没有从\WorkTitle{圣经}中得出这样的形象时更为喜悦，尽管事实相同，知识也相同，这是另一个问题，而且非常难以回答。然而，没有人怀疑这样的事实：在某些情况下，通过形象获得知识更加愉快，而寻求时的困难会带来发现时的更大喜悦。——因为那些寻求却找不到的人遭受饥饿。而那些根本不寻求，只因所需就在身旁而拥有的人，常常因饱足而倦怠。这两种情况导致的软弱都应避免。因此，圣神以令人钦佩的智慧和对我们福祉的关怀，如此安排了\WorkTitle{圣经}：以较明晰的章节满足我们的饥饿，以较晦涩的章节激发我们的胃口。因为几乎没有什么从那些晦涩段落中挖掘出来的东西，不能用最直白的语言在其他地方找到。\par

% structure: p0016 source=h2 target=subsection reason=relative-heading-rank; base=h2

\subsection{第七章——达至智慧的步骤：第一，敬畏；第二，虔敬；第三，知识；第四，决心；第五，谋略；第六，心的净化；第七，终点或止境，智慧。}

9. 首先，我们必须因对天主的\Emph{敬畏}而受引导，去寻求认识祂的旨意，祂命令我们渴望什么、避免什么。这敬畏必将激发我们想到自己的必死性和面临的死亡，并将一切骄傲的冲动钉在十字架上，仿佛我们的肉身被钉在木头上。其次，我们的心必须为\Emph{虔敬}所驯服，不可违抗\WorkTitle{圣经}，无论它被理解时如何攻击我们的一些罪过，或者当我们不理解时，觉得自己更有智慧，能发出更好的命令。我们必须宁可相信：无论那里写着什么，即使隐藏着，也胜过我们凭自己的智慧所能想出的任何东西。\par

10. 在这两步――敬畏与虔敬――之后，我们来到第三步，即\Emph{知识}，我现在正要讨论它。因为在此，每一位认真研读圣经的人都要操练自己，旨在圣经中只寻找这一目的：天主应为祂自身而受爱，近人应为天主而受爱；并且当以全心、全灵、全意爱天主，并爱近人如己——也就是说，我们所有对近人的爱，如同我们对自己的爱，都当指向天主。\BibleRef{玛 22:37}我在前一卷讨论事物时已提及这两条诫命。因此，每个人必须先在圣经中发现，由于他纠缠于对这个世界——即暂时之物——的爱，已远离了圣经所命令的那种对天主和对近人的爱。然后，那引导他思考天主审判的敬畏，以及那让他别无选择只得相信并服从圣经权威的虔敬，迫使他哀叹自己的处境。因为良好希望的知识不是使人自夸，而是使人忧伤。怀着这种心境，他恒切祈祷，求得天主助佑的安慰，以免在绝望中被压垮，于是逐渐进入第四步——即\Emph{力量}和\Emph{决心}——在此他饥渴慕义。因为在此心境中，他摆脱了各种对短暂事物致命的喜乐，转而将情感专注于永恒之物，即那不变而合一的天主圣三。\par

11. 当他竭尽全力凝视从远处闪耀的这一对象，并因自己视力的软弱而无法承受那无比的光芒时，他就在第五步——即在\Emph{慈悲的劝告}——中洁净自己那被粗暴搅动、被卑鄙欲望困扰的灵魂，除去它所沾染的污秽。在这一阶段，他勤勉地操练爱近人；当他达到爱仇敌的地步时，满怀希望，力量不衰，他便登上第六步，在此他\Emph{净化那能看见天主的眼睛本身}，\BibleRef{玛 5:8}达到天主能被那些尽可能死于这个世界的人看见的程度。因为人看见祂，恰恰是当他们向这个世界死去；而他们向这个世界活着时，就看不见祂。然而，尽管那光开始显现得更清晰，不仅更可忍受，甚至更令人喜悦，但我们仍说我们是在镜子中模糊地观看，因为我们行走是因信德，而非因目睹，我们在此世漂泊为客，即使我们的交往是在天上。在这一阶段，人如此净化他情感的眼光，以致不将近人置于真理之前，甚至不与真理等同，因此也不将自己置于真理之前，因为他不爱那爱他如己的人。因此，那位圣洁的人将如此专一、心地纯洁，以致他不会偏离真理，无论是为取悦人，还是为避免此生所遭遇的任何烦恼。这样的子女上升到\Emph{智慧}，这是第七步，也是最后一步，他在平安和宁静中享受这智慧。因为敬畏天主是智慧的开端。那么，从那开端直到达致智慧本身，我们的道路就是通过现在所描述的这些步骤。\par

% structure: p0020 source=h2 target=subsection reason=relative-heading-rank; base=h2

\subsection{第八章：正典书卷}

12. 但现在让我们回到刚才提到的第三步，因为我已立定心志，要照着主赐予我的智慧来谈论和推理它。那么，最熟练的圣经诠释者，将是这样的人：他首先阅读了全部圣经，并将其存记于心，尽管尚未完全理解，但至少拥有阅读所能给予的那种知识——至少那些被称为\Emph{正典}的书卷。因为当他在对真理的信仰中被建立起来后，他将更安全地阅读其他书卷，免得它们首先占据一颗软弱的心灵，并以危险的虚假和妄想欺骗它，使之充满不利于健全理解的偏见。至于正典圣经，他必须遵循多数天主教会的判断；其中，当然要高度尊重那些被认为堪当宗徒之座并接纳书信的教会。因此，在正典圣经中，他将按以下标准判断：宁取所有天主教会都接纳的书卷，胜过一些教会不接纳的。再者，凡非所有教会都接纳者，他宁取获得更多教会与更高权威认可者，胜过较少教会与较低权威所持守者。然而，若他发现某些书卷为多数教会所持守，另一些则为更高权威的教会所持守（尽管这事不大可能发生），我认为在此情形下，双方的权威应视为相等。\par

13. 如今，我们说要行使判断的全部圣经正典，包含在以下书卷中：梅瑟五书，即\BibleBook{}{创世纪}、\BibleBook{}{出谷纪}、\BibleBook{}{肋未纪}、\BibleBook{}{户籍纪}、\BibleBook{}{申命纪}；农之子若苏厄书一卷；\BibleBook{}{民长纪}一卷；一短卷名\BibleBook{}{卢德传}，它似乎更属于列王纪的开端；其次，列王纪四卷，即\BibleBook{}{撒慕尔纪上}、\BibleBook{}{撒慕尔纪下}、\BibleBook{}{列王纪上}、\BibleBook{}{列王纪下}；编年纪两卷，即\BibleBook{}{编年纪上}、\BibleBook{}{编年纪下}——这两卷并非前后衔接，而是平行叙述，可以说覆盖同样的内容。刚才提到的书卷是历史书，包含连续的时代叙事，遵循事件顺序。另有其他书卷似乎没有固定顺序，与前书顺序无关，彼此也无关，如\BibleBook{}{约伯传}、\BibleBook{}{多俾亚传}、\BibleBook{}{艾斯德尔传}、\BibleBook{}{友弟德传}、\BibleBook{}{玛加伯上}、\BibleBook{}{玛加伯下}、\BibleBook{}{厄斯德拉上}、\BibleBook{}{厄斯德拉下}，最后这两卷更像是连续正规历史的续篇，该历史以列王纪和编年纪终结。接下来是先知书，其中有\BibleBook{}{圣咏集}一卷；撒罗满三书，即\BibleBook{}{箴言}、\BibleBook{}{雅歌}、\BibleBook{}{训道篇}；因为两卷书，一名\BibleBook{}{智慧篇}、一名\BibleBook{}{德训篇}，因风格类似而被归诸撒罗满，但更可能的意见是它们由息辣之子耶稣所写。不过它们仍应算在先知书内，因为它们已被承认为具有权威。其余是严格称为先知书的书卷：十二卷独立的先知书，彼此相连，从未被分开，故算为一卷；这些先知的名字如下：\BibleBook{}{欧瑟亚}、\BibleBook{}{岳厄尔}、\BibleBook{}{亚毛斯}、\BibleBook{}{亚北底亚}、\BibleBook{}{约纳}、\BibleBook{}{米该亚}、\BibleBook{}{纳鸿}、\BibleBook{}{哈巴谷}、\BibleBook{}{索福尼亚}、\BibleBook{}{哈盖}、\BibleBook{}{匝加利亚}、\BibleBook{}{玛拉基亚}；然后有四大先知：\BibleBook{}{依撒意亚}、\BibleBook{}{耶肋米亚}、\BibleBook{}{达尼尔}、\BibleBook{}{厄则克耳}。旧约的权威包含在这四十四卷书的范围内。新约的权威则包含在以下书卷中：福音书四卷，即\BibleBook{玛窦}{福音}、\BibleBook{马尔谷}{福音}、\BibleBook{路加}{福音}、\BibleBook{若望}{福音}；宗徒保禄的书信十四封——\BibleBook{}{罗马书}、\BibleBook{}{格林多前书}、\BibleBook{}{格林多后书}、\BibleBook{}{迦拉达书}、\BibleBook{}{厄弗所书}、\BibleBook{}{斐理伯书}、\BibleBook{}{得撒洛尼前书}、\BibleBook{}{得撒洛尼后书}、\BibleBook{}{哥罗森书}、\BibleBook{}{弟茂德前书}、\BibleBook{}{弟茂德后书}、\BibleBook{}{弟铎书}、\BibleBook{}{费肋孟书}、\BibleBook{}{希伯来书}；\BibleBook{}{伯多禄前书}、\BibleBook{}{伯多禄后书}；\BibleBook{}{若望一书}、\BibleBook{}{若望二书}、\BibleBook{}{若望三书}；\BibleBook{}{犹达书}；\BibleBook{}{雅各伯书}；\BibleBook{}{宗徒大事录}一卷；\BibleBook{}{若望默示录}一卷。\par

% structure: p0023 source=h2 target=subsection reason=relative-heading-rank; base=h2

\subsection{第九章 我们应如何研读圣经}

14. 在所有这些书卷中，敬畏天主并怀有温良虔敬之心的人寻求天主的旨意。在进行这项探索时，要遵守的首要规则是，如我所说，了解这些书卷，即使尚未理解，仍要阅读它们以便记在脑中，或至少不致对它们完全不知。其次，那些在其中明确陈述的事，无论是生活的规则还是信德的规则，都要更仔细、更勤勉地探究；一个人发现得越多，他的理解力就越强。因为在圣经中明确陈述的事中，可以找到所有关乎信德和生活方式的真理——即望德和爱德，这些我已在前书中谈论过了。此后，当我们对圣经的语言有了一定程度的熟悉后，就可以着手解开和研究晦涩的段落，并在此过程中从较明确的表述中提取例子以照亮较晦涩的段落，并利用无争议的段落的证据消除对可疑段落的任何犹豫。在这件事上记忆的作用很大；但如果记忆有缺陷，没有任何规则可以弥补。\par

% structure: p0025 source=h2 target=subsection reason=relative-heading-rank; base=h2

\subsection{第十章 未知或含混的记号阻碍圣经的理解}

15. 现在，阻止所写内容被理解的原因有两个：被掩盖在未知的记号下，或被掩盖在含混的记号下。记号要么是本义的，要么是比喻义的。当我们用它们来指向它们原初设计指向的对象时，它们被称为本义的，例如我们说\Emph{bos}，意指一头牛，因为所有与我们一样使用拉丁语的人都这样称呼它。当那些我们用本义名称所指的事物本身被用来意指别的东西时，记号就是比喻义的，例如我们说\Emph{bos}，并通过这个音节理解通常被称为牛的动物；然后更进一步，通过这头牛理解福音的宣讲者，正如圣经按照宗徒的解释所表明的，当它说：\BibleQuote{你不可笼住正在打谷的牛的嘴。}\par

% structure: p0027 source=h2 target=subsection reason=relative-heading-rank; base=h2

\subsection{第十一章 语文知识，尤其是希腊文和希伯来文的知识，为消除对记号的不知是必要的}

16. 矫正对正确符号无知的最大良药是语言的知识。操拉丁语的人——即我所要教导的那些人——为了认识圣经，还需要另外两种语言，即希伯来语和希腊语，以便当拉丁译者的无穷差异性使他们产生疑惑时，他们可以求助于原文。尽管我们确实发现圣经中常有一些希伯来语词未被翻译，例如\Quote{阿们}、\Quote{阿肋路亚}、\Quote{辣哈}、\Quote{贺三纳}等同类词语。其中一些词，虽然本可翻译，却因其更神圣的权威而保留了原形，例如\Quote{阿们}和\Quote{阿肋路亚}。另一些词则被认为无法译成另一种语言，我提到的另外两个词便属此类。因为在某些语言中，有些词无法译成另一种语言的习语。这尤其发生在感叹词上，这类词更多地表达心灵的情感，而非思想中的任何部分。所举的两个词据说属于此类：\Quote{辣哈}表达愤怒之人的呼喊，\Quote{贺三纳}表达喜乐之人的呼喊。但掌握这些语言是必要的，并非为了少数这样的词——它们很容易标记并询问——而是如前所述，因为译者的差异性。因为从希伯来语译成希腊语的圣经译本可以数得过来，但拉丁译者却数不胜数。因为在信仰的初期，每个偶然得到希腊文手稿、并自认为对两种语言稍有知识的人，都冒险从事翻译工作。\par

% structure: p0029 source=h2 target=subsection reason=relative-heading-rank; base=h2

\subsection{第十二章——多种解释有益。由含混词语产生的错误。}

17. 如果读者不是粗心大意，这种情况反而会有助于而非妨碍对圣经的理解。因为对多种文本的查考常常阐明了一些较隐晦的段落；例如，在依撒意亚先知的那段经文中，一位译者读作：\Quote{不可轻视你苗裔的家人}；另一位读作：\Quote{不可轻视你自己的肉}。这两者相互证实。因为一个被另一个所解释；因为\Quote{肉}可取其字面意义，使人明白他受劝诫不可轻视自己的身体；而\Quote{你苗裔的家人}可比喻地理解为基督徒，因为他们由与我们同一的苗裔——即圣言——属灵地出生。当两位译者的意思比较起来时，一个更可能的词义浮现出来，即命令是不要轻视我们的亲属，因为当人将\Quote{你苗裔的家人}与\Quote{肉}联系起来时，亲属最自然地出现在脑海中。因此，我认为宗徒那话是此意，他说：\BibleQuote{或者我能激起我的肉身的同胞发愤，救他们一些人。}\BibleRef{罗 11:14}——即藉着嫉妒那些已经相信的人，他们中有些人也能相信。他称犹太人为他的\Quote{肉}，是由于血统关系。又如，同一位先知依撒意亚的那段话：\Quote{你们若不相信，必不得了解}，另一位译作：\Quote{你们若不相信，必不得存留}。两者中哪一个为直译，若不参照原文就无法确定。然而对有知识的读者来说，两者都含有伟大的真理。因为译者们很难分歧到在一点上毫无交集。因此，在这里，既然了解在于看见且是持久的，而信德在我们尚处于时间之事的摇篮中如同婴孩时用奶喂养我们（\BibleQuote{因为我们现今是藉着信德行走，并非藉着目睹}——\BibleRef{格后 5:7}），此外，除非我们藉信德行走，我们绝不会达至那不会过去而是持久的看见，我们的了解因持守真理而得以净化——因此一人说：\Quote{你们若不相信，必不得了解}；另一人说：\Quote{你们若不相信，必不得存留}。\par

18. 很常见的是，一位译者，对意思不很清楚时，被原文中的含混所欺骗，而将一段话完全置于与作者意思相异的构造中。例如，有些文本读作：\BibleQuote{他们的脚是\Emph{锐利的}，以便流血}\BibleRef{罗 3:15}；因为希腊语中\OriginalTerm{锐利的}{\GreekText{ὁζύς}}一词既意为\Emph{锐利的}也意为\Emph{迅速的}。因此那位译作\Quote{他们的脚是迅速的，以便流血}的人见到了真义。另一位则因取了含混词错误的意思而陷于错误。这样的翻译并非含混，而是错误的；两者之间有极大差别。因为我们学习并非去解释，而是去纠正这类文本。出于同样原因，因为希腊词\OriginalTerm{牛犊}{\GreekText{μόσχος}}意为牛犊，有些人未理解\OriginalTerm{嫩枝}{\GreekText{μοσχεύματα}}\BibleRef{智 4:3}是树的嫩枝，而译作\Quote{牛犊}；这一错误已渗入众多文本，以致你几乎找不到其他写法。然而意思非常清楚；因为随后的词句显明了这一点。因为\Quote{奸夫的栽种不会扎根深厚}比\Quote{牛犊}更合适；因为牛犊用脚行走于地面，并非由根固定于土中。在此段中，事实上，上下文其余部分也证明了这一翻译。\par

% structure: p0032 source=h2 target=subsection reason=relative-heading-rank; base=h2

\subsection{第十三章——如何修正错误的解释。}

19. 然而，因为我们若不查阅译者们所译的原语言，就无法清楚看到各位译者按其各自能力和判断所努力表达的实际思想；又因为译者若非学识渊博，往往偏离作者原意，所以我们必须努力掌握圣经译成拉丁文时所依据的那些语言，或必须获取那些紧扣原文字面的译本——并非因为它们足够，而是因为我们可以用它们来纠正其他译者在其翻译中选择既遵循词句又遵循含义的自由或错误。因为不仅个别词语，而且整个短语有时被译成拉丁文，而任何一位希望遵循说拉丁文的前辈用法的人根本无法用拉丁语习语翻译它们。虽然这些有时不妨碍对段落的理解，但它们冒犯了那些在事物本身甚至其符号保持纯洁时感到更大喜悦的人。因为所谓的\OriginalTerm{语法错误}{solecism}无非是按照与那些有权威的前辈不同的规则组合词语。无论我们说\Quote{\Emph{inter homines}}（在人们中）还是\Quote{\Emph{inter hominibus}}，对于一个只想知道事实的人来说都无关紧要。同样，\OriginalTerm{发音错误}{barbarism}除了以与说拉丁文的前辈不同的方式发音一个词之外又是什么？因为\OriginalTerm{宽恕}{ignoscere}（宽恕）这个词的第三音节应读长还是读短，对于一个以任何方式恳求天主宽恕其罪过的人来说，并不太重要。那么，言语的纯洁除了保留由前人权威所确立的语言习俗之外，又是什么呢？\par

20. 人在这类事情上越软弱，就越容易被冒犯；而他们之所以软弱，是因为他们希望显得有学问，不是在那些有助于建立德性的事物的知识上，而是在符号的知识上——这符号的知识难以不令人骄傲，\BibleRef{格前 8:1}——因为甚至事物的知识也常常会使我们昂首挺胸，若不是被我们师傅的轭所压制的话。因为以下这段经文如此表达，又如何妨碍我们理解呢？\Quote{\Emph{Quæ est terra in quo isti insidunt super eam, si bona est an nequam; et quæ sunt civitates, in quibus ipsi inhabitant in ipsis?}}我更倾向于认为这不过是另一种语言的习语，而非有任何更深的意义。同样，那个我们如今无法从歌唱者口中除去的短语：\Quote{\Emph{Super ipsum autem floriet sanctificatio mea}}，肯定未损及意义。然而，一个更有学问的人会更喜欢将其改正，说\Quote{\Emph{florebit}}而不是\Quote{\Emph{floriet}}。除了歌唱者的习惯之外，没有任何东西阻碍这一改正。这类错误，如果一个人不选择完全避免，很容易忽视它们，因为它们不妨碍正确的理解。但另一方面，看看宗徒的话：\Quote{\Emph{Quod stultum est Dei, sapientius est hominibus, et quod infirmum est Dei, fortius est hominibus}}。如果有人在此处保留希腊语习语，说\Quote{\Emph{Quod stultum est Dei, sapientius est hominum et quod infirmum est Dei fortius est hominum}}，一个敏捷而细心的读者确实会努力达至真实含义，但一个较迟钝的人要么完全不懂，要么会得出完全错误的理解。因为这种表达方式不仅在拉丁语中有语法错误，而且模棱两可，仿佛其含义可以是：人的愚拙或人的软弱比天主的更聪明或更强。但甚至\Quote{\Emph{sapientius est hominibus}}（比人更聪明）这一表达，即使没有语法错误，也并非没有歧义。因为\Quote{\Emph{hominibus}}是作为与格复数还是作为夺格复数，除非参照含义，否则无法辨明。因此，最好说\Quote{\Emph{sapientius est quam homines}}和\Quote{\Emph{fortius est quam homines}}。\par

% structure: p0035 source=h2 target=subsection reason=relative-heading-rank; base=h2

\subsection{第十四章——如何发现未知词语和习语的含义。}

21. 关于歧义的记号，我会在后面讨论。目前我在处理未知的记号，就词语而言，有两种情况。因为读者不认识某个词或习语，就会停滞不前。如果这些属于外语，我们必须向说这些语言的人询问，或者如果有闲暇，我们自己学习这些语言，或者查阅并比较多种译本。然而，如果我们不熟悉的是我们自己语言中的词或习语，我们通过习惯阅读或聆听逐渐认识它们。没有什么比将那些我们不知道意义的词和短语记在记忆中更好的了，这样当我们遇到一个更有学问的人可以询问时，或者遇到一个通过前文或后文或两者显示我们所不知短语的力量和意义的段落时，我们可以借助记忆轻松地关注此事并了解一切。然而，习惯的力量如此之大，甚至在学问方面也是如此，那些以某种方式在圣经研读中滋养和长大的人，会对其他形式的语言感到惊讶，认为它们的拉丁语不如他们从圣经中学到的拉丁语纯正，而这些圣经中的拉丁语在拉丁作家中是找不到的。在这方面，众多的译本也提供了非常大的帮助，如果仔细比较和讨论它们的文本。只是必须消除所有明显的错误。对于那些渴望知道的人，圣经首先应该运用他们校正文本的技能，使未校正的文本让位于校正的文本，至少当它们是同一译本的抄本时。\par

% structure: p0037 source=h2 target=subsection reason=relative-heading-rank; base=h2

\subsection{第十五章——在各译本中，优先选用\OriginalTerm{七十贤士译本}{Septuagint}和\OriginalTerm{意大利古译本}{Itala}。}

22. 现在，在译本本身中，意大利文（\Emph{Itala}）应优先于其他译本，因为它更贴近原文，同时无损于表达的清晰度。为了校正拉丁文，我们必须使用希腊文译本，其中七十贤士译本的权威在旧约方面是卓越的；因为所有更有学识的教会都报道说，七十位译者在他们的翻译工作中如此充满圣神的临在和能力，以致于在那群人中间只有一种声音。而且，如果正如所报道的，以及许多值得信赖的人所断言的那样，他们在翻译工作中被分开，每人各自在一间小屋里，然而没有在任何一人的手稿中发现任何与其他所有手稿在相同词语和相同词语顺序中不同的内容，那么谁敢将任何东西与这样的权威相比，更不用说偏好什么了？即使他们共同商议，结果是从他们共同的劳动和判断中产生了全体一致的意见，即使如此，任何人，无论他多么有经验，也不应该或不适合去纠正许多可敬和有学识之人的一致意见。因此，即使在原始希伯来文中发现了与这些人所表达的不同形式的东西，我认为我们必须顺服天主眷顾的安排，祂使用这些人来实现：犹太民族因宗教顾忌或嫉妒而不愿让其他国家知道的书籍，在托勒密王权力的帮助下，很早就让将来要信主的国家知道了。因此，他们有可能以圣神（祂在他们内工作，并赐给他们同一声音）认为最合适外邦人的方式翻译。但是，尽管如此，正如我上面所说，那些最贴近原文的译者的比较也常常有助于澄清意思，并非没有价值。因此，旧约的拉丁文文本，正如我正要说的，如果必要的话，应根据希腊文的权威来校正，尤其是那些尽管有七十位却被称为以同一声音翻译的希腊文权威。至于新约的书籍，如果因拉丁文文本的多样性而产生任何困惑，我们当然必须屈服于希腊文，尤其是那些在更有学识和研究的教会中发现的希腊文。\par

% structure: p0039 source=h2 target=subsection reason=relative-heading-rank; base=h2

\subsection{第十六章——语言和事物的知识有助于理解比喻性表达。}

23. 至于比喻性记号，如果读者因不认识其中任何一个而停滞不前，其意义应部分通过语言知识，部分通过事物知识来追寻。例如，\Place{史罗亚池}——我们的主曾用唾沫和泥抹在那人眼睛上，并吩咐他去那里洗——就具有比喻性意义，并且无疑传递着秘密的含义；但是，如果那位福音作者没有解释那个名字（见\BibleRef{若 9:7}），如此重要的意义就不会被注意到。同样，我们毫不怀疑，许多希伯来名字没有被那些书的作者解释，如果有人能解释它们，将在解决圣经的谜题方面大有价值和服务。许多精通该语言的人已经为后人带来了不小的益处，他们解释了所有这些词语，而不参照它们在圣经中的位置，告诉我们\Person{亚当}的意思、\Person{厄娃}的意思、\Person{亚巴郎}的意思、\Person{梅瑟}的意思，以及地名：\Place{耶路撒冷}、\Place{熙雍}、\Place{西乃}、\Place{黎巴嫩}、\Place{约但河}，以及在该语言中我们不熟悉的任何其他名字。当这些名字被研究和解释后，圣经中的许多比喻性表达就变得清晰了。\par

24. 对事物的无知也会使比喻性的表达变得晦涩难懂，例如我们不知道圣经中经常用来作比较的动物、矿物或植物的性质。比如，关于蛇的一个众所周知的事实——为了保护头部，它会将整个身体暴露给攻击者——这为理解主的命令带来了何等的光照：我们要像蛇一样机警；\BibleRef{玛 10:16}也就是说，为了我们的头——基督，我们应甘愿将身体交给迫害者，以免因保全身体而否认天主，从而使基督信仰在我们内被摧毁。或者，另一个说法：蛇通过狭窄的洞穴挤压自己，脱去旧皮，从而获得新力量——这与效法蛇的机敏、脱去旧人（如宗徒所说）以便穿上新人的教导是多么吻合；\BibleRef{弗 4:22}并且也要通过狭窄之处脱去旧人，正如主所说：\BibleQuote{你们要从窄门进去！}\BibleRef{玛 7:13}因此，对蛇本性的认识为圣经从这动物引申的许多隐喻提供了光照；同样，对其他同样经常用作比较的动物的无知，对读者来说是极大的障碍。对于矿物和植物也是如此：例如，对红玉髓的认识——它在黑暗中发光——照亮了书中许多以隐喻方式使用的晦涩之处；而对绿柱石或金刚石的无知则常常关闭知识之门。我们之所以容易理解鸽子带回方舟的橄榄枝象征永久的和平，\BibleRef{创 8:11}唯一的原因是我们知道橄榄油的柔滑不易被其他液体破坏，而且橄榄树本身是常青的。此外，许多人由于对牛膝草的无知，不知道它洁净肺部的功效，也不知道它据称能穿透岩石（尽管它是一种矮小不起眼的植物），因此无法理解为什么经文说：\BibleQuote{求你用牛膝草洁净我，我就干净。}\par

25. 对数字的无知，也阻碍我们理解圣经中以比喻和神秘方式记载的事物。一个坦诚的心灵，如果我可以这样说的话，必然想知道，例如，梅瑟、厄里亚和主自己都禁食四十天这一事实意味着什么。除非通过对数字的认识和思考，否则无法克服解释这一行动所包含的比喻的困难。因为四十这个数字包含十的四倍，表示对万物的知识，以及这种知识与时间的交织。因为昼夜和一年的循环都是由四个周期完成的：昼夜由早晨、中午、傍晚和夜间四个时辰完成；一年由春、夏、秋、冬四个月完成。当我们生活在时间中时，我们必须弃绝并禁戒一切时间中的欢乐，为了我们渴望在其中生活的永恒；尽管通过时间的流逝，我们恰恰学会了轻视时间、寻求永恒这一教训。此外，数字十表示对造物主和受造物的知识，因为造物主内有三位；数字七表示受造物，因为生命和身体。生命由三部分构成，因此天主应当以全心、全灵、全意去爱；身体显然由四种元素构成。因此，当十这个数字与时间联系起来呈现给我们时，即当它被乘以四次时，我们被劝诫要不受玷污地生活，不参与任何时间中的欢乐，即禁食四十天。这由法律所象征的梅瑟、由先知所象征的厄里亚以及主亲自劝诫我们；主仿佛接受了法律和先知的见证，在山上出现在其他两人之间，他的三位门徒惊异地观看。接下来，我们要同样探究，从四十这个数字如何产生五十这个数字，后者在我们的宗教中因五旬节而具有非同寻常的神圣意义；以及这个数字如何因时间的三分（法律之前、法律之下、恩宠之下）或可能因圣父、圣子、圣神之名以及加上三位本身，与至圣教会的奥迹相关联，并达到一百五十三条鱼的数目——这些鱼是在主复活后，当网撒在船的右边时捕获的。\BibleRef{若 21:11}同样，在圣经中，许多其他数字和数字的组合以比喻的方式传达教导，而数字的无知常常使读者无法获得这种教导。\par

26. 对音乐的无知也关闭了许多事物，并使它们变得晦涩。例如，有人并非不熟练地根据瑟和竖琴的区别解释了一些比喻。而且，学者们讨论以下问题并非不合时宜：是否有某种音乐规律迫使十弦瑟恰好有那么多弦？或者，如果没有这样的规律，这个数字本身是否反而更应被视为具有神圣意义——或指法律十诫（若对这数字再有疑问，我们只能将其归因于造物主和受造物），或如上文所解释的十这个数字本身。再者，福音中提到建殿的年数\BibleRef{若 2:20}——即四十六年——具有某种难以言喻的音乐声响，当它联系到主的身体（殿所指的对象）时，迫使许多异端者承认主所取的并非虚假的，而是真实的人性身体。在圣经的几处地方，我们都可以看到数字和音乐受到尊崇。\par

% structure: p0044 source=h2 target=subsection reason=relative-heading-rank; base=h2

\subsection{第十七章——九缪斯传说的起源。}

27. 我们切不可听信异教迷信的虚妄，它们将九缪斯说成是朱庇特与墨丘利的女儿。瓦罗驳斥了这些说法，我怀疑在异教徒中是否有比他对此类事情更为好奇或更有学问的人。他说，某个城邦（我不记得名字了）曾向三位艺术家各订制一组缪斯雕像，作为奉献物置于阿波罗神庙中，意图从他们那里选买雕刻得最美的雕像。结果这三位艺术家创作的作品同样优美，全部九尊雕像都令城邦满意，便全部买下献于阿波罗神庙；他又说，后来诗人赫西俄德给她们都取了名字。因此，并非朱庇特生了九缪斯，而是三位艺术家各创作了三尊。城邦当初订制三尊，并非因它们在异象中被看见，也非因它们以那个数目出现在任何公民眼前，而是因为显而易见，所有声音——即歌曲的质料——按本性有三种。因为声音要么由嗓音产生，如那些不用乐器而用口歌唱的人；要么由吹气产生，如喇叭和笛子；要么由敲击产生，如竖琴和鼓，以及其他所有通过敲击发声的乐器。\par

% structure: p0046 source=h2 target=subsection reason=relative-heading-rank; base=h2

\subsection{第十八章——即使来自世俗来源，也不可轻视任何帮助。}

28. 然而，无论事实是否如瓦罗所述，我们仍不应因异教的迷信而放弃音乐，如果我们能从中获取任何有助于理解圣经的东西；也不应因为我们着手研究竖琴及其他可能帮助我们把握属灵事物的乐器，就必然要忙于他们的戏剧俗物。因为我们不应因他们说墨丘利发明了文字而拒绝学习文字；也不应因他们为正义与美德修建神庙，并宁可以石头形式崇拜那些本应存于心中的事物，而因此离弃正义与美德。相反，让每一个良善而真实的基督徒明白，无论真理在何处被发现，都属于他的主人；当他在他们的宗教文献中认出并承认真理时，让他拒绝迷信的虚构，并为那些人哀伤并避开他们，他们\BibleQuote{虽然认识天主，却没有以祂为天主而光荣祂，也不感谢祂；反而在妄想中变得愚昧，他们自负为智者，反而成了愚妄，将不朽天主的光荣，改换为必朽的人、飞禽、走兽和爬虫形状的偶像。}\BibleRef{罗 1:21}\par

% structure: p0048 source=h2 target=subsection reason=relative-heading-rank; base=h2

\subsection{第十九章——异教知识的两种类型。}

29. 但为更详细地解释这整个主题（因为这是不可忽略的），异教徒中流行两种知识。一种是关于人所设立之物的知识，另一种是关于他们所观察到的——无论是过去发生的事，还是天主所设立之事——的知识。前一种，即处理人类制度的知识，一部分是迷信的，一部分则不是。\par

% structure: p0050 source=h2 target=subsection reason=relative-heading-rank; base=h2

\subsection{第二十章——人类制度的迷信性质。}

30. 所有为制造和崇拜偶像而作的人为安排都是迷信的，它们或涉及崇拜受造物或其一部分为神，或涉及与魔鬼的预兆和盟约的咨询和安排，例如巫术中所使用的，诗人通常不是教导而是歌颂这些。属于此类的，还有肠卜师和占卜官的书，它们以更大胆的欺骗著称。在此类中，我们还必须包括所有医学所谴责的护身符和疗法，无论它们在于咒语，或在于他们称为\Emph{字符}的标记，或在于悬挂或系上甚至以某种方式舞动某些物品，这些做法与身体状况无关，而是与某些隐藏或明显的记号有关；他们用较不冒犯的名称\OriginalTerm{自然疗法}{physica}称呼这些疗法，以便显得不是从事迷信活动，而是利用自然的力量。例如，每只耳朵顶部的耳环，或手指上的鸵鸟骨戒指，或告诉你打嗝时用右手握住左手大拇指。\par

31. 对此，我们还可加上数千种最无聊的习俗，它们需在身体某部分跳动时，或当朋友挽臂同行时，有石头、狗或孩子从中间穿过时遵守。踢一块石头，仿佛它是朋友间的分界者，其害处尚不及打一个无辜的男孩——若他碰巧从并肩行走的人中间跑过。但有趣的是，男孩们有时被狗所报复；因为人们常常迷信到竟敢打从他们中间跑过的狗——然而并非没有代价，因为狗有时会使其攻击者匆忙逃去找真正的医生。同样属于此类的还有以下规则：出门时跨过门槛；穿鞋时若有人打喷嚏，就回到床上；前往某处时若绊倒，就返回家中；衣服被老鼠咬破时，对即将来临的不幸的恐惧甚于对当前损失的悲伤。由此有加图那句妙语：有人来咨询他，说老鼠啃了他的靴子，他回答说：\Quote{那不奇怪，但如果靴子吃了老鼠，那才真是奇怪了。}\par

% structure: p0053 source=h2 target=subsection reason=relative-heading-rank; base=h2

\subsection{第二十一章——占星家的迷信。}

32. 我们也不能从这类迷信中排除那些因关注生日而被称为\Emph{\OriginalTerm{生辰星命家}{genethliaci}}，但现在通常称作\Emph{\OriginalTerm{占星家}{mathematici}}的人。因为这些人，虽或辛苦寻求我们出生时星辰的真位置，甚至有时可能找到，但只要他们由此试图预测我们的行为或行为的后果，就犯下严重错误，并将无经验的人卖入悲惨的束缚。因为任何自由人去到这类占星家那里，他付钱是为了能脱身成为\OriginalTerm{火星}{Mars}或\OriginalTerm{金星}{Venus}的奴隶，或更确切地说，或许是所有星辰的奴隶——那些最初堕入此错误并传给后世的人，给这些星辰起了名字，或因与兽类相似而用兽名，或为赐予那些人荣誉而用人名。这不值得惊奇，当我们想到，甚至在较近我们自己的时代，罗马人曾试图将我们称为\OriginalTerm{晨星}{Lucifer}的星奉献给凯撒的名誉。这或许已经实行，且名字传至久远，只是他的先祖\OriginalTerm{金星}{Venus}已先于他给了这颗星自己的名字，且不能依法将她生前从未拥有、也未曾寻求拥有的事物传给她的继承人。因为，在有一个空位或未因某个先前的亡者而受尊的地方，通常的做法就施行了。例如，我们把月份名\OriginalTerm{Quintilis}{Quintilis}和\OriginalTerm{Sextilis}{Sextilis}改为\OriginalTerm{七月}{July}和\OriginalTerm{八月}{August}，以尤利乌斯·凯撒和奥古斯都·凯撒之名称呼；由此，任何有心者容易看出，上述星辰在天上原本并无它们如今所持的名字。但因人们或被王权强迫、或被人的愚昧驱动去纪念的亡者已死，他们似乎以为，将亡者的名字放在星辰上，就是将他们升到天上。但无论人们如何称呼它们，仍然是天主按照自己的意愿创造并安排的星辰，它们有固定的运行，季节由此区分和变化。当任何人出生时，容易通过那些被圣经以如下语句斥责的人所发现和制定的规则，观察到这运行所到达的节点：\BibleQuote{因为他们既然能深知，能测度世界，为何不更迅速地找主？}\BibleRef{智 13:9}\par

% structure: p0055 source=h2 target=subsection reason=relative-heading-rank; base=h2

\subsection{第二十二章——观察星辰以预测人生事件的愚昧}

33. 但是，想通过这样的观察来预测出生者的品性、行为与命运，是极大的迷妄与疯狂。至少对于但凡对这类事情有所了解的人（事实上这些事情只适合忘掉），这种迷信已经被驳倒，无可置疑。因为观察的是星辰的位置（他们称之为星座），在那人出生之时——这些可怜的人被他们更可怜的被骗者咨询。现在可能发生，在双胞胎的情况下，一个紧随另一个出离母腹，其间几乎没有时间间隔可以被察觉并在星座位置上标记。由此必然得出，双胞胎在许多情况下是在相同星辰下出生，然而他们的行为和遭遇并不等同，反而常如此不同，其中一个最为幸运，另一个最为不幸。例如，我们得知\Person{厄撒乌}和\Person{雅各伯}是孪生兄弟，出生如此接近，以致后生的雅各伯被发现用手抓住在他前面出生的兄弟的脚跟。\BibleRef{创 25:24} 诚然，这两人的出生日与时辰无法以任何方式标记而不给出相同的星座。但这两人在品性、行为、辛劳和命运上有何等不同，圣经如今广传，为万民所诵，业已证实。\par

34. 说双胞胎出生之间分隔的那极小极短暂的时间，在大自然和天体的极速运动中产生重大影响，这是没有道理的。因为，虽然我可以承认它确实产生极大影响，但占星家不能就此在星座中发现它，而他自称是通过观察星座来解读命运的。那么，如果他审视星座时并未发现差异（当然，无论被咨询是\Person{雅各伯}或他的兄弟，星座都是一样的），那么天上存在差异——他鲁莽且轻率地诋毁了它——而他的星盘（他焦虑却徒劳地查看）却没有差异，这于他何益？所以，这些观念，其源头在于人凭妄断随意固定在事物的某些符号上，应归于同一类，正如与魔鬼的盟约与契约。\par

% structure: p0058 source=h2 target=subsection reason=relative-heading-rank; base=h2

\subsection{第二十三章——我们为何拒斥占卜技艺}

35. 因为就这样，贪恋邪恶之事的人，因天主隐秘的判断，被交付去受戏弄和欺骗，作为他们邪恶欲望的公正报应。因为他们被假天使所迷惑和欺骗，这些假天使管辖着世界的最低部分，这是由天主眷顾的律法所安排，并依照祂最可钦佩的万物秩序。这些迷惑和欺骗的结果是，借着这些迷信和有害的占卜方式，许多过去和未来的事得以知晓，并且正如所预言的那样应验；在那些实行迷信仪式的人身上，许多事也符合他们的仪式，他们被这些成功所诱惑，变得更加急切地探究，进一步陷入最有害的错误的迷宫。对我们有益的是，天主圣言并不对这灵魂的淫乱保持沉默；它警告灵魂不要追随这些行为，并非因为那些自称这些行为的人说谎言，而是说：\\Quote{即使他们告诉你们的事应验了，也不要听他们。}\\BibleRef{申 13:1}因为死去的撒慕尔的亡魂预言了真理给撒乌耳王，但这并不意味着那些召唤他亡魂的渎神仪式不那么可憎；虽然\\WorkTitle{宗徒大事录}中那个腹语女子为主的门徒们作了真实的见证，但保禄宗徒并没有因此放过那恶神，而是斥责并将它赶出，从而使那女子洁净了。\\BibleRef{宗 16:16}\par

36. 因此，所有这类技艺，要么是虚妄的，要么是有罪的迷信的一部分，源于人与魔鬼之间有害的交往，基督徒必须完全拒绝和避免它们，如同虚假奸诈的友谊的盟约。宗徒说：\\Quote{并非偶像算得什么，}宗徒说；\\Quote{而是因为他们所祭献的，是祭献给魔鬼，并非祭献给天主；我不愿意你们与魔鬼有交往。}\\BibleRef{格前 10:19}现在，宗徒关于偶像以及为尊敬偶像所献的祭物所说的话，我们也应同样看待所有那些假想的记号，它们或引人崇拜偶像，或引人崇拜受造物或其部分以代替天主，或与注意药符和其他仪式有关——因为这些并非是天主设立作为促进爱天主和爱近人的公共手段，而是浪费可怜之人的心，在私下和自私地追求暂时之物。因此，关于所有这些知识部门，我们必须害怕并回避与魔鬼的交往，魔鬼与他们的首领魔鬼，只致力于关闭并阻挡我们回归的大门。所以，正如从天主创造并安排的星辰中，人们依照自己的幻想画出了虚妄的预兆，同样，从那些在天主眷顾治理下出生或以任何其他方式存在的事物中，如果恰巧发生了一些不寻常的事——例如，一匹骡子生仔，或一个物体被雷电击中——人们常常通过自己的猜测画出预兆，并将它们记录下来，仿佛它们是按规律画出来的。\par

% structure: p0061 source=h2 target=subsection reason=relative-heading-rank; base=h2

\subsection{第24章——迷信仪式所维持的与魔鬼的交往和协议。}

37. 所有这预兆之所以有效，仅仅取决于先前在心灵中与魔鬼达成的默契，这默契可以说是共同的语言，但它们都充满了有害的好奇心、折磨人的焦虑和致命的奴役。因为它们之所以受到注意，并非因为它们本身有意义，而是通过注意和记录它们，它们才获得了意义。因此，对于不同的人，根据他们各自的想法和偏见，这些预兆被弄得不同。因为那些一心欺骗的灵体，会注意为每个人提供与他们自己的猜测和先入之见已经卷入其中的相同种类的预兆。例如，字母X的相同形状，呈十字形，在希腊人中表示一件事，在拉丁人中表示另一件事，这不是由于自然，而是由于关于其含义的约定和预先安排；因此，任何通晓两种语言的人在写给希腊人时使用这个字母的意义，与写给拉丁人时不同。相同的读音\\Emph{beta}，在希腊人中是一个字母的名称，在拉丁人中则是一种蔬菜的名称；当我说\\Emph{lege}，这两个音节对希腊人意味着一件事，对拉丁人意味着另一件事。现在，正如所有这些记号根据每个社群中的安排影响心灵（各人的安排不同，故影响不同人的心灵也不同），而且，人们并不是因为它们已经具有意义才约定它们作为记号，相反，它们之所以现在有意义，是因为人们已经约定了它们；同样，那些维持与魔鬼的毁灭性交往的记号，其意义也完全取决于各人的观察。这在占卜师的仪式中显得非常清楚；因为他们在观察预兆之前和完成观察之后，都努力避免看到鸟的飞行或听到鸟的叫声，因为这些预兆在观察者心中没有预先的安排之外是没有意义的。\par

% structure: p0063 source=h2 target=subsection reason=relative-heading-rank; base=h2

\subsection{第25章——在非迷信的人类制度中，有些是多余的，有些是方便和必要的。}

38. 但当这一切都从基督徒的脑海中切除并根除之后，我们必须审视那些非迷信的人类制度，也就是说，并非与魔鬼结合而设立，而是人与人之间相互结合所设立的制度。因为所有在人类中有效力的安排，由于人们相互同意它们应当有效，都是人类制度；其中有些是多余与奢侈之事，有些是便利与必要之事。因为如果演员在舞蹈中所做的那些记号是出于本性而有效，而非出于人的安排与同意，那么过去传令官就不会在哑剧演员跳舞时向迦太基人民宣布他想要表达什么——这事仍被许多老年人记得，我们常从他们那里听到。我们完全可以相信这一点，因为即使是现在，如果有一个不习惯这种愚蠢行为的人走进剧院，除非有人告诉他这些动作的含义，否则他将白费心思地关注它们。然而，所有人在选择记号时都追求一定程度的相似性，以使记号尽可能像它们所表示的事物。但由于一件事物可能在许多方面与另一件事物相似，除非人们相互约定，这些记号在人类中并不总是具有相同的意义。\par

39. 但关于图画、雕像以及其他这类旨在表现事物的作品，没有人会弄错，尤其是如果它们是由技艺精湛的艺术家所制作；每个人一看到这些相似物，就认出它们所相似的事物。而整个这一类都应算作人类的多余发明，除非在涉及其中任何一件作品时，询问它为何、在何处、何时、由谁授权制作是重要的事。最后，那成千上万以谎言取悦人的寓言和虚构故事，都是人类的发明；任何虚假和撒谎的东西，没有什么比这更被认为是人类所特有的、源于自身的。在人与人之间便利而必要的安排中，应计入那些为了区分性别或等级而在身体服饰和装饰上所做的不同选择；以及无数种记号，没有这些记号，人类的交往要么根本无法进行，要么会带来极大的不便；还有关于度量衡、货币的铸造和称重的安排，这些安排因每个国家和民族而异，以及其他同类事物。这些若非人类的发明，就不会在不同民族中有所不同，也不能在特定民族中由其统治者的意愿而改变。\par

40. 基督徒绝不可忽视这整类为了生活必要交往而便利的人类安排，反而应对它们给予足够的关注，并牢记于心。\par

% structure: p0067 source=h2 target=subsection reason=relative-heading-rank; base=h2

\subsection{二十六、我们当采纳与当避免的人力创造}

因为某些人类的制度在某种程度上是自然事物的表现和相似物。其中，与魔鬼的交往有关的，如前所述，必须彻底弃绝并厌恶；另一方面，与人类相互交往有关的，只要不是奢侈和多余之事，就应当采纳，尤其是阅读所必需的字母形式，以及所需的各种语言——我在上文中已谈到此事。速记字符也属于此类，熟悉它们的人被称为速记员。所有这些都有用，学习它们并无不合法的，也不使我们陷入迷信，或因奢侈而削弱我们，只要它们占据我们的心不致妨碍更重要的目标——它们应当服务于这些目标。\par

% structure: p0069 source=h2 target=subsection reason=relative-heading-rank; base=h2

\subsection{二十七、某些知识领域（并非纯粹人力创造）有助于我们诠释圣经}

41. 但是，谈到下一点，我们不应把那些人所传给我们、不是作为他们自己的安排，而是作为对过去事件和天主上智安排的研究结果，算作人类制度。这些事物中，有些涉及身体感官，有些涉及理智。通过身体感官达到的，我们要么相信见证，要么在它们被指出时察觉，要么从经验中推论。\par

% structure: p0071 source=h2 target=subsection reason=relative-heading-rank; base=h2

\subsection{二十八、历史在多大程度上有助益}

42. 因此，我们从历史中学到的任何关于过去年代年表的事物，即使是在教会之外作为儿童教育的内容而学的，也极大地帮助我们理解圣经。因为我们经常借助奥林匹亚纪和执政官的名字来寻求各种信息；而对我们的主出生和受难时执政官的无知，使一些人误以为祂受难时是四十六岁，因为犹太人告诉祂，圣殿（祂以此象征祂的身体）建造了四十六年。\BibleRef{若 2:19}现在我们凭福音作者的权威知道，祂受洗时约三十岁。\BibleRef{路 3:23}祂此后生活的年数，虽然我们可以通过将祂的行事结合起来推算出来，但为了不产生任何疑点，通过将世俗历史与福音进行比较，可以更清楚、更确定地查明。然而，仍显而易见的是，圣殿建造四十六年的说法并非毫无目的，由此可以理解那更为隐秘的身体的塑造，为了我们，天主独生子——万物由祂造成——屈尊披上这身体。\par

43. 至于历史的用处，且不提希腊人，我们自己的\Person{盎博罗削}解决了多大的一个问题！因为，当\Person{柏拉图}的读者和赞美者们敢于诽谤性地断言我们的主耶稣基督所有那些他们不得不钦佩和赞扬的言论，都是从\Person{柏拉图}的著作中学来的——因为他们坚称，无可否认\Person{柏拉图}生活在我们主来临之前很久！——这位杰出的主教，通过他对世俗历史的研究，发现\Person{柏拉图}在\Person{耶肋米亚}先知在埃及的时候曾旅行到那里，难道没有表明，\Person{柏拉图}更可能是通过\Person{耶肋米亚}被引入我们的文献，从而能够教导和写出那些如此公正地受到赞扬的观点吗？因为就连\Person{毕达哥拉斯}本人（这些人声称\Person{柏拉图}从他的后继者那里学到了神学）的生活年代也不早于那个希伯来民族的书籍，在一神崇拜兴起于那个民族之中，我们的主按肉身也是出自他们。因此，当我们思考这些年代时，变得更为可能的是，那些哲学家们从我们的文献中学到了他们所说的一切美好而真实的东西，而不是我们的主耶稣基督从\Person{柏拉图}的著作中学到了什么——相信这一点是极其愚蠢的。\par

44. 而且，即使在历史叙述中描述了前人的人类制度，历史本身也不应被算作人类制度；因为那些已经过去、无法挽回的事情应被算作属于时间的过程，而天主是时间的创造者和治理者。因为讲述已经发生的事是一回事，指示应当做什么是另一回事。历史忠实地、有益地叙述了已经发生的事；而占卜者的书籍以及所有同类著作，则旨在教导应当做什么或遵守什么，使用的是建议者的胆大妄为，而非叙述者的忠实。\par

% structure: p0075 source=h2 target=subsection reason=relative-heading-rank; base=h2

\subsection{第二十九章——自然科学在何种程度上是释经的辅助。}

45. 还有一种类似于描述的叙述，其中不是过去的状态，而是现存的状态被告知给那些不知道它的人。属于这一类的是所有关于地方位置、动物、树木、草药、石头及其他物体的性质的作品。关于这一类我在上面已经讨论过，并且表明这类知识有助于解决圣经中的难题，这不是说这些物体要按照某些标记用作药物或迷信的工具；因为那种知识我已经将其与现在谈论的合法且自由的知识区分开来。因为说\Quote{如果你碾碎这种草药并喝下它，它会消除你胃里的疼痛}是一回事；说\Quote{如果你把这草药挂在脖子上，它会消除你胃里的疼痛}是另一回事。在前一种情况下，有益健康的混合物得到认可，在后一种情况下，迷信的咒语受到谴责；尽管确实，在不使用咒语、呼求和标记的情况下，常常不确定是否系在身上或以任何方式固定以治愈身体的东西，是通过自然效力起作用（在这种情况下可以自由使用），还是通过某种咒语起作用（在这种情况下，基督徒应当更小心地避免它，因为它似乎越有效）。但是，当一件东西为什么有效的原因不明显时，使用它的意图就非常重要，至少在治疗或调理身体方面，无论是在医学还是农业中。\par

46. 再者，关于星辰的知识，不是叙述的事情，而是描述的事情。然而，圣经中提到的星辰很少。月亮的运行，通常用于庆祝我们主受难的周年纪念，大多数人都知道；至于其余天体的升起、落下及其他运动，则很少有人完全了解。这种知识本身虽然不涉及迷信，但对解释圣经的帮助甚少，几乎没有任何帮助，而且由于无益地吸引注意力，反而是一种妨碍；而且它与占卜命运者非常有害的错误密切相关，因此忽略它更为方便和得体。此外，除了对现存状态的描述，它还包含某种类似过去叙述的东西；因为人们可以从星辰的当前位置和运动出发，依据规则追溯它们过去的运动。它还涉及对未来的规则性预测，不是以预感和预兆的方式，而是通过确定的计算；其目的不是像\OriginalTerm{占星家}{genethliaci}那样荒谬地从它们中获取关于我们自身行为和命运的信息，而只是关于天体本身的运动。因为，就像计算月龄的人，当他发现今天的月龄后，可以告诉任何年前月龄是多少，或任何年后月龄是多少一样，擅长这种计算的人通常也能对每一个天体回答类似的问题。关于所有这些知识，就其用处而言，我已经陈述了我的观点。\par

% structure: p0078 source=h2 target=subsection reason=relative-heading-rank; base=h2

\subsection{第三十章——机械工艺对释经学有何贡献。}

47. 此外，至于其余的艺术，无论是那些通过它们制造某种东西，当工匠的努力结束后，其作品仍然存在，如房屋、长凳、盘子及其他类似之物；还是那些可以说是辅助天主运作的，如医学、农业和航海；或是那些唯一的结果是动作的，如舞蹈、赛跑和摔跤——在所有这些艺术中，经验教导我们从过去推断未来。因为任何精通这些艺术的人，在操作中移动肢体时，都不会不将过去的记忆与未来的期望联系起来。对于这些艺术，应获得非常肤浅和粗略的知识，不是为了实践它们（除非有义务迫使我们，这个问题我目前不触及），而是为了对它们形成判断，这样当圣经使用源自这些艺术的修辞时，我们就不至于完全不知道圣经想要传达什么。\par

% structure: p0080 source=h2 target=subsection reason=relative-heading-rank; base=h2

\subsection{三十一、论辩术的使用。论谬误。}

48. 剩余的是那些不属于身体感官，而属于理智的知识分支，其中推理科学和数字科学是主要的。推理科学对于探究和阐明圣经中出现的各种问题大有裨益，只是在使用它时，我们必须警惕好争论的倾向，以及以陷阱诱捕对手的幼稚虚荣。因为有许多所谓的\Emph{诡辩}，它们是虚假的推理结论，却如此逼真地模仿真实结论，以至于不仅欺骗迟钝之人，就连聪慧之人，若不加提防，也会受骗。例如，一人对其谈话对象提出命题：\Quote{凡我所是，你则不是。}对方表示赞同，因为该命题部分为真，一人狡猾，另一人单纯。随后第一人加上：\Quote{我是一人；}对方也赞同后，第一人得出结论：\Quote{那么你便不是人。}对于这类诱骗性论证，依我判断，圣经在那处表达了憎恶，其中记载：\Quote{有一人在言语中显示智慧，却仍被憎恨；}不过，那种并非旨在诱骗，而仅追求言语修饰超过庄重目的的风格，确实也被称为诡辩的。\par

49. 也存在有效的推理过程，它们通过将与之辩论之人的错误推向其逻辑结论，而引向虚假结论；有时，一位良善博学之人会得出这些结论，目的是使那错误源之人因这些结果而感到羞愧，并因此引导他放弃自己的错误，因为他发现如果他想保留原有意见，他就必须也持有其他他所谴责的意见。例如，宗徒在说\BibleQuote{那么基督就没有复活}，以及\BibleQuote{那么我们的宣讲就是徒劳的，你们的信德也是徒劳的}；\BibleRef{格前 15:13}时，并没有得出真实的结论；随后他还得出其他全然虚假的推论；因为基督已经复活，宣扬这一事实之人的宣讲并非徒劳，那些相信之人的信德也并非徒劳。但是所有这些虚假推论都是从那说没有死人复活之人的意见中有效推论出来的。既然这些推论被斥为虚假，那么，既然如果死人没有复活它们就会为真，就必有死人复活。因此，有效的结论不仅可以从真实的命题中得出，也可以从虚假的命题中得出；推理的有效法则很容易在教会之外的学校中学习到。但命题的真实性则必须在教会的圣经中探究。\par

% structure: p0083 source=h2 target=subsection reason=relative-heading-rank; base=h2

\subsection{三十二、有效的逻辑序列并非人所发明，而仅是人所观察。}

50. 然而，逻辑序列的有效性并不是人所发明的，而是人所观察并记录的，以便他们能够学习和教授它；因为它永恒地存在于事物的理性中，且源于天主。因为，正如叙述事件次序的人并非自己创造该次序；描述地域位置、或动物、植物或矿物的本性的人，并非描述人的安排；指出星辰及其运动的人，并非指出他自己或任何他人所命定的东西——同样，说\Quote{当后件为假时，前件也必为假}的人，说出了最真实的话；但他并非自己使之如此，而只是指出它是如此。我刚才引用的宗徒保禄的推理正是基于这一规则。因为前件是\Quote{没有死人复活}——这是宗徒所要推翻的犯错之人的立场。接着，从这个前件（即没有死人复活的断言）必然推出的后果是\Quote{那么基督就没有复活}。但这一后果是假的，因为基督已经复活；因此前件也是假的。而前件就是没有死人复活。因此我们得出结论：必有死人复活。现在这一切可以简洁地表达如下：如果没有死人复活，那么基督就没有复活；但基督已经复活，因此必有死人复活。因此，这一规则——即当后件被否定时，前件也必须被否定——不是人所创造的，而只是人所指出的。这一规则涉及推理的有效性，而非陈述的真实性。\par

% structure: p0085 source=h2 target=subsection reason=relative-heading-rank; base=h2

\subsection{三十三、虚假推论可以从有效推理中得出，反之亦然。}

51. 然而，在这段论述复活之处，推论的原则是有效的，而得到的结论也是真实的。但在假结论的情况下，也有一种推论的有效性，如下所示。假设某人承认：如果蜗牛是动物，它就有声音。一旦承认了这一点，那么，当证明了蜗牛没有声音时，就得出结论（因为当后件被证明为假时，前件也为假），蜗牛不是动物。现在，这个结论是假的，但它却是从虚假的承认中得出的真实而有效的推论。因此，陈述的真实性有其自身的价值；推论的有效性取决于与己争论之人的陈述或承认。因此，正如我上面所说，可以由一个有效的推理过程得出一个错误的推论，以便我们想要纠正其错误的人，在看到其逻辑结论完全站不住脚时，为他承认了前件而感到遗憾。因此，很容易理解，正如在见解错误时推论可以是有效的一样，在见解正确时推论也可能是不健全的。例如，假设某人提出一个陈述：\Quote{如果这人是义人，他就是善的}，我们承认其真实性。然后他补充：\Quote{但他不是义人}；当我们也承认这一点时，他得出结论：\Quote{因此他不是善的}。现在，虽然这些陈述中的每一个可能都是真实的，但推论的原则却是不健全的。因为，正如当后件被证明为假时前件也为假，但当前件被证明为假时后件却不一定为假，这并不成立。因为这一陈述是真的：\Quote{如果他是演说家，他就是人}。但如果我们在其后加上：\Quote{他不是演说家}，则不能推出\Quote{他不是人}。\par

% structure: p0087 source=h2 target=subsection reason=relative-heading-rank; base=h2

\subsection{Chapter 34.— 知道推论的原则与知道见解的真实是两回事。}

52. 因此，知道推论的原则是一回事，知道见解的真实是另一回事。在前者中，我们学习什么是后承，什么是非后承，以及什么是不兼容。后承的例子是：\Quote{如果他是演说家，他就是人}；非后承的例子是：\Quote{如果他是人，他就是演说家}；不兼容的例子是：\Quote{如果他是人，他就是四足兽}。在这些例子中，我们判断联系。然而，关于见解的真实，我们必须考虑命题本身，而不是它们之间的相互联系；但是，当我们不确定的命题通过有效的推论与真实而确定的命题联系起来时，它们本身也必然成为确定的。现在，有些人一旦确定了推论的有效性，就自夸好像这也包含了命题的真实性。另外，许多持有真实见解的人，因为不懂推论的原则而无端地自卑；然而，知道有死人复活的人，肯定比只知道\Quote{如果没有死人复活，那么基督就没有复活}这一推论的人更好。\par

% structure: p0089 source=h2 target=subsection reason=relative-heading-rank; base=h2

\subsection{Chapter 35.— 定义的科学不是假的，虽然它可以应用于虚假事物。}

53. 再者，定义、划分和分立的科学，虽然经常应用于虚假事物，但它本身并非虚假，也不是由人的谋划所构成，而是从事物的理性中演化出来的。因为虽然诗人们已将其应用于他们的虚构，而假哲学家，甚至异端者——即假基督徒——也将其应用于他们的错误教义，但这并不成为它虚假的理由；例如，在定义、划分或分立中，不得包含任何与当前事务无关的东西，也不得遗漏任何与之有关的东西。这是真的，即使所要定义或划分的事物并非真实。因为连虚假本身也被定义，当我们说虚假是对事物状态的陈述，而该状态并不如我们所陈述的那样；这个定义是真的，尽管虚假本身不可能是真的。我们也可以对它进行划分，说有两种虚假：一种涉及根本不可能真实的事物，另一种涉及虽非真实但可能真实的事物。例如，说七加三是十一的人，说的任何情况下都不可能是真的；但说在一月朔日下了雨的人，尽管事实可能并非如此，但说的却是可能发生过的事。因此，对虚假事物的定义和划分可能完全真实，尽管虚假本身当然不可能是真的。\par

% structure: p0091 source=h2 target=subsection reason=relative-heading-rank; base=h2

\subsection{Chapter 36.— 雄辩术的规则是真的，虽然有时用于说服人相信虚假的事。}

54. 还有一种更加丰富的论证规则，称为雄辩术，这些规则并不因其可用于说服人相信虚假之事而减少其真理性；但由于它们也可用于支持真理，所以应当受到指责的不是这种能力本身，而是那些将其用于不良目的的人的乖僻。并非由于人类的一种安排，使得情感的表达获得听者的好感，或者简短清晰的叙述有效果，或者变化吸引人的注意力而不使其厌倦。其他同类规则也是如此：无论它们所应用的事业是真是假，只要它们能有效地产生知识或信念，或者感动人心使其欲求或厌恶，它们本身便是真实的。人们是发现了这些事实，而不是制定了它们应如此。\par

% structure: p0093 source=h2 target=subsection reason=relative-heading-rank; base=h2

\subsection{Chapter 37.— 修辞学与辩证法的运用。}

55. 然而，这门技艺一旦学得，与其说用于确定含义，不如说在含义确定后用于阐述它。但先前所论及的那门涉及推理、定义和分类的技艺，在发现含义方面有极大助益，前提是人们不陷入这样的错误：以为学得这些便掌握了幸福生活的真正秘诀。然而，有时人们发现，获得学习这些科学所要达到的目标，比经历这些规则极为复杂棘手的训练反而更容易。这好比一个人想制定行走的规则，他提醒你不可在放下前脚之前抬起后脚，然后详细描述你应该如何活动关节和膝盖的枢纽。因为他说的是实情，行走别无他法；但人们发现，实际行走时执行这些动作比在行走过程中注意它们，或听人讲述时理解它们更容易。反之，那些不能行走的人，对这些指示更不关心，因为他们无法通过尝试来验证它们。同样，一个聪明人常常在理解推理规则之前，更快地看出推理不成立；而一个愚钝的人则看不出不成立，更谈不上掌握规则。至于所有这些规律，我们从中获得的快乐更多来自它们是真理的展示，而非对辩论或形成意见的帮助，或许除了它们使理智得到更好的训练之外。然而，我们必须小心，不能让它们同时使理智更倾向于作恶或虚荣——也就是说，不能让学得它们的人有倾向用似是而非的言论和引人上钩的提问来误导他人，或使他们以为自己获得了某种了不起的东西，足以胜过善人和无辜者。\par

% structure: p0095 source=h2 target=subsection reason=relative-heading-rank; base=h2

\subsection{第三十八章——数的科学并非人所创造，而仅为所发现。}

56. 现在论及数的科学，即使最迟钝的理解力也能明白，它并非人所创造，而是藉研究发现的。因为，尽管维吉尔可以凭自己的喜好把\Emph{Italia}的第一个音节读长，而古人则读短，但没有任何人能凭自己的喜好规定三乘三不是九，或不构成一个平方，或不是三的三倍，也不是六的一倍半，或断定由于奇数没有一半，因而它们不是任何数的两倍这一说法不成立。那么，无论数字就其本身来考虑，还是应用于图形、声音或其他运动的规律，它们都有固定的法则，这些法则并非人所制定，而是由机智者的敏锐所揭示的。\par

57. 然而，那将这些东西看得如此之高，以致倾向于自夸为有学问之人，而不是去探求他所知觉为真的事物从何获得其真理性，以及他所知觉为不变的事物又从何获得其真理性与不变性；并且，从形体现象上升到人的心智，发现心智也是可变的（因为它有时受教，有时未受教），尽管它介于其上不变真理与其下可变事物之间，却不努力使一切事物都归向对唯一天主的赞美和爱慕——他知道万物皆由他而存在——，我再说，如此行事的人，或可显得有学问，但绝不能被视为有智慧。\par

% structure: p0098 source=h2 target=subsection reason=relative-heading-rank; base=h2

\subsection{第三十九章——应对上述哪些学科予以关注，并以何种精神}

58. 因此，我认为宜于警告那些敬畏天主、寻求人生幸福的勤勉有才的年轻人，不可轻率地投身于在基督的教会之外流行的那些学问分支，仿佛它们能确保他们获得所寻求的幸福；而要清醒而谨慎地辨别它们。如果他们发现其中某些学问是因创立者喜好的不同而变异、因错误的臆测而晦涩难解，尤其是涉及藉着关于记号的盟约与契约而与魔鬼共融，就应完全弃绝并深恶痛绝。也要让年轻人从那些不必要且奢侈的人为制度中收回注意力。但为了今生的必需，我们不可忽视那些使我们能与旁人交往的人为安排。然而，我认为在外教人那里找到的其他学问分支中，除了关于物体（无论过去或现在）与身体感官相关的知识（其中也包括实用机械技艺的实验和结论），以及推理科学和数的科学之外，没有什么是可用的。关于这一切，我们必须遵守这一格言：\Quote{凡事不可过分}，尤其是在那些属于感官、受制于空间和时间关系的事物上。\par

59. 那么，有些人在圣经中出现的所有词语和名称上——不论是希伯来语、叙利亚语、埃及语还是其他语言——所做的，是如何将他们留在圣经中未经解释的词语逐个取出并加以解释；以及\Person{欧瑟伯}在过去的史上所做的，为的是解决圣经中那些需要历史知识才能解答的问题——我说，这些人在这类事情上所完成的，使基督徒不必为了少许知识而在许多课题上耗尽精力，同样的事，我想，也可以在其他事上做到，若是有胜任者愿意本着慈爱之心为弟兄们的益处承担这项工作。这样，他就可以把圣经中提到的未知地点、动物、植物、树木、石头、金属以及其他种类的事物分类排列，只处理这些，并将他的记述写下来。对于数字也可以这样做，解释那些只出现在圣经中的数字的理论并记录下来。也许这些事已经部分或全部被人完成（正如我发现许多我从未想过的事情已被善良博学的基督徒完成并写了下来），但它们或是被粗心的人群所遗失，或是被嫉妒者所隐藏。我不确定同样的事是否能用于推理的理论；但在我看来不能，因为推理像神经一样贯穿圣经的全部结构，因此在解释歧义段落（我将在以后讨论）方面，比确定未知符号的意义（我现在讨论的主题）更有用。\par

% structure: p0101 source=h2 target=subsection reason=relative-heading-rank; base=h2

\subsection{第四十章——外教人所说的一切正确之事，我们都应取为己用。}

60. 还有，倘若那些被称为哲学家、尤其是柏拉图学派的人，说过任何真实且与我们信仰相符的话，我们不仅不应回避，反而应从那些非法占有它们的人那里取来为我们所用。因为，正如埃及人不仅拥有以色列子民所憎恶并逃避的偶像和重担，也拥有金银器皿和衣物，这同一子民在离开埃及时据为己有，预备作更好的用途，并非凭自己的权柄，而是出于天主的命令；埃及人自己则在无知中提供了他们自己未曾善用的物品；同样，外教学问的各分支不仅包含虚妄而迷信的幻想和不必要的辛劳重担——我们每个人在基督的领导下从外教人的团契中出来时，都应当憎恶并躲避这些——也包含更适宜于真理使用的自由教育，以及一些极其卓越的道德规范；甚至关于唯一真神敬拜的某些真理也存在于它们中间。这些，可以说，是他们的金银，他们自己并未创造，而是从处处散布的天主眷顾的矿藏中挖掘出来，却邪僻而非法地用于魔鬼的崇拜。因此，基督徒在精神上从这些人悲惨的团契中分离出来时，应当从他们那里取走这些，并用于宣讲福音的正当用途。他们的衣物——即那些适应于此生不可或缺的人际交往的人类制度——我们也必须取来转为基督化的用途。\par

61. 我们弟兄中有多么好的忠信之人还做了什么呢？我们难道没有看到，\Person{西彼廉}，那位最具说服力的教师和最蒙福的殉道者，当他从埃及出来时，带着多少金银和衣物吗？\Person{拉克坦提乌斯}带来了多少？还有\Person{维克托里努斯}、\Person{奥普塔图斯}、\Person{希拉里}，更不必说还活着的人！希腊人中有数不清的人借用了多少？而在所有这些之前，天主最忠信的仆人\Person{梅瑟}已经做了同样的事；因为经上记载他学会了埃及人一切的智慧。\BibleRef{宗 7:22} 外教人的迷信（尤其是在那些年代，它反抗基督的轭而迫害基督徒）绝不会向这些人提供它认为有用的知识分支，若是它怀疑他们即将把它们用于敬拜唯一真神，从而推翻对偶像的虚妄崇拜。但他们把自己金银和衣物给了正要离开埃及的天主子民，并不知道他们所给的东西将如何转为基督的侍奉。因为出埃及时所发生的事，无疑是一个预表，预示了现在所发生的事。我说这话，并不贬低任何其他同样好或更好的解释。\par

% structure: p0104 source=h2 target=subsection reason=relative-heading-rank; base=h2

\subsection{第四十一章——研读圣经需要何种精神。}

62. 但是，当研读圣经的学生，按我所指示的方式准备好，开始他的研究时，让他时常默想宗徒的那句话：\BibleQuote{知识使人自高自大，爱德却立人。}\BibleRef{格前 8:1} 因为他会感到，无论他从埃及带出的财富有多少，除非他过了逾越节，否则他不能安全。现在，基督我们的逾越节，为我们被祭献，\BibleRef{格前 5:7} 而基督的祭献没有比祂亲自向那些祂看见在埃及法郎手下劳苦重担的人所发出的召唤更清楚地教导我们：\BibleQuote{凡劳苦和负重担的，你们都到我这里来，我要使你们安息。你们背起我的轭，跟我学习，因为我是良善心谦的：这样你们的灵魂必要找到安息。因为我的轭是柔和的，我的担子是轻省的。}\BibleRef{玛 11:28} 对这轭，除了那些良善心谦的人，谁会觉得轻省呢？知识不会使他们自高自大，而爱德却立人。那么，让他们记得，那时那些在预像和阴影中庆祝逾越节的人，当被命令用羔羊的血涂在门框上时，他们便用\OriginalTerm{牛膝草}{hyssop}来涂。\BibleRef{出 12:22} 这是一种柔和谦卑的草，然而没有什么比它的根更强壮、更深入；为使我们在爱德中生根建基，才能和众圣徒一同领悟那宽、长、高、深，\BibleRef{弗 3:17}——即领悟我们主的十字架，其宽度由横木（双手伸展其上）表示，其长度由从地面到横木的部分（整个身体从头以下被钉在其上）表示，其高度由从横木到顶端（头所在处）表示，其深度由埋入地下的隐藏部分表示。而藉着这十字架的标记，一切基督徒的行为都被象征化了，即：在基督内行善，恒心依持祂，盼望天堂，且不亵渎圣事。被这基督徒的行为净化后，我们甚至能认识\BibleQuote{基督那超越知识的爱}，祂与圣父同等，万有是藉祂而造成的，\BibleQuote{使我们充满天主的一切富裕}\BibleRef{弗 3:19}。此外，牛膝草还有一种洁净的功效，使人的心胸不致因那自高自大的知识而膨胀，也不致虚夸从埃及带出的财富。圣咏作者说：\BibleQuote{求你用牛膝草洗涤我，我就洁净；洗涤我，我就比雪更白。求你使我听到喜乐和欢欣的声音。}然后他立刻加上，以表明牛膝草所指的是从骄傲中洁净：\BibleQuote{你所折断的骨头，必要欢欣。}\par

% structure: p0106 source=h2 target=subsection reason=relative-heading-rank; base=h2

\subsection{第四十二章 圣经与世俗作者相比}

63. 然而，以色列子民从埃及带出的金银衣服的财富，与后来在耶路撒冷所获得、并在所罗门王时代达到顶峰的财富相比，是多么贫乏；同样，所有从外邦书籍中收集的有用知识，与圣经的知识相比，也是多么贫乏。因为无论一个人从其他来源学到什么，若是有害的，在那里必受谴责；若是有用的，在其中已经包含。而且，虽然每个人都可以在那里找到他随处学到的有用之物，但他在那里会找到更丰富的、其他地方找不到的东西，这些只能在圣经那奇妙崇高又奇妙简朴中学习。\par

当读者掌握了上述指示，使未知的记号不再成为他的障碍；当他良善心谦，服从基督轻松的轭，背负祂轻省的担子，在信德中生根、奠基、建造，以致知识不能使他自高自大时，让他进而思考并讨论圣经中那些模棱两可的记号。关于这些，我将在第三卷中，按主所愿赐予的，尽力加以说明。\par

\SourceRecord{Translated by James Shaw.；Nicene and Post-Nicene Fathers, First Series；Vol. 2.；Edited by Philip Schaff.；Buffalo, NY: Christian Literature Publishing Co.,；1887.；<http://www.newadvent.org/fathers/12022.htm>.}

% structure: p0001 source=h1 target=section reason=page-title

\section{\WorkTitle{论基督教教义（第三卷）}}

% structure: p0002 source=h2 target=subsection reason=relative-heading-rank; base=h2

\subsection{第一章  前书摘要及本书范围}

一、敬畏天主的人勤勉在圣经中寻求祂旨意的知识。当他藉着虔诚而变得温和，以致不爱争辩；当他具备了语言知识，不被生僻的词语和表达方式所阻碍，并具备某些必要事物的知识，不致对借喻所用之物的力量和性质无知；此外，还得到经文准确性的帮助——这准确性是通过校勘上的技巧和细心获得的——如此预备之后，让他着手探究和解决圣经中的歧义。为了不使他被暧昧的表记引入歧途，就我所能给予的指示而言（然而可能因他智力高超或所享之光更明亮，他会嘲笑我将指出的方法为幼稚），但正如我所说，就我所能给予的指示而言，让那处于可受我指教心态的人知道，圣经的歧义要么在于本义词，要么在于比喻义，这些类别我在第二卷中已经描述过。\par

% structure: p0004 source=h2 target=subsection reason=relative-heading-rank; base=h2

\subsection{第二章  通过标点消除歧义的规则}

二、但当本义词造成圣经歧义时，我们首先要看我们在标点或发音上是否无误。因此，如果注意经文后，仍不确定应如何标点或发音，让读者咨询他从圣经较明显段落及教会权威所收集的信德规则——此事我在第一卷谈论事物时已充分处理。但如果两种读法（或所有读法，若不止两种）都给出符合信德的意义，则需查看上下文，包括前后文，看众多解释中哪一种与上下文契合并被其允许。\par

三、现在看一些例子。异端对标点如下：\BibleRef{若 1:1} \Quote{\Emph{在起初已有圣言，圣言与天主同在，圣言是}}，以便下一句变成 \Quote{\Emph{这圣言在起初就与天主同在}}，这源于不愿承认圣言是天主。但必须根据信德规则拒绝，该规则关于天主圣三的平等，指示我们说：\Quote{\Emph{圣言是天主}}；然后加上：\Quote{\Emph{这圣言在起初就与天主同在}}。\par

四、但以下标点歧义无论你如何理解都不违背信仰，因此必须由上下文决定。这是宗徒所说的：\Quote{我不知该选择什么：我夹在两者之间，有离世的渴望，与基督同在，这实在好得多；但留在肉身中，对你们更为必要。} \BibleRef{斐 1:22} 现在不确定我们是否应读作 \Quote{\Emph{ex duobus concupiscentiam habens}}（渴望两件事），或 \Quote{\Emph{compellor autem ex duobus}}（我夹在两者之间）；然后加上 \Quote{\Emph{concupiscentiam habens dissolvi, et esse cum Christo}}（有离世与基督同在的渴望）。但由于接着有 \Quote{\Emph{multo enim magis optimum}}（因为这实在好得多），显然他说他渴望那更好的；因此，虽然他夹在两者之间，但他渴望其一，并看到另一的必要性；即渴望与基督同在，而必须留在肉身中。现在这个歧义由随后一个词解决，这个词译作 \Emph{enim}（因为）；省略了这个虚词的译者宁愿采用使宗徒似乎不仅夹在两者之间，而且渴望两事的解释。因此我们必须这样标点句子：\Quote{\Emph{et quid eligam ignoro: compellor autem ex duobus}}（我不知该选择什么：我夹在两者之间）；然后在此句号后接着：\Quote{\Emph{concupiscentiam habens dissolvi, et esse cum Christo}}（有离世与基督同在的渴望）。并且，似乎有人问他为什么更渴望这个，他补充说：\Quote{\Emph{multo enim magis optimum}}（因为这实在好得多）。那么，他为什么夹在两者之间？因为他有留下的需要，他用这些话补充：\Quote{\Emph{manere in carne necessarium propter vos}}（但留在肉身中，对你们更为必要）。\par

五、然而，当歧义既不能由信德规则又不能由上下文澄清时，我们不妨按所建议的任何方法标点句子。如致格林多人书那段：\Quote{所以，亲爱的，我们既有这些恩许，就当洁净自己，除去肉体和心灵的一切污秽，在敬畏天主中成圣。容纳我们；我们没有亏负过谁。} \BibleRef{格后 7:1} 我们不确定应读作 \Quote{\Emph{mundemus nos ab omni coinquinatione carnis et spiritus}}（洁净自己，除去肉体和心灵的一切污秽），与经文 \Quote{好使她身心圣洁} \BibleRef{格前 7:34} 一致；或读作 \Quote{\Emph{mundemus nos ab omni coinquinatione carnis}}（洁净自己，除去肉体的一切污秽），使下一句为 \Quote{\Emph{et spiritus perficientes sanctificationem in timore Dei capite nos}}（并在敬畏天主中成全心灵的圣洁，容纳我们）。因此，这类标点歧义留给读者自行判断。\par

% structure: p0009 source=h2 target=subsection reason=relative-heading-rank; base=h2

\subsection{第三章  发音如何帮助消除歧义。不同的询问方式}

6. 关于歧义标点的一切指示，对于不确定的发音也应同样遵守。这些发音问题，除非归咎于读者的粗心，否则都要凭信仰的准则，或参照前后文来纠正；如果这两种方法都不成功，它们就会保持不确定，但无论读者如何发音，都不致犯错误。例如，倘若我们的信德（即天主必不指控祂的选民，基督必不定祂选民的罪）没有拦阻，这段经文：\Quote{谁能控告天主所拣选的人呢？}就可以这样发音，使接下来的一句成为对这问题的回答：\Quote{天主使人成义}，并构成第二个问题：\Quote{谁能定罪呢？}回答是：\Quote{基督耶稣死了。}\BibleRef{罗 8:33}但若相信这一点，那真是疯狂至极；因此，该段经文将会这样发音：使第一部分成为探询性问题，第二部分成为修辞性反问。古人说，探询与反问的区别在于：探询可以有许多答案，而反问的答案只能是\Quote{不}或\Quote{是}。那么，该段应如此发音：在探询\Quote{谁能控告天主所拣选的人呢？}之后，接下来以反问形式提出：\Quote{使人成义的天主吗？}——隐含的回答是\Quote{不}。同样，我们将有探询\Quote{谁能定罪呢？}而答案又以反问形式出现：\Quote{是那死了的基督吗？不，祂已复活了，现今在天主右边，也为我们转求。}——对每一问题隐含的回答都是\Quote{不}。另一方面，在宗徒所说的那段话：\Quote{那么，我们说什么呢？那未追求义的外邦人却得到了义。}\BibleRef{罗 9:30}除非在探询\Quote{我们说什么呢？}之后，接下来的话作为对这问题的回答：\Quote{那未追求义的外邦人却得到了义}，否则就与下文不一致了。至于纳塔乃耳的那句话：\Quote{从纳匝肋能出什么好事吗？}\BibleRef{若 1:47}无论一个人选择用肯定的语气来读（这样只有\Quote{从纳匝肋}属于问句），还是用怀疑和犹豫的语气来问整句话，我看不出有什么差别。但两种意思都不违背信仰。\par

7. 还有一种歧义源于音节读音的不确定，这当然与发音有关。例如，在经文中：\Quote{\OriginalTerm{我的骨头}{os meum}没有向祢隐藏，祢在暗中制造了它}，读者不清楚该将\OriginalTerm{拉丁词}{os}一词读作短音还是长音。若读作短音，它是\OriginalTerm{骨头}{ossa}的单数；若读作长音，它是\OriginalTerm{口}{ora}的单数。这类困难可通过查看原文来解决，因为在希腊文中我们找到的不是\OriginalTerm{口}{\GreekText{στόμα}}，而是\OriginalTerm{骨}{\GreekText{ὀστέον}}。因此，俗语在传达意思上往往比知识分子的纯净语言更有用。因为我宁愿有这句粗鄙的话：\Quote{\OriginalTerm{你的骨头没有向你隐藏}{non est absconditum a te ossum meum}}，也不愿经文用更优美的拉丁语却意思不够清楚。但有时当音节发音不确定时，可由同一句中邻近的词来确定。例如，宗徒的那句话：\Quote{关于这些事，\OriginalTerm{我预先告诉你们}{prædico}，正如\OriginalTerm{我以前说过}{prœdixi}的，行这样事的人不能继承天主的国。}\BibleRef{迦 5:21}如果他只说：\Quote{关于这些事，\OriginalTerm{我预先告诉你们}{quæ prædico vobis}}，而没有加上\Quote{\OriginalTerm{正如我以前说过}{sicut prœdixi}}，那么除非回到原文，我们无法知道\OriginalTerm{拉丁词}{prædico}一词的中间音节应读作长音还是短音。但既然这样，显然应读作长音；因为他不是说\OriginalTerm{正如我预告过}{sicut prœdicavi}，而是说\OriginalTerm{正如我以前说过}{sicut prædixi}。\par

% structure: p0012 source=h2 target=subsection reason=relative-heading-rank; base=h2

\subsection{第四章——如何解决歧义}

8. 不仅这些，而且那些既不涉及标点也不涉及发音的含混之处，也应以同样方式加以审视。例如\WorkTitle{得撒洛尼人书}中的这一处：\Emph{因此，我们在你们中得到了安慰，弟兄们}。现在\Emph{弟兄们}是呼格还是宾格有疑问，而从信德的观点看，两种理解都不违背。但在希腊语中，两种格的形式不同；因此，当我们查看原文时，显示为呼格。如果译者选择说：\Emph{因此，我们在你们中获得了安慰，弟兄们}，他就不那么逐字直译，但含义的疑问会少一些；或者，如果他加上\Emph{我们的}，那么听到\Emph{因此，我们在你们中得到了安慰，我们的弟兄们}时，几乎无人会怀疑是呼格。但这是一个相当危险的自由。然而，在\WorkTitle{格林多人书}的那段经文里，宗徒说：\BibleQuote{我指着我们在我们的主基督耶稣内所有的夸耀\Emph{per vestram gloriam}起誓，我天天冒死}\BibleRef{格前 15:31}。有一位译者译作\Emph{我指着你们的夸耀起誓}，因为希腊文中起誓的形式没有歧义。因此，在\OriginalTerm{本义词}{proper words}的情况下，至少就圣经而言，几乎很难发现任何上下文（显示作者的意图）、或译文对照、或原文查考所不能解释的歧义。\par

% structure: p0014 source=h2 target=subsection reason=relative-heading-rank; base=h2

\subsection{第五章——将圣经的比喻表达按字面理解是一种可悲的奴役。}

9. 但是，比喻词语的含混之处，我接下来要谈论，需要格外的谨慎和勤勉。首先，我们必须警惕将比喻表达按字面理解。因为宗徒的话在此也同样适用：\BibleQuote{文字使人死，神却使人活}\BibleRef{格后 3:6}。因为当比喻说的话被当作字面话理解时，就会被按血肉的方式领会。没有什么比这更恰当地被称为灵魂的死亡：当灵魂中使之超越禽兽的部分，即理智，因盲目依附文字而屈服于血肉时。因为跟随文字的人把比喻词当作本义词，而不将本义词所指示的东西引申到其派生意义；例如，他听到安息日，只想到七天中按固定顺序重复的那一天；他听到祭献，思想便不越出从羊群取来的惯常牺牲和地上的果实。把标记当作事物，不能将心灵的眼光提升到形体和受造物之上，以畅饮永恒之光，这无疑是灵魂的一种可悲的奴役。\par

% structure: p0016 source=h2 target=subsection reason=relative-heading-rank; base=h2

\subsection{第六章——犹太人奴役的益处。}

10. 这种奴役，在犹太人民的情形中，与在其他民族中截然不同；因为，前者虽然束缚于暂时的事物，却以这样一种方式：在所有这一切中，独一的天主被置于他们心灵面前。虽然他们关注精神现实的标记而忽略了现实本身，不知道这些标记所指，但他们仍有此信念根植于心中：他们服从这种奴役是在取悦那独一不可见的天主。宗徒描述这种奴役如同孩童在启蒙师指导下的奴役。那些顽固依附这些标记的人不能忍受我们的主在启示时刻到来时对他们的忽视；因此，他们的首领控告他在安息日治病，而那些百姓，把这些标记当作现实紧紧抓住，不能相信一位拒绝像犹太人那样遵守这些标记的人是天主或来自天主。但那些相信的人——最初的耶路撒冷教会就是从他们中形成的——清楚地表明，被启蒙师如此引导是何等大的益处：标记，曾暂时施加于服从者，使观察者的思想固定在敬拜那创造天地的独一真天主上。这些人，因他们已非常接近精神事物（即使在暂时和血肉的祭献和预表中，虽然他们不清楚地理解其精神意义，却已学会敬拜那独一永恒的天主），竟充满了如此多的圣神，以至于他们变卖所有财物，将价银放在宗徒脚前分给穷人\BibleRef{宗 4:34}，并完全献身于天主，作为新的圣殿，而他们之前所事奉的旧圣殿不过是地上的预表。\par

11. 如今，没有记载任何外邦人的教会这样做，因为那些以人手所造的偶像为神的人，并没有如此接近精神事物。\par

% structure: p0019 source=h2 target=subsection reason=relative-heading-rank; base=h2

\subsection{第七章——外邦人无用的奴役。}

如果任何时候他们中的任何人试图证明他们的偶像只是记号，但他们仍将这些记号用于对受造物的敬拜和崇敬。举例来说，对我来说有何区别呢？即如尼普顿的像本身不被视为神，而只代表辽阔的海洋以及所有其他由泉源涌出的水？正如他们中一位诗人所描述的，若我记得不错：\Quote{你，尼普顿父神，你那灰白的额角被喧响的海洋环绕，你的胡须是无尽流淌的浩瀚大洋，你的发丝是蜿蜒的江河。}这豆荚在甜美的外壳内摇动碎石，但它不是人的食物，而是猪的。知道福音者明白我的意思。\BibleRef{路 15:16}那么，尼普顿的像用于这种解释对我有何益处呢？除非结果是我既不崇拜那像也不崇拜那解释所指的东西？因为任何你选取的雕像，对我来说，和那辽阔的海洋一样都是神。然而，我承认，那些以人的作品为神的人，比那些以天主的作品为神的人更低。但诫命是我们要爱慕和事奉独一的天主，祂是这一切的创造者，外邦人将其像作为神或作为神的记号与代表来崇拜。如果为了一个有用目的而设立的记号，被当作它本来要指代的事物本身来接受，那是肉体的奴役，那么把旨在代表无用之物的记号当作事物本身，岂不更是如此！因为即使你回到这些记号所指的事物本身，并让你的心智沉浸于对这些事物的崇拜，你也丝毫不会摆脱肉体奴役的重担和制服。\par

% structure: p0021 source=h2 target=subsection reason=relative-heading-rank; base=h2

\subsection{第八章 犹太人以一种方式脱离奴役，外邦人以另一种方式。}

12. 因此，由基督而来的自由将那些在有用记号之下受奴役的人（可以说是接近这自由的）找到，并解释他们所束缚的记号，藉着将他们提升到这些记号所指的实在，使他们获得自由。由这样的人形成了以色列圣者的教会。另一方面，它将那些在无用记号之下受奴役的人找到，不仅使他们从这种记号的奴役中解放出来，而且将这些记号本身完全废除和清除，从而使外邦人从众多假神的败坏（圣经常常且恰当地称之为淫乱）转离，去崇拜独一的天主：不是使他们现在陷入有用记号的奴役，而是使他们的心智操练于这些记号的属灵理解。\par

% structure: p0023 source=h2 target=subsection reason=relative-heading-rank; base=h2

\subsection{第九章 谁受制于记号，谁不受制}

13. 凡是使用或尊敬任何有意义之物却不明白其所指者，便是受制于记号；反之，凡使用或尊敬天主所立的有用记号，且明白其力量与意义者，所尊敬的并非那可见而短暂的记号，而是所有这类记号所指的实在。这样的人即使在受奴役之时也是属神而自由的，那时尚未适宜向属肉心的心灵启示那些记号，因为他们的肉性需要藉着服从这些记号而被克服。圣祖和先知，以及所有在以色列人民中藉着他们的服务圣神将圣经的助益与安慰赐予我们的人，都属于这类属神的人。但在现今，在我们主的复活中，我们自由的明证已如此清楚地照耀之后，我们不再受重负压迫，要留心那些我们现在已明白的记号；反而我们的主自己，以及宗徒的实践，将少数仪式代替许多传给了我们，这些仪式既极易施行，其意义又极为庄严，在遵守上最为神圣；例如，圣洗圣事和主的身体与血的庆典。任何人一看这些仪式，就知道它们所指为何，因此不是以肉体的奴役，而是以属神的自由来尊敬它们。按字面行事，将记号当作它们所指的事物，是软弱和奴役的标志；而错误地解释记号，则是被谬误误导的结果。然而，凡不明白一个记号所表示的是什么，却知道它是一个记号的人，并不处于奴役。即使受制于未知但有用的记号，也比通过错误地解释记号，使脖颈从奴役的轭下脱出，却插入谬误的缠绕中更好。\par

% structure: p0025 source=h2 target=subsection reason=relative-heading-rank; base=h2

\subsection{第十章 如何辨别一个短语是否为比喻性的}

14. 但除了前述规则（防止我们将比喻性的言谈当作字面来理解）之外，我们还必须注意另一条规则，即告诫我们不可将字面形式的言谈当作比喻性的。首先，我们必须指出发现一个短语是字面还是比喻的方法。这个方法当然如下：凡在天主圣言中，若按字面理解不能归于生活的纯洁或教义的健全，你便可判定为比喻。生活的纯洁涉及爱天主和爱近人；教义的健全涉及认识天主和认识近人。此外，每个人都在自己的良心中怀有希望，因为他觉察自己已达到对天主和近人的爱和认识。这一切已在第一卷中论述过了。\par

15. 然而，由于人倾向于不是根据罪本身的邪恶性来评判罪，而是根据他们自己的习俗来评判，因此常常发生这样的情况：一个人只会认为他自己国家和时代的人们习惯谴责的事是可指责的，只会认为他同伴的习俗所认可的事是值得赞扬或认可的；于是便出现这样的情形：如果圣经要么命令了与听众习俗相违背的事，要么谴责了与听众习俗不相违背的事，而同时圣经的权威又抓住了他们的心，他们就认为这种表达是比喻性的。实际上，圣经除了爱德之外不命令任何事，除了私欲之外不谴责任何事，并以此方式塑造人的生命。同样，如果一种错误见解占据了心灵，人们就认为圣经对此所断言的一切都必须是比喻性的。其实，圣经所断言的只是公教信仰，涉及过去、将来和现在的事。它是过去的叙述、将来的预言、现在的描述。但这一切都是为了滋养和坚固爱德，并克服和根除私欲。\par

16. 我所说的爱德，是指那种以天主本身为享受对象，并为了天主而享受自己和近人的心灵情感；我所说的私欲，是指那种不顾天主而只顾享受自己和近人以及其他肉体之物的心灵情感。再者，当私欲未被克制时，它所做的败坏自己灵魂和身体的事，称为\OriginalTerm{恶习}{vice}；但它所做的伤害他人的事，称为\OriginalTerm{罪行}{crime}。这就是所有罪可能被划分的两个类别。恶习先出现；因为当恶习耗尽了灵魂，使其陷入某种贫乏时，它便容易滑向罪行，以便消除障碍或为其恶习寻求帮助。同样，爱德为了自己的益处所做的，称为\OriginalTerm{智德}{prudence}；但为了近人的益处所做的，称为\OriginalTerm{仁爱}{benevolence}。这里智德在先；因为没有人能将自己未拥有的益处给予他人。私欲的统治被摧毁多少，爱德的统治就建立多少。\par

% structure: p0029 source=h2 target=subsection reason=relative-heading-rank; base=h2

\subsection{第11章——解释似乎将严厉归于天主和圣人的语句的规则。}

17. 因此，凡圣经中归于天主或其圣人的任何严厉和明显的残酷（无论是言语还是行为），都有助于摧毁私欲的统治。如果其含义明确，我们就不应将其视为比喻性的说法而另作引申。例如，宗徒的这句话：\BibleQuote{但你固执而不悔改，是在为自己积蓄忿怒，以备忿怒的日子和天主公义审判的显现；祂要照各人的行为报应各人：凡恒心行善，寻求荣耀、尊贵和不朽的，就以永生报应他们；至于那些好争斗、不顺从真理、反顺从不义的人，就以忿怒和恼恨、患难和困苦报应他们，加给一切作恶的人，先是犹太人，后是外邦人。}\BibleRef{罗 2:5} 但这是对那些不愿克制自己私欲、因而自身陷入私欲毁灭的人说的。然而，当私欲在其曾统治的人身上被推翻时，便用了这个明确的表达：\BibleQuote{凡属于基督的人，已把肉身连同邪情私欲钉在十字架上了。}\BibleRef{迦 5:24} 只是即使在这些例子中，也有些词语是比喻性地使用的，例如\Quote{天主的忿怒}和\Quote{钉在十字架上}。但这些并不那么多，也不至于遮蔽含义而使之成为寓言或谜语——这才是真正称为\OriginalTerm{比喻性表达}{figurative}的那类说法。但在对耶肋米亚说的话中：\BibleQuote{看，我今天立你于万民和列国之上，为要拔除、拆毁、毁灭、推翻，}\BibleRef{耶 1:10} 毫无疑问，整个语言都是比喻性的，应指向我所说的目的。\par

% structure: p0031 source=h2 target=subsection reason=relative-heading-rank; base=h2

\subsection{第12章——解释归于天主和圣人、但在无经验者看来似乎邪恶的那些言论和行为的规则。}

18. 再者，那些在无经验者看来似乎有罪的事（无论是仅凭言语还是实际行为），若归于天主或那些其圣德被放在我们面前作为榜样的人，则完全是比喻性的，其中隐含的意义内核要像食物一样被取出，以滋养爱德。任何人若使用短暂之物比其周围人的习俗更不自由，他要么是节制的，要么是迷信的；反之，若使用这些事物而越过了周围善人的习俗界限，那么他要么是在他所做的事中有更深含义，要么是有罪的。在所有这些事情中，该受责备的不是对事物的使用，而是使用者的私欲。例如，任何人头脑清醒都不会相信，当主的脚被那女人用珍贵的香液傅抹时（\BibleRef{若 12:3}）是出于那些我们憎恶的宴会上奢侈放荡之人惯常涂抹的目的。因为芬芳的香气意味着善行生活赢来的好名声；而赢得这名声的人，在跟随基督的足迹时，可谓用最珍贵的香液傅抹了祂的脚。所以，那些在其他人身上常常是罪的事，当归于天主或先知时，便成了某种伟大真理的标记。例如，与娼妓交往，在放荡的品行中是另一回事，而在欧瑟亚先知履行预言时所做的（\BibleRef{欧 1:2}）则是另一回事。因为在醉酒和放荡的宴会上赤身露体是极其邪恶的事，但这并不意味着在浴室中赤身是罪。\par

19. 因此，我们必须慎重考虑何者适合时间、地点与人，不可轻易指控他人犯罪。因为智者享用最精美的食物，也可能毫无饕餮或贪食之罪；而愚人可能以最可憎的贪欲渴求最粗劣的食物。任何心智健全的人，都宁愿效法主吃鱼，而不愿效法厄撒乌吃扁豆，或效法牛吃大麦。因为有些野兽以较普通的食物为生，但这并不表示它们比我们更有节制。在这一切事上，决定我们行为可赞或可责的，不是所用之物的本性，而是我们使用它们的理由和寻求它们的方式。\par

20. 古代圣人，在尘世王国的形式下，预表预言了天国。因需要众多后裔，那时一人多妻的风俗无可指责；同理，一妻多夫却不适宜，因为妇女不会因此更生育，相反，通过滥交寻求利益或后代乃是卑劣的淫乱。关于此类事，无论当时圣人毫无情欲地做了什么，圣经都未加责备，尽管他们做了现今除非出于情欲不能做的事。凡此一切记载，我们不仅应按历史字面理解，也应按象征预言意义理解，并解释为最终指向爱天主或爱近人、或兼爱二者的目标。正如古罗马人穿长及脚跟并带袖子的长袍是可耻的，但现在有身份的人不穿这样的长袍却是可耻的；同样，我们也要注意其他事物，勿让情欲掺杂于其中；因为情欲不仅滥用所居之地的风俗以达到邪恶目的，而且常常越出习俗的界限，在可耻的爆发中暴露其隐藏在流行风尚下的丑恶。\par

% structure: p0035 source=h2 target=subsection reason=relative-heading-rank; base=h2

\subsection{第十三章——继续同一主题。}

21. 凡与我们所处时代的人们的习惯相符之事——这些人，我们或因必要被迫，或因义务自愿，与之共度此生——善良伟大的人都应将其导向某种明智或仁爱的目的，或直接，如我们应尽的本分，或象征，如先知所允许的。\par

% structure: p0037 source=h2 target=subsection reason=relative-heading-rank; base=h2

\subsection{第十四章——认为没有绝对是非之人的错误。}

22. 然而，那些除了自己的生活方式外对其他生活方式一无所知的人，若遇到此类行为的记载，除非他们受权威约束，就会视之为罪；他们没有考虑到，他们自己在婚姻、筵席、衣着或其他人生必需品与装饰上的习俗，在其他民族和其他时代的人看来，也是罪。他们被这无穷无尽的习俗差异所困惑，有些人半睡半醒——也就是说，既未沉入愚昧的沉睡，又不能醒来进入智慧之光——便认为根本没有绝对的是非，每个民族都以自己的习俗为是非；既然每个民族有不同的习俗，而是非必须恒常不变，那么显然根本不存在是非。这些人没有意识到，仅举一例，诫命\BibleQuote{凡你们愿意人给你们做的，你们也要照样给人做}不能因任何民族习俗的差异而改变。这条诫命，当指向爱天主时，摧毁一切恶习；当指向爱近人时，杜绝一切罪行。因为没有人愿意玷污自己的住所；因此，他也不应玷污天主的住所，即他自己。没有人愿意他人伤害自己；因此，他自己也不应伤害他人。\par

% structure: p0039 source=h2 target=subsection reason=relative-heading-rank; base=h2

\subsection{第十五章——解释比喻表达之规则。}

23. 情欲的暴政既已推翻，爱德便以其至公至义的法则统治：为天主本身而爱祂，并为天主而爱自己和近人。因此，对于比喻表达，应遵循如下规则：用心反复思索我们所读的，直到找到一个有助于巩固爱德统治的解释。如果照字面理解立刻就能得出这样的意义，那么该表达就不应视为比喻。\par

% structure: p0041 source=h2 target=subsection reason=relative-heading-rank; base=h2

\subsection{第十六章——解释命令与禁令之规则。}

24. 若语句是命令，或禁止罪行或恶习，或命令智慧或仁爱之举，则非比喻。然而，若它似乎命令一项罪行或恶习，或禁止智慧或仁爱之举，则是比喻。基督说：\Quote{你们若不吃人子的肉，不喝他的血，在你们内便没有生命。}（\BibleRef{若 6:53}）这似乎命令一项罪行或恶习；因此它是一个比喻，命令我们应当有分于（\OriginalTerm{分享}{communicandem}）主的苦难，并应保留一种甜蜜而有益的回忆（\OriginalTerm{在记忆中}{in memoria}），即他的肉曾为我们被刺穿并被钉十字架。圣经说：\Quote{你的仇人若饿了，就给他吃；若渴了，就给他喝。}这无疑是一项命令行善。但接着：\Quote{因为你这样做，是将炭火堆在他头上。}有人会认为这是在命令行恶。那么，不要怀疑这表达是比喻性的；并且，尽管有两种解释可能，一种指向加害，另一种指向显示优越，爱德却相反地要召唤你回到仁爱，并将炭火解释为痛悔的燃烧呻吟，由此治愈人的骄傲，他哀叹自己曾是那在困苦中帮助他的人之仇敌。同样，当主说：\Quote{爱惜自己性命的，必丧失性命。}我们不应认为他禁止人应照顾自己生命的智慧，而是他用比喻说：\Quote{让他丧失性命}——即让他摧毁并失去他目前对自己性命的扭曲而反常的运用，借此他的欲望专注于暂时之事，以致他不留心永恒。经上记载：\Quote{施舍给虔诚人，不要帮助罪人。}这句话的后半似乎禁止仁爱；因为它说：\Quote{不要帮助罪人。}因此，要明白\Quote{罪人}是比喻地指罪恶，即你不可帮助的是他的罪。\par

% structure: p0043 source=h2 target=subsection reason=relative-heading-rank; base=h2

\subsection{第十七章——有些命令是普遍给予所有人的，其他是特别给予某类人的。}

25. 此外，常有这样的事：一个人达到或自认为达到更高层次的神修生活，便认为给予较低层次之人的命令是比喻性的；例如，若他选择了独身生活，并因天国而自阉，他便主张圣经中关于爱妻与治妻的命令不应按字面理解，而应按比喻理解；若他决心使女儿守贞，他便试图对这话作比喻解释：\Quote{应将你的女儿出嫁，你就完成了一件大事。}（\BibleRef{德 7:27}）因此，我们理解圣经的另一条规则如下——承认有些命令是普遍给予所有人的，有些是特别给予某类人的，使药物不仅作用于整体健康状态，也作用于每个成员的特别软弱之处。因为那不能提升到更高层次的，必须在自身层次内得到照顾。\par

% structure: p0045 source=h2 target=subsection reason=relative-heading-rank; base=h2

\subsection{第十八章——我们必须考虑某事被享有或允许的时间。}

26. 我们还必须当心，不要以为旧约中因着那时代的条件而按字面理解而非比喻理解也不算罪行或恶习的事，可以转用于现今作为生活习性。因为除非私欲管辖他，并试图在原本要推翻它的圣经本身中寻找支持，否则无人会这样做。那不幸的人并未察觉，这些事之所以记录，有这有益的目的：使怀有好望的人学到有益的教训，既知道他们所轻视的习俗可以用于善，也知道他们所拥抱的习俗可被用于定罪，若前者之运用伴随着爱德，后者之运用伴随着私欲。\par

27. 因为，若一人能以贞洁使用许多妻子，则另一人能以私欲使用一个妻子。我更赞许那为着更远的目的而使用多妻之生育能力的人，胜于那为着自身而享受一个妻子身体的人。在前者，那人追求符合时代情况的实用目的；在后者，他满足专注于暂时享受的私欲。那些因着他们的不节制，宗徒以宽恕允许他们有一妻的人，不如那些虽各有许多妻子，却如智者使用食物和饮料仅为了身体健康那样，仅为了后嗣而使用婚姻的人亲近天主。因此，若这些人仍活到我们主的来临，那时不再是抛掷石头而是聚集石头的时候（\BibleRef{训 3:5}），他们会立即因天国而自阉。因为除非在享受中有私欲，否则禁欲并不困难。确实，我所说的那些人知道，即使在妻子方面，放纵也是滥用和不节制，这由多俾亚与妻子结婚时的祈祷所证明。因为他祈祷说：\Quote{我们的祖先的天主，你是应受赞美的；你的圣名永远应受赞美和光荣；诸天和你的所有造物都要赞美你。你造了亚当，又造了厄娃作他的妻子和助手。……上主，现在你知道我娶这个妹妹不是由于情欲，而是出于真诚。所以求你怜悯我们，上主。}（\BibleRef{多 8:5}）\par

% structure: p0048 source=h2 target=subsection reason=relative-heading-rank; base=h2

\subsection{第十九章——恶人以己度人。}

28. 然而，那些放纵情欲的人，或四处漫游沉溺于各种淫乱，或即使对待一位妻子，不仅超出了生育子女所需的尺度，而且以某种奴隶般自由的无耻放纵，堆积更兽性的污秽，这样的人不相信古时的人节制地使用多位妻子，只着眼于当时环境所必需的繁衍种族的责任；而他们自己，被情欲的网罗所缠，连对一位妻子都做不到的事，他们认为对多位妻子根本不可能。\par

29. 但同样是这些人，可能会说，连尊崇和赞美善良圣洁的人也不应该，因为他们自己受到尊崇和赞美时，就会骄傲自大，因为频繁且广泛地听受阿谀之辞，就越是渴求最空洞的荣耀，变得如此轻浮，以致一阵谣言之风，无论看似有利还是不利，都会把他们卷入恶习的漩涡，或撞在罪恶的岩石上。那么，让他们学习，要避免被赞美的诱饵所捕获，或被侮辱的刺所刺痛，是多么困难和艰巨的事；但不要以己度人。\par

% structure: p0051 source=h2 target=subsection reason=relative-heading-rank; base=h2

\subsection{第二十章——善人在一切外在环境中始终如一。}

相反，让他们相信，我们信仰的宗徒们既不会因受人的尊崇而高傲，也不会因受人的鄙视而沮丧。当然，这两种试探对那些伟人来说都不缺乏。因为他们既受到信徒们高声赞美的颂扬，也受到迫害者诽谤的贬低。但宗徒们按时机利用这一切，却不被腐蚀；同样，古代的圣人们按他们时代的需要运用他们的妻子，并不像那些拒绝相信这些事的人那样受情欲的捆绑。\par

30. 因为如果他们曾受任何这类情欲的影响，他们绝不可能克制自己，不对他们的儿子怀有刻骨的仇恨，因为他们知道儿子们曾勾引和奸污他们的妻子和妾。\par

% structure: p0054 source=h2 target=subsection reason=relative-heading-rank; base=h2

\subsection{第二十一章——达味虽犯奸淫，却非好色。}

但当达味王遭受他那不敬虔和逆性的儿子如此伤害时，他不仅容忍了他的疯狂情欲，还在他死后哀悼他。他确实没有陷入肉体的嫉妒之网，因为感动他的不是他自己的损害，而是他儿子的罪恶。正是为此，他曾下令，如果他的儿子在战败时，不可杀死他，以便他降服后能有悔改的机会；当这个计划落空时，他哀悼他儿子的死，不是因为自己的损失，而是因为他知道如此不敬虔的通奸者和弑亲者被投入了何等的惩罚。因为在这之前，在另一个没有犯罪的儿子身上，虽然他在患病时极其担忧，但在儿子死后，他安慰了自己。\BibleRef{撒下 12:19}\par

31. 那些人在使用他们的妻子时有何等的节制和克己，这主要表现在：当这位同一位国王，被情欲的狂热和暂时的顺境所冲昏头脑，非法占有了一个女人，并下令杀死了她的丈夫时，一位先知指责了他的罪恶；这位先知来向他指明他的罪，给他讲了那个穷人的比喻，这穷人只有一只母羊羔，而他的邻舍虽有许多，但当客人来到时，却不肯从自己的羊群中取用，反而把穷人的那只母羊羔摆给客人吃。达味对那人发怒，命令将他处死，并叫那穷人得到四倍的偿还；这样，他无意中定了自己故意所犯的罪。\BibleRef{撒下 12:1}当先知向他指明这点，并宣布了天主的惩罚后，他以深切的忏悔洗去了他的罪。但在那个比喻中，穷人的母羊羔只表示奸淫；关于杀害那女人的丈夫——即拥有那只母羊羔的穷人的谋杀——在比喻中并未提及，因此定罪判决仅针对奸淫。由此我们可以理解，当他因在一位女人身上犯过而被迫惩罚自己时，他拥有多位妻子是多么有节制。但在他的情况下，过度的欲望并未在他里面定居，只是过客。因此，那不合法的欲望甚至被那位控诉的先知称为\Quote{客人}。因为他没有说，他取了穷人的母羊羔为他的王摆设宴席，而是为他的客人。至于他的儿子撒罗满，这种情欲却不像客人一样来去，而是像君王一样统治。关于他，圣经并不沉默，而是指责他喜爱外邦女子；因为在他统治之初，他燃烧着对智慧的热望，但通过属灵的爱获得智慧后，却因肉体的情欲而失去了它。\par

% structure: p0057 source=h2 target=subsection reason=relative-heading-rank; base=h2

\subsection{第二十二章——关于圣经中那些表达对现今善人所谴责之行为的赞许的段落，所应遵循的规则。}

32. 因此，虽然旧约中记载的全部或几乎全部事件，不仅要按字面理解，也要按喻意理解，然而，即使读者按字面理解了那些事件，并且虽然事件的作者受到赞扬，但这些事件与自吾主降世以来作为天主诫命守护者的善人的习惯相悖，那么，让他将喻象指向其解释，但不要将行为转移到自己的生活习惯中。因为许多在那些时代作为义务而做的事，如今只能出于情欲而做。\par

% structure: p0059 source=h2 target=subsection reason=relative-heading-rank; base=h2

\subsection{第二十三章——关于伟人罪恶叙述的规则。}

33. 当他读到伟人的罪过时，虽然他或许能在其中看出并追溯未来事物的预象，但也应以此历史事实来教导他，不敢因自己的善行而自夸，并因自己的义而轻视他人为罪人，因为他看到在如此杰出的人物身上，既有应避免的风暴，也有应哀悼的沉船。因为这些人的罪过被记载下来，为的是使人们无论何时何地都对宗徒的那句话感到战惊：\Quote{所以，那自以为站得稳的，要小心免得跌倒。}\BibleRef{格前 10:12}因为圣经中几乎没有一页不明确写着天主拒绝骄傲的人，却赐恩宠给谦卑的人。\par

% structure: p0061 source=h2 target=subsection reason=relative-heading-rank; base=h2

\subsection{所用言辞的特质是最重要的。}

34. 因此，我们试图理解的任何表达，主要应探究的是，它是字面的还是比喻的。因为当确定它是比喻的，就容易通过应用我们在第一书中所讨论的物之律法，以各种方式转换，直到达成真正的解释，尤其是当我们凭借虔敬的锻炼而加强的经验来协助时。现在我们通过注意上述的考虑来发现一个表达是字面的还是比喻的。\par

% structure: p0063 source=h2 target=subsection reason=relative-heading-rank; base=h2

\subsection{同一个字并不总是指同一事物。}

当它被证明是比喻的，表达它的言辞将被发现取自相似的对象或具有某种相似性的对象。\par

35. 但是，由于事物之间有许多相似的方式，我们不应认为存在这样的规则：在一处以相似性表示的事物在所有其他地方都必须表示同一事物。因为我们的主使用了\OriginalTerm{酵母}{leaven}，既有坏的意义，如祂说：\Quote{你们要谨慎法利塞人的酵母}，也有好的意义，如祂说：\Quote{天国好像酵母，女人取来藏在三斗面里，直到全部发了酵。}\BibleRef{路 13:21}\par

36. 关于这种变化的规则有两种形式。因为表示一种事物又表示另一种的事物，要么表示相反的事物，要么仅仅表示不同的事物。例如，当它们在一个地方用于好的意义，在另一个地方用于坏的意义时，表示相反的事物，如上述的\OriginalTerm{酵母}{leaven}。另一个例子是：\OriginalTerm{狮子}{lion}代表基督，如经上所说：\Quote{犹大支派的狮子已经得胜了}；\BibleRef{默 5:5}又代表魔鬼，如经上所说：\Quote{你们的仇敌魔鬼，如同咆哮的狮子，巡游寻找可吞吃的人。}\BibleRef{伯前 5:8}同样，\OriginalTerm{蛇}{serpent}用于好的意义：\Quote{你们要机警如蛇}；\BibleRef{玛 10:16}又用于坏的意义：\Quote{蛇用它的狡猾诱惑了厄娃。}\BibleRef{格后 11:3}\OriginalTerm{饼}{bread}用于好的意义：\Quote{我是从天上降下的生活的饼}；\BibleRef{若 6:51}用于坏的意义：\Quote{偷来的饼是甘甜的。}\BibleRef{箴 9:17}如此等等其他许多例子。我所引用的例证在意义上确实毫无可疑之处，因为只应使用明确的例子。然而，有些段落，在应如何理解的意义上是不确定的，例如：\Quote{上主手里有杯，酒色赤红，满斟混合之酒。}现在不确定这是否表示天主的义怒，但不至于最后的惩罚，即\Quote{到渣滓}；或者表示圣经的恩宠从犹太人转移到外邦人，因为\Quote{祂推倒一个，又立起另一个}——某些礼仪仍然保留在犹太人中间，他们以血肉的方式理解这些礼仪，因为\Quote{其中的渣滓还没有倒尽。}以下是同一对象被理解为并非相反，而只是不同意义的例子：\OriginalTerm{水}{water}表示群众，如我们在\BibleBook{}{默示录}中所读，\BibleRef{默 17:15}也表示圣神，例如：\Quote{从他的腹中要流出活水的江河。}\BibleRef{若 7:38}其他许多事物除了水以外，也应根据其所在之处来解释。\par

37. 同样，其他对象也不是单一的意义，而是每一个不仅表示两个，有时甚至表示几个不同的事物，取决于上下文。\par

% structure: p0068 source=h2 target=subsection reason=relative-heading-rank; base=h2

\subsection{晦涩的段落应由更清晰的段落来解释。}

现在，从这些词意义更明显的地方，我们必须收集它们在晦涩段落中应如何理解的意义。例如，理解向天主所说的话\Quote{拿起\OriginalTerm{盾牌}{shield}和\OriginalTerm{圆盾}{buckler}，为我起来帮助}，最好的办法就是参考我们读到的段落：\Quote{上主，祢以恩宠如同盾牌般环绕我们。}然而，我们不应这样理解，以为无论何处遇到盾牌表示任何保护，我们就必须认为它只表示天主的恩宠。因为我们也听到关于信德的\Emph{\OriginalTerm{胸甲}{breastplate}}，宗徒说：\Quote{你们要拿起信德作盾牌，使你们能熄灭恶者的一切火箭。}\BibleRef{弗 6:16}另一方面，关于这类属灵武器，我们也不应将信德只归于盾牌；因为我们在另一处读到信德的\Emph{\OriginalTerm{胸甲}{breastplate}}：宗徒说：\Quote{穿上信德和爱德的胸甲。}\par

% structure: p0070 source=h2 target=subsection reason=relative-heading-rank; base=h2

\subsection{一段文字可有多种解释。}

38. 再者，若对圣经同一句话不是只有一个解释，而是有两个或更多解释，即使作者的原意未被发现，只要能由圣经其他段落证明其中的任何一种解释合乎真理，就没有危险。一个人研读圣经，努力领会圣神藉其发言的作者的原意，无论他在这努力中成功，还是从词句中得出不同的含义，但若这含义不违背健全的教义，只要他有圣经其他段落的见证支持，便无可责难。因为作者或许已经看到，我们试图解释的词句中正包含着这层含义；而藉着他讲出这些话的圣神，肯定预见到读者会产生这一解释，甚至促成它产生，因为这解释也建立在真理之上。天主对待圣经，还能有比这更宽厚、更丰盛的安排吗？即同一句话可以有多种理解，而且这些理解都得到同样神圣的其他段落的一致见证的认可。\par

% structure: p0072 source=h2 target=subsection reason=relative-heading-rank; base=h2

\subsection{第二十八章——用圣经其他段落解释可疑之处比用理性更稳妥。}

39. 然而，当所推出的含义无法用圣经确凿的证据澄清时，我们便只好用理性的证据来阐明它。但这是危险的做法。因为行走在圣经的光照下要稳妥得多；这样，当我们想要审视那些被隐喻表达所遮蔽的段落时，要么获得毫无争议的含义，要么一旦发生争论，便通过从同一圣经各处寻找的见证来解决它。\par

% structure: p0074 source=h2 target=subsection reason=relative-heading-rank; base=h2

\subsection{第二十九章——认识修辞格的必要性。}

40. 此外，我希望有学问的人知道，我们圣经的作者使用文法学家以希腊名称\OriginalTerm{修辞格}{tropes}称呼的所有那些表达形式，而且使用得比不熟悉圣经、从其他著作学会这些修辞格的人所能想象或相信的更自由、更多样。然而，懂得这些修辞格的人在圣经中认出它们，并因对它们的认识而大大有助于理解圣经。但这不是向无知者教授这些的地方，以免显得我是在教文法。不过，我确实建议在别处学习它们，正如我在第二卷中已经建议过的那样——即，在我论述必要语言知识的地方。因为文法本身得名的文字（希腊文中\Quote{字母}一词为\OriginalTerm{文字}{\GreekText{γράμματα}}）是我们说话时清晰的声音的符号。关于这些修辞格，我们在圣经中不仅找到例子（我们拥有它们全部的例子），而且还找到它们本身的名称：例如，\OriginalTerm{寓意}{allegory}、\OriginalTerm{谜语}{enigma}和\OriginalTerm{比喻}{parable}。然而，几乎所有据说通过自由教育学会的这类修辞格，甚至在没有学过文法、只满足于使用俗语的人的日常言语中也能找到。有谁不说\Quote{祝您如此繁荣？}？这就是称为\OriginalTerm{隐喻}{metaphor}的修辞格。有谁不说一个没有鱼、不是为鱼而挖、却因鱼而得名的\Quote{鱼池}？这就是称为\OriginalTerm{换称}{catachresis}的修辞格。\par

41. 若这样逐一列举其余所有修辞格将十分冗长；因为俗人的言语使用它们全部，甚至包括那些意义与字面相反的更奇特的修辞格，例如称为\OriginalTerm{反讽}{irony}和\OriginalTerm{反语}{antiphrasis}者。在反讽中，我们通过语调来表达我们想传达的含义；就如我们对一个行为恶劣的人说：\Quote{你做得很好。}但反语并非通过语调来表明与字面相反的含义；而是要么其中使用的词语与其词源意义相反，如林丛（\OriginalTerm{林丛}{lucus}）因其缺乏光亮而得名；要么习惯使用某种表达方式，尽管它通过相反律以\Quote{是}代\Quote{否}，例如我们在一个地方询问那里没有的东西，得到的回答是\Quote{多得很。}；或者我们添加一些词语，明确表示我们说的恰恰相反，如\Quote{提防他，因为他是个好人。}有哪个没有学问的人不使用这类表达，尽管他对这些修辞格的性质和名称一无所知？然而，认识它们在澄清圣经的难点时是必要的；因为当按字面理解词语得出荒谬含义时，我们应当立即查问它们是否可能以某种我们不熟悉的比喻意义被使用；许多晦涩的段落就这样得到了阐明。\par

% structure: p0077 source=h2 target=subsection reason=relative-heading-rank; base=h2

\subsection{第三十章——考辨多纳图派提科尼乌斯的规则。}

42. 有一位\Person{提科尼乌斯}，他本人虽是多纳徒派，却极其有力地撰文反对多纳徒派（在此显示出极不一致的性情，因为他不愿完全放弃他们）。他写了一部书，称为\WorkTitle{规则之书}，因为他在其中提出了七条规则，这些规则就像是开启圣经奥秘的钥匙。这些规则中，第一条涉及主与祂的身体，第二条涉及主身体的双重划分，第三条涉及恩许与法律，第四条涉及\OriginalTerm{种类}{species}与\OriginalTerm{属类}{genus}，第五条涉及时间，第六条涉及重述，第七条涉及魔鬼与其身体。作者所阐述的这些规则，若仔细考量，确实对洞悉圣经的奥秘大有助益，但它们并未解释所有难解之处，因为还需要其他多种方法，这些方法远非这七条规则所能涵盖，以至于作者本人解释许多晦涩段落时并未使用任何一条规则；他发觉其实并无此必要，因为那些段落并不存在他的规则所适用的困难。例如，他探讨在\BibleBook{}{默示录}中，\Person{若望}受命写信的七个教会的天使应当如何理解；经过大量而多样的论证，他得出结论：天使就是教会本身。在整篇详尽而充实的讨论中，尽管所探讨的问题确实非常晦涩，却完全没有用到这些规则。此例足以说明问题，因为若要收集正典圣经中所有无需这七条规则即可澄清的晦涩段落，那将过于冗长繁琐。\par

43. 然而，作者本人在推荐这些规则时，赋予了它们如此大的价值，以至于似乎一旦彻底了解并恰当运用它们，我们就能解释法律中——即圣经中——所有晦涩的段落。因为他在这本书的开始如此写道：\Quote{在我想到的所有事情中，我认为没有比写一本小规则书更必要的，仿佛为法律的神秘之处制造钥匙、开设窗户。因为有一些神秘的规则，它们掌握着整部法律秘密处的钥匙，并使那对许多人不显明的真理宝藏变得可见。若这规则体系如我所传授的那样被毫无嫉妒地接受，那么封闭的将被打开，晦涩的将被澄清，以致一个人在先知预言广大的森林中旅行时，若他循着这些规则如同光明的路径，必将免于走入歧途。} 现在，如果他曾说：\Quote{有一些神秘的规则掌握着法律某些秘密的钥匙}，甚或\Quote{掌握着法律重大秘密的钥匙}，而不是他实际所说的\Quote{整部法律秘密处的钥匙}；如果他未曾说\Quote{封闭的将被打开}，而是说\Quote{许多封闭的将被打开}，他就说了真话，并且他不会因为他给予他那精心实用之作超出事实的价值，而使读者产生错误的期待。我认为说这些是恰当的，既为了使这部书能被好学之人阅读（因为它在理解圣经上非常有助益），也为了使人对它不抱超过它实际内容的期待。当然，阅读时必须谨慎，不仅因为作者作为人会犯错，更主要因为他作为多纳徒派所提出的异端学说。现在我将简要说明这七条规则教导或建议什么。\par

% structure: p0080 source=h2 target=subsection reason=relative-heading-rank; base=h2

\subsection{\Emph{第三十一章 \Person{提科尼乌斯}的第一条规则}}

44. 第一条是有关\Emph{主与祂的身体}的，即这样：既然我们知道头和身体——即基督和祂的教会——有时被指示为同一位（因为对信徒们说\BibleQuote{你们就是亚巴郎的后裔}不是徒然的，见：\BibleRef{迦3:29}，而亚巴郎只有一个后裔，就是基督），那么当从头部转到身体或从身体转到头部，而说话的对象并未改变时，我们不必感到困难。因为同一个人说话：\BibleQuote{祂装饰了我如同新郎，佩戴了首饰；又装扮了我如同新娘，戴上了珠宝。}然而，这当然需要解释，这两处哪是指头，哪是指身体，即哪是指基督，哪是指教会。\par

% structure: p0082 source=h2 target=subsection reason=relative-heading-rank; base=h2

\subsection{\Emph{第三十二章 \Person{提科尼乌斯}的第二条规则}}

45. 第二条规则是关于\Emph{主身体的二分}；但这其实不是一个恰当的名称，因为真正不属于基督的那些部分不会与祂同在永恒。因此，我们应当说，这条规则是关于主之真实与混杂的身体，或真实与虚假的身体，或诸如此类的名称；因为，且不说永恒，就是现在，假善人也不能说是在祂内，尽管他们看似在祂的教会内。因此，这条规则可称为：关于\Emph{混杂的教会}。这条规则要求读者注意，当圣经现在转向或谈论另一群不同的人时，似乎仍在像以前一样对同样的人说话或谈论，就好像这两群人因暂时在圣事上共同参与而构成了一个身体似的。例如《雅歌》中的这句话：\BibleQuote{我虽黑，却秀美，如刻达尔的帐幕，如撒罗满的帷幔。}（\BibleRef{歌 1:5}）因为经上不是说，我曾黑如刻达尔的帐幕，但现在却秀美如撒罗满的帷幔。教会宣称自己目前两者都是；这是因为好鱼和坏鱼暂时同在一网中（\BibleRef{玛 13:47}）。刻达尔的帐幕属于依市玛耳，他\BibleQuote{不可与自由妇人的儿子一同承继产业}（\BibleRef{迦 4:30}）。同样，当天主论及教会中善的部分说：\BibleQuote{我要领瞎子走不认识的路，引他们走不认识的途径；我要使黑暗在他们面前化为光明，崎岖之地变成平坦；这些事我必实行，决不撇弃他们。}（\BibleRef{依 42:16}）他随即论及另一部分，即混杂在善人中的恶人说：\Quote{他们必退后。}这些话所指的是一群与前者完全不同的人；但因这两群人目前同在一个身体内，他便说话如同句子的主语没有改变。然而，他们不会永远在一个身体内；因为其中之一就是福音中记载的那个恶仆，他的主人回来时，\BibleQuote{要把他斩断，使他和假善人遭受同样的命运。}（\BibleRef{玛 24:50}）\par

% structure: p0084 source=h2 target=subsection reason=relative-heading-rank; base=h2

\subsection{第三十三章 提科尼乌斯的第三条规则}

46. 第三条规则关乎\Emph{恩许与法律}，也可以用其他术语称之为精神与文字，这正是我撰写论此主题的书时所使用的名称。它也可以名为恩宠与法律。然而，在我看来，这本身就是一个重大问题，而非仅用于解决其他问题的规则。正是对此问题缺乏清晰认识，才起源或至少大大加剧了白拉奇异端。提科尼乌斯在澄清这一点上作的努力虽好，却不完全。因为，在讨论信德与行为的问题时，他说行为是天主赐予我们作为信德的赏报，但信德本身却是我们自己的，不是来自天主；他没有记住宗徒的话：\BibleQuote{愿平安与爱德、连同信德，由天主父和主耶稣基督赐予弟兄们。}（\BibleRef{弗 6:23}）然而，他没有遇到在我们时代兴起的这种异端，它使我们付出了许多辛劳和麻烦来捍卫那藉着我们的主耶稣基督而来的天主的恩宠；而正如宗徒所说：\BibleQuote{你们中间也不免有异端，好让那些经得起考验的人显明出来。}（\BibleRef{格前 11:19}）这使我们更加警醒和勤奋地发现圣经中提科尼乌斯所忽略的东西——他因没有敌人要防备，在这方面就不太专注和迫切——即：连信德本身也是那\BibleQuote{按各自信德的分量分赐给各人}（\BibleRef{罗 12:3}）的天主的恩赐。因此，对某些信徒说过：\BibleQuote{因为你们为了基督，不仅得以相信祂，而且得以为祂受苦。}（\BibleRef{斐 1:29}）那么，当有人从这段经文中得知并相信这两者都是赐予的，谁还能怀疑它们都是天主的恩赐呢？此外还有许多其他证据证实这一点。但我现在不是在讨论这一教义。我已在别处多次论及它。\par

% structure: p0086 source=h2 target=subsection reason=relative-heading-rank; base=h2

\subsection{第三十四章 提科尼乌斯的第四条规则}

47. 提科尼乌斯的第四条规则关乎\Emph{类型与类别}。因为他如此称呼，其用意是：类型应理解为部分，类别则是整体——他所称的类型是其中的一部分：例如，每一个城市都是万民大社会的一部分；他称城市为类型，全体民族构成类别。这里没有必要应用逻辑学家常用的那种细微区分，他们极为敏锐地讨论部分与类型之间的差别。如果圣经中所发现的情况不是关于单个城市，而是关于某个省份、部落或王国，这条规则自然也一样适用。例如，不仅关于耶路撒冷，或关于外邦人的某些城市，如提洛或巴比伦，圣经所说的一些话超越了该城市的界限，更适用于所有民族；而且关于犹大、埃及、亚述，或你任意选择的、包含许多城市但仍不是整个世界而只是其一部分的任何民族，也有所说的话超出了该国的界限，更适切地适用于其所属的整体；或按我们的作者所说，适用于其作为类型的类别。因此，这些话已成为常识，甚至未受过教育的人也能理解在某个帝国命令中哪些是特别规定的，哪些是普遍适用的。同样的情况也发生在人身上：例如，关于撒罗满所说的一些话，其范围远远超出他，唯有应用于基督和祂的教会（撒罗满是其一部分）时才能正确理解（\BibleRef{撒下 7:14}）。\par

48. \OriginalTerm{寓体}{species}并非总是被超越，因为事物常被说成这样，以致显然也适用于它，或者甚至只适用于它。但是，当圣经在某种程度上一直谈论寓体，从该点由寓体转向\OriginalTerm{本体}{genus}时，读者必须谨慎提防，不要在寓体中寻求他在本体中可以更好、更确定地找到的东西。例如，先知厄则克耳所说的话：\BibleQuote{当以色列家族住在自己的土地上时，他们以自己的道路和行为玷污了它；他们的道路在我面前如同经期妇人的不洁。因此我向他们倾泄我的愤怒，因为他们在地上所流的血，以及他们用以玷污那地的偶像；我将他们分散在外邦人中，他们被散布在各国；我按他们的道路和行为审判了他们。}\BibleRef{则 36:17} 现在很容易理解，这适用于宗徒所说的那个以色列家族：\BibleQuote{你们看按血肉而生的以色列；}\BibleRef{格前 10:18} 因为按血肉而生的以色列人民确实做和忍受了这里所提及的一切。紧接的经文也可以理解为适用于同一人民。但当先知开始说：\BibleQuote{我要圣化我被异民所亵渎的大名，就是你们在他们中所亵渎的；外邦人将知道我是主，}\BibleRef{则 36:23} 读者现在应当仔细观察寓体如何被超越，本体如何被采纳。因为他接着说：\BibleQuote{我要在你们身上，在他们眼前显我为圣。我要从外邦人中领出你们，从各国聚集你们，领你们回到你们的土地。那时我要在你们身上洒清水，你们就洁净；我要洗净你们的一切污秽和一切偶像。我还要赐给你们一颗新心，在你们五内注入一种新的精神；我要从你们的肉体内除去那石心，并赐给你们一颗肉心。我要将我的神放在你们心中，使你们遵行我的法令，遵守并实行我的诫命。你们要安居在我赐给你们祖先的地方；你们要做我的子民，我要做你们的天主。我也要从一切不洁中拯救你们。}\BibleRef{则 36:23} 现在，这是新约的预言，不仅属于那唯一民族的遗民——正如别处所说：\BibleQuote{以色列子民的数量虽如海沙，但只有遗民得以得救，}\BibleRef{依 10:22}——而且也属于应许给他们的祖先和我们祖先的其他民族；这里应许了那重生的圣洗，正如我们所见，现在正赐给万民，对此，任何审视此事的人都不怀疑。而那位宗徒在赞美新约的恩宠及其与旧约相比的卓越时所说的话：\BibleQuote{你们是我们的书信……不是用墨水写的，而是用永生天主的圣神写的；不是写在石版上，而是写在心肉版上，}\BibleRef{格后 3:2} 显然指向先知所说的这句话：\BibleQuote{我还要赐给你们一颗新心，在你们五内注入一种新的精神；我要从你们的肉体内除去那石心，并赐给你们一颗肉心。}\BibleRef{则 38:26} 现在，宗徒所说的\Quote{心肉版}所取自的肉心，先知打算指出它与石心不同，因为拥有感觉生命；而感觉他理解为理性生命。因此，\OriginalTerm{属神的以色列}{spiritual Israel}不是由一个民族组成，而是由应许给祖先在祂的后裔中，即在基督内的所有民族组成。\par

49. 因此，\OriginalTerm{属神的以色列}{spiritual Israel}与由一个民族组成的\OriginalTerm{属血肉的以色列}{carnal Israel}的区别，在于恩宠的新颖，而非血统的高贵；在于感觉，而非种族；但先知在其深意中，在谈论属血肉的以色列时，未指明过渡便转而谈论属神的以色列，并且虽然现在谈论的是后者，却似乎仍在谈论前者；这并不是他吝啬我们清晰理解圣经，仿佛我们是敌人，而是他如同医生对待我们，给予我们有益的精神锻炼。因此，我们应该将这句话\BibleQuote{我要领你们回到你们的土地}以及他稍后如同重复自己所说的话\BibleQuote{你们要安居在我赐给你们祖先的地方}，不是按字面理解为指按血肉而生的以色列，而是按精神理解为指属神的以色列。因为那无玷无瑕、从万国中聚集、注定与基督永远为王的教会，本身就是有福之地、永生之地；我们应当理解，当这应许赐给祖先时，它就已经赐给了他们，因为祖先所相信的将在其时机赐予的东西，由于应许和旨意的不变性，对他们来说就如同已经赐予一样；正如宗徒在写给弟茂德的书信中所说的赐给圣人的恩宠：\BibleQuote{不是按照我们的行为，而是按照祂的旨意和恩宠，这恩宠是在万世以前，在基督耶稣内赐予我们的；如今借着我们救主的显现，已显明出来。}\BibleRef{弟后 1:9} 他称这恩宠为在那些将要领受的人尚不存在时就已赐下；因为他视之为在天主的安排和旨意中已经完成的事，这将在其时机发生，而他自己说它如今已显明出来。然而，这些话也可能指未来的时代之地，那时将有新天新地，不义之人不能居住。因此，确实对义人说，那地本身是他们的，其中没有任何部分属于不义之人；因为这就如同它已被赐予一样，当它被坚定地决定将要赐予时。\par

% structure: p0090 source=h2 target=subsection reason=relative-heading-rank; base=h2

\subsection{第三十五章 —— \Person{提科尼奥}的第五条规则}

50. 第五条规则——提科纽称之为\Quote{论时期}的规则——是我们可以藉此频繁地发现或推测圣经中未明确提及的时间量的规则。他说这条规则适用于两种方式：要么用于称为 \OriginalTerm{提喻}{synecdoche} 的修辞格，要么用于合法的数字。提喻这种修辞格要么以部分代整体，要么以整体代部分。例如，关于我们的主在山上仅仅在三位门徒面前显圣容，使祂的脸面明亮如太阳，衣裳洁白如雪的时间，一位福音记载此事发生在\Quote{八天以后}\BibleRef{路 9:28}，而另一位则记载发生在\Quote{六天以后}。现在，这两个关于天数的陈述不可能都正确，除非我们假定那位说\Quote{八天以后}的作者，将基督说预言那天的后半部分和祂成全预言那天的前半部分算作两个整天；而那位说\Quote{六天以后}的作者，只计算了这两天之间所有完整的那些日子。这种以部分代整体的修辞格，也解释了关于基督复活的那个大问题。因为除非我们把祂受难那天的后半部分和前一夜连接起来，算作一整天，又把祂复活的那一夜的后半部分与刚破晓的主日连接起来，也算作一整天，否则我们就无法得出祂预言自己将在地心里三天三夜的时间。\BibleRef{玛 12:40}\par

51. 其次，我们的作者称那些圣经特别推崇的数字为 \OriginalTerm{合法数字}{legitimate numbers}，例如七、十、十二，或任何细心研读圣经者很快会认识的其他数字。这类数字常常代表永恒的时间；例如，\Quote{我一日要赞美祢七次} 的意思与\Quote{祂的赞美常在我口中} 完全相同。它们的效力完全相同，无论是乘以十，如七十和七百（由此，耶肋米亚书中提到的七十年\BibleRef{耶 25:11} 可按属灵意义理解为教会旅居于外人中的整个时期），还是自乘，如十乘十得一百，十二乘十二得一百四十四，这最后一个数字在《默示录》中用来代表全体圣人。\BibleRef{默 7:4} 由此可见，这些数字不仅解决关于时间的问题，而且其意义适用范围更广，延伸至许多主题。例如《默示录》中上面提到的那个数字，不是指向时间，而是指向人。\par

% structure: p0093 source=h2 target=subsection reason=relative-heading-rank; base=h2

\subsection{第36章——提科纽的第六条规则。}

52. 第六条规则——提科纽称之为 \OriginalTerm{重述}{recapitulation}——在充分留心时可以在圣经中难解的地方发现。因为某些事件被如此叙述，以至于叙事似乎遵循时间顺序或事件的连续性，而实际上它未经说明就回到了先前在适当位置被遗漏的事件。如果我们不理解这一点，就会因为不应用这里所说的规则而犯错。例如，在《创世纪》中我们读到：\Quote{上主天主在东方的伊甸种植了一个乐园，并将祂所形成的人安置在里面。上主天主使地面生出各种好看而好吃的树。}\BibleRef{创 2:8} 这里似乎表示最后提到的事件发生在上主形成人并把他安置在乐园之后；而实际上，这两件事被简短提及，即上主种植了一个乐园，并将祂所形成的人安置在里面，然后叙事以重述的方式回头讲述之前被遗漏的事，即乐园是如何被种植的：上主使地面生出各种好看而好吃的树。接着有：\Quote{生命树也在乐园中央，还有知善恶树。} 接着是浇灌乐园的河流，它分为四支，是四条水源；所有这些都涉及乐园的布置。当这一切结束时，又重复了已经讲过的事实，但按事件的严格顺序是在这一切之后：\Quote{上主天主将人领来，安置在伊甸的乐园内。}\BibleRef{创 2:15} 因为这一切其他事完成后，人才被安置在乐园里，正如现在从叙事本身顺序所显示的那样：并非人安置在那里之后才做了其他事，正如先前的陈述可能暗示的那样，如果我们没有准确识别并理解叙事借此回溯先前被省略之事的重述。\par

53. 在同一部书中，当记述\Person{诺厄}儿子的世代时，这样说：\BibleQuote{这些是\Person{含}的儿子们，按照他们的家族、按照他们的语言、在他们的地方和他们的邦国。}\BibleRef{创 10:20} 又当列举\Person{闪}的儿子们时说：\BibleQuote{这些是\Person{闪}的儿子们，按照他们的家族、按照他们的语言、在他们的土地、按照他们的邦国。}\BibleRef{创 10:31} 并且论及他们全体时加上了：\BibleQuote{这些是\Person{诺厄}儿子们的家族，按照他们的世代、在他们的邦国；在这些当中，邦国在洪水之后在地上划分。全地都是一样的语言和一样的话语。} 现在加上这句话：\BibleQuote{全地都是一样的语言和一样的话语}，似乎表明在邦国分散到地上之时，他们全都拥有共同的一种语言；但这显然与前面的话\Quote{按照他们的家族，按照他们的语言}不相一致。因为如果全都有共同的一种语言，就不能说每个家族或邦国有自己的语言。所以，正是以\OriginalTerm{总括}{recapitulation}的方式加上了这句话：\BibleQuote{全地都是一样的语言和一样的话语}，这里的叙述在此回溯，未指明变化，讲述事情的经过：从拥有共同的一种语言，邦国被分为多种语言。因此，我们立即读到建塔的事，以及这惩罚因天主对傲慢的审判而降临在他们身上；正是在此之后，他们才按照自己的语言分散到全地。\par

54. 这种总括出现在更隐晦的形式中；例如，我们的主在福音中说：\BibleQuote{\Person{罗特}从\Place{索多玛}出去的那一天，火从天上降下，消灭了他们全体。人子显现的日子也要这样。在那一天，在屋顶上的人，他的器物在屋里，不要下来拿；在田里的人也同样不要回来。你们要记得\Person{罗特}的妻子。}\BibleRef{路 17:29} 难道是在我们的主显现之后，人才应当留心这些话语，而不回头看，即不渴望他们已弃绝的过去生活吗？岂不更应在现在留心它们，好叫当主显现时，每个人按照他所留心或轻视的事而得到赏报吗？然而，因为圣经说\Quote{在那一天}，人们就会认为留心这些话语的时间是主显现的时候，除非读者警醒而有智慧，以至理解总括（\OriginalTerm{总括}{recapitulation}），在这方面，另有经文帮助他，即使在宗徒时代就已宣告说：\BibleQuote{孩子们，现在是末时期了。}\BibleRef{若一 2:18} 因此，从福音被宣讲之时起，直到主显现之时为止，正是人们应当留心这些话语的那一天；因为主在这所说的显现，正属于同一天，而这一天将由审判之日终结。\BibleRef{罗 2:5}\par

% structure: p0097 source=h2 target=subsection reason=relative-heading-rank; base=h2

\subsection{第37章——\Person{蒂科尼乌斯}的第七条规则}

55. \Person{蒂科尼乌斯}的第七条规则，也是最后一条，是关于\Emph{魔鬼及其身体}的。因为他是恶者的首领，这些恶者在某种意义上就是他的身体，注定要与他一起进入永火的刑罚，正如基督是教会——祂的身体——的头，注定要与祂一起在永恒的国度与荣耀中。因此，正如被称为\Emph{主及其身体}的第一条规则指示我们，当圣经论及同一个位格时，要努力分辨陈述的哪一部分适用于头，哪一部分适用于身体；同样，这最后一条规则向我们显示，有些关于魔鬼的陈述，其真理性在关于他自身方面并不如关于他的身体方面那样明显；而他的身体不仅由那些明显走迷的人组成，也包括那些虽然实际上属于他，却暂时与教会混杂在一起的人，直到他们离开此世，或直到在最后的筛选中糠秕从麦子中被分离出来。例如，\Person{依撒意亚}说：\BibleQuote{怎么你从天上坠落了，早晨之子，光明者！}以及上下文在\Place{巴比伦}王的形象下同一位所说的话，当然都应理解为指魔鬼；然而同一处所说：\BibleQuote{他被压碎在地上，向所有民族派遣使者}，并不完全贴切地适用于头本身。因为，虽然魔鬼向所有民族派遣他的使者，然而被压碎在地上的是他的身体，而不是他自己，除非他本身存在于他的身体中，这身体被碾碎如地上的尘土，被风从地面吹起。\par

56. 以上所有这些规则，除了关于应许和法律的规则外，都使得在表达另一种意思时理解一种意思，这正是比喻措辞的特点；在我看来，这种措辞过于广泛，任何人都无法完全把握。因为，无论何时说出某件事，其意图是让另一件事被理解，我们就有了比喻表达，即使这个比喻名称在修辞学中找不到。当这种表达出现在通常可以找到的地方时，理解起来并不困难；但是，当它出现在不常见的地方时，理解起来就费力了，有些人费力多，有些人费力少，正如人们从天主那里获得智力恩赐的多少不同，或者他们获得的外部帮助多少不同。而且，正如我上面讨论的专用词的情况一样，事物被理解的方式与它们被表达的方式完全一致；同样，在比喻词的情况下，表达一件事而需要理解另一件事，我刚刚讲完了我认为足够的内容，这些令人敬畏的文献的学习者应该被劝告，不仅要熟悉圣经中常用的表达形式，仔细观察它们，准确记住它们，而且，特别和最重要的是，祈求天主使他们能够理解它们。因为在他们专注研究的这些书中，他们读到：\Quote{主赐予智慧：从祂口中发出知识和明智；}\BibleRef{箴 2:6}正是从祂那里，他们获得了对知识的渴望，如果这种渴望与虔诚相结合的话。但关于标记，就词语而言，我现在已经说得够了。剩下的，在下一本书中，按照天主赐给我的光照，讨论与他人交流思想的方法。\par

\SourceRecord{Translated by James Shaw.；Nicene and Post-Nicene Fathers, First Series；Vol. 2.；Edited by Philip Schaff.；Buffalo, NY: Christian Literature Publishing Co.,；1887.；<http://www.newadvent.org/fathers/12023.htm>.}

% structure: p0001 source=h1 target=section reason=page-title

\section{\WorkTitle{论基督圣道（第四卷）}}

% structure: p0002 source=h2 target=subsection reason=relative-heading-rank; base=h2

\subsection{第一章——本书并非修辞学论著}

1. 我这本名为\WorkTitle{论基督圣道}的著作，起初分为两部分。在序言中，我预先答复了可能反对这部著作的人，然后我说：\Quote{全部圣经的解释都取决于两件事：确定正确意义的方法，以及将已确定的意义传达出来的方法。我将先讨论确定意义的方法，其次讨论传达意义的方法。}因此，既然我已经就确定意义的方法讲了很多，并将三卷书用在这部分主题上，那么，我将只就传达意义的方法说几句话，以便尽可能把它们全部纳入一卷书中，从而以四卷书完成整部著作。\par

2. 因此，首先我愿藉此序言暂时平息那些读者可能有的期待——他们以为我将陈述我在世俗学校中学过并教过的修辞学规则，并警告他们说，不必从我这里寻求任何这类规则。并非我认为这些规则一无是处，而是说它们无论有何用处，都该从别处学习；倘若某位善人恰好有闲暇学习它们，他不可要求我在本著作或其他任何著作中教授它们。\par

% structure: p0005 source=h2 target=subsection reason=relative-heading-rank; base=h2

\subsection{第二章——基督徒教师使用修辞学艺术是合法的}

3. 修辞学这门艺术既可用于支持真理，也可用于支持谬误，谁敢说真理的捍卫者当赤手空拳地对抗谬误？例如，那些试图说服人接受虚假之事的人，知道如何引入话题，使听者处于友善、专注或受教的心态，而真理的捍卫者却不懂这门艺术？前者简短、清晰且可信地讲述他们的谎言，而后者却以令人厌倦、难以理解、总之不易相信的方式讲述真理？前者以诡辩的论据反对真理、保卫谬误，而后者却既不能捍卫真理，也不能反驳谬误？前者以错误的意见灌输听者的心灵，藉其言辞力量威吓、软化、激励、振奋他们，而后者在捍卫真理时却迟钝、冷漠、昏昏欲睡？谁如此愚蠢，竟认为这明智？既然雄辩的能力可用于双方，并在支持错误或正确方面大有助益，那么当恶人用它来赢得邪恶无价值的事业、助长不义和错误时，善人为何不研究如何让它为真理服务？\par

% structure: p0007 source=h2 target=subsection reason=relative-heading-rank; base=h2

\subsection{第三章——获得修辞技巧的合适年龄和合适方法}

4. 但关于这一主题的理论和规则（若再加上由练习和习惯而充分精于使用众多词汇和众多言辞装饰的舌辩，你便获得了所谓的\Emph{雄辩}或\Emph{演说术}）可以在我的这些著作之外学习，只要在适当年龄留出合适的时间。但只有那些能快速学会它们的人如此；因为罗马雄辩大师们自己也不避讳说，任何不能快速学会这门艺术的人，永远无法彻底掌握它。这是真是假，何须追问？因为即使这门艺术偶尔最终能被智力较慢的人掌握，我也不认为它重要到希望已成年的人花时间学习它。让少年人注意它就够了；即便他们中也不是所有将来要在教会中服务的人，而只是那些尚未从事更紧迫、或显然应优先于它的职业的人。因为智力和热情充沛的人，通过阅读和聆听雄辩演说家，比遵循雄辩规则更容易变得善辩。甚至在我们的大幸、被置于可靠权威地位的圣经正典之外，也不乏教会著作，有才能的人在阅读它们时，即使并不刻意追求，而只专注于所讨论的议题，也会沾染其中所写的雄辩色彩；当然，特别地，如果他再练习写作、口述，最后还讲论他基于虔敬和信德所形成的思想。然而，如果缺乏这种才能，修辞规则要么不被理解，要么在耗费巨大努力加以推行后，它们才被略微理解，也证明毫无用处。因为即使那些学会了规则并能流利优雅地讲话的人，在讲话时也不总是能想到这些规则以遵循它们，除非他们正在讨论规则本身。事实上，我认为几乎没有人能同时做到两件事——即，讲得好，并且为了做到这一点，在讲话时想着讲话的规则。因为我们必须小心，在我们思考如何按照艺术规则讲话时，不要让要说的话溜走。然而，在雄辩者的演说中，我们发现了被遵循的雄辩规则，而演讲者在讲话时并没有将它们视为雄辩的辅助，无论他们曾经学过规则，还是从未遇到过它们。因为他们之所以成为雄辩者，才体现了这些规则；并非他们使用这些规则来成为雄辩者。\par

5. 因此，既然婴孩唯有从会说话的人那里学习词语和短语才能学说话，那么人为何不能仅仅通过阅读和学习雄辩者的演说并尽力模仿他们，而未受任何言辞艺术的教育，就成为雄辩者呢？从实际的例子中，我们在这方面看到了什么？我们知道许多人，虽不熟悉修辞规则，却比许多学过这些规则的人更雄辩；但我们不知道有哪个人是没有读过和听过雄辩者的演说和辩论而成为雄辩者的。因为即使文法的艺术教导言语的正确性，男童也不必学习，如果他们有在说正确言语的人中间成长和生活的优势。因为即使不知道任何过错的名称，他们从习惯于正确言语，就会抓住所听到的任何言语中的错误并加以避免；正如城里人，即使不识字，也能抓住乡野之人的错误。\par

% structure: p0010 source=h2 target=subsection reason=relative-heading-rank; base=h2

\subsection{第四章 基督徒教师的职责}

6. 因此，圣经的解释者和教师、真信仰的捍卫者和谬误的反对者，其职责是既教导正确的事，又驳斥错误的事，并在执行此任务时，安抚敌意者，激励疏忽者，告知无知者现今所发生的事以及未来可能发生的事。然而，一旦听者变得友好、专注且乐于学习，无论是他原本如此，还是教师使他们如此，其余的目标便视情况所需而进行。如果听者需要教导，就必须通过叙述使所讨论的事完全明了。另一方面，澄清可疑之点则需要推理和证明的展示。然而，如果听者需要的是激励而非教导，以便他们勤奋地去做已经知道的事，并使他们的情感与他们所承认的真理和谐一致，那就需要更有力的言辞。在此，恳求与责备、劝诫与斥责，以及所有其他激发情感的方法，都是必要的。\par

7. 我所提及的一切方法，在言语作为媒介的场合，几乎人人常用。\par

% structure: p0013 source=h2 target=subsection reason=relative-heading-rank; base=h2

\subsection{第五章 智慧比口才更为重要——对基督徒教师而言}

然而，有些人运用这些方法时粗俗、不雅、冰冷，而另一些人则运用得敏锐、优雅、有气势，那么我所说的这项工作应该由那些能够以智慧——即使没有口才——进行辩论和演讲的人来承担，并使听者受益，即使他受益的程度不如他同时拥有口才时那么大。但是我们必须提防那充斥口才而无内容的人，尤其是当听者喜欢那些不值得听的内容，并认为因为演说者口才好，所以他说的必定是真的时候。这种看法甚至也被那些认为修辞艺术应该教授的人所持有；因为他们承认：\Quote{智慧虽无口才，对国家少有助益；然而口才而无智慧，往往造成实际损害，从无益处。}如果教导口才原理的人，即使在论述口才的著作中，也被真理所迫而承认这一点，尽管他们不知道真正的智慧——即从光明之父那里降下的天上智慧——那么，我们这些更高智慧的儿子和仆人，岂不更应感受到这一点！现在，一个人的讲话智慧或多或少，正好比他在圣经知识上的进步；我不是指他大量阅读和记忆，而是指正确理解并仔细探究其含义。因为有人阅读却忽略；他们记住词语，却不关心理解含义。显然，我们必须将那些不善于记忆词语、却能以心之眼洞察圣经之心的人，远远置于这些人之上。然而，比这两种人都更好的是那种人：当他愿意时，既能重复词语，同时又能正确领会其含义。\par

8. 现在，对于那些必须以智慧讲话，即使不能以口才讲话的人，特别需要记住圣经的话。因为他越看出自己言语的贫乏，就越应汲取圣经的丰富，以便他用自己的话所说的，能用圣经的话来证明；而他自己，虽在自己的话中渺小软弱，却能从伟人的确证见证中获得力量和权能。因为他的证明能带来喜悦，当他不能用说话风格取悦人时。但是，如果有人不仅渴望以智慧讲话，还渴望以口才讲话（若他能两者兼具，肯定更有助益），我宁愿送他去阅读、聆听和练习模仿雄辩者，而不建议他与修辞教师共度时间；尤其当他所读所听的人，被公正地称赞为不仅以口才，而且以智慧讲话，或习惯于如此讲话时。因为口才雄辩者听来令人愉悦，智慧者听来令人受益。因此，圣经不说群众有口才，而是说\BibleQuote{智慧群众乃是世界的幸福。}\BibleRef{智 6:24}正如我们必须常常吞下有益健康的苦药，同样我们必须始终避免有害健康的甜食。但是，有什么比有益健康的甜味或甜蜜的健康更有益呢？因为我们越是努力使这些东西变甜，就越容易使它们的健康效用发挥作用。因此，教会的著作家们，不仅以智慧，也以口才诠释了圣经；而阅读这些著作的时间，并不超过那些勤勉且有闲暇的人足以穷尽它们的时间。\par

% structure: p0016 source=h2 target=subsection reason=relative-heading-rank; base=h2

\subsection{第六章 圣经作者将口才与智慧相结合}

这里，可能有人会问，那些其受天主默感的著作构成正典的作者们——该正典具有最有益的权威——是只应被视为有智慧，还是也应被视为有雄辩呢？这个问题对我来说，以及对于与我持相同看法的人，是很容易解决的。因为在我理解这些作者的地方，我似乎不仅觉得没有什么比他们更有智慧，而且觉得没有什么比他们更雄辩。我敢断言，所有真正理解这些作者所言的人，同时也会领悟到，这些话不可能以任何其他方式被恰当地说出。因为正如有一种雄辩更适合青年，有一种更适于老年，若不适合说话者的身份，就不能称为雄辩，同样，有一种雄辩适合于那些正义地要求最高权威、且明显受天主默感的人。他们就是用这种雄辩说话；没有其他雄辩适合他们；而这种雄辩本身也不适合其他任何人，因为它符合他们的性格，同时（不是由于空洞的浮夸，而是由于实在的价值）它超越他人的雄辩，正如它看似低于他们一样。然而，在我不能理解这些作者的地方，尽管他们的雄辩不那么明显，我毫不怀疑那雄辩与我理解的部分属于同一种类。这些神圣有益的话语的晦涩本身，也是某种雄辩的必要元素，这种雄辩旨在助益我们的理解力，不仅通过发现真理，也通过操练其能力。\par

然而，我若有时间，可以向那些赞扬自己语言形式优于我们作者的人（不是因为其崇高，而是因其浮夸）指出，他们引以为傲的所有那些雄辩的力量与美，都能在天主出于其美善而提供给我们的圣经中找到，以塑造我们的品格，并引导我们从这邪恶的世界走向上方的福乐世界。但是，并非这些作者与外教演说家和诗人共有的特质，使我在他们的雄辩中得到难以言喻的喜悦；我更惊叹的是，凭借他们特有的雄辩，他们如此运用我们的这种雄辩，使其不论在出现或缺席时都不显眼：因为谴责它或炫耀它都不适合他们；若他们避开它，便是做了前者；若他们突出它，则可能显得在行后者。而在博学之士注意到其出现的那些段落中，所谈论的题材都是如此，以致于表达它们的词语似乎不是说话者刻意寻找的，而是自发呈现出来的；好像智慧走出自己的房屋——即智者的胸怀——而雄辩则如同不可分离的侍者，不请自随。\par

% structure: p0019 source=h2 target=subsection reason=relative-heading-rank; base=h2

\subsection{第七章——从保禄书信与亚毛斯预言中汲取的真雄辩范例。}

因为谁会看不见宗徒在以下经文中想说什么，以及他说得多么智慧呢：\BibleQuote{我们也因患难而夸耀，因为知道患难生忍耐，忍耐生老练，老练生盼望，盼望不至于羞耻，因为天主的爱借着赐予我们的圣神，已倾注在我们心中。}\BibleRef{罗 5:3} 如今，若有人以不学无术的博学（若我可这样说）来争论说宗徒在此遵循了修辞规则，岂不每个基督徒，无论有学无学，都会嘲笑他？然而我们在此发现了那个在希腊文中称为\OriginalTerm{层进法}{\GreekText{κλίμαζ} (climax)}、拉丁文中有些人称之为\OriginalTerm{递进}{gradatio}（因为他们不愿称之为\OriginalTerm{梯状}{scala (a ladder)}）的修辞格，即词语和思想具有相互依存的联系，正如我们在此看到忍耐出自患难，老练出自忍耐，盼望出自老练。这里还发现另一种修饰：因为在某些以单一语调完成的陈述（我们称之为\OriginalTerm{分句与断句}{membra et cæsa}，希腊人称为\OriginalTerm{从句}{\GreekText{κῶλα}}和\OriginalTerm{片语}{\GreekText{κόμματα}}）之后，紧接着一个圆周句（\OriginalTerm{圆周句}{ambitus sive circuitus}），希腊人称为\OriginalTerm{圆周句}{\GreekText{περίοδος}}，其分句在说话者声调中悬置，直到整个句子由最后一个分句完成。因为那些先于圆周句的陈述，第一个分句是：\Quote{知道患难生忍耐；}第二个：\Quote{忍耐生老练；}第三个：\Quote{老练生盼望。}然后随后附属的圆周句由三个分句完成，其中第一个是：\Quote{盼望不至于羞耻；}第二个：\Quote{因为天主的爱已倾注在我们心中；}第三个：\Quote{由赐予我们的圣神。}但这些和同类的其他事物都包含在口才艺术的教学中。因此，正如我不肯定宗徒受雄辩规则指导，我亦不否认他的智慧自然地产生并伴有雄辩。\par

12. 在致格林多人的第二封书信中，他又驳斥了那些从犹太人中间出来、企图诋毁他名声的假宗徒。他被迫谈论自己，尽管将此视为愚妄，但他的言谈何其智慧、何其雄辩！智慧是他的向导，雄辩是他的侍从；他追随前者，后者追随他，但他并不因后者追随而轻视它。他说：\BibleQuote{我再言，不要有人以为我是愚妄的；即使如此，也请当作愚妄人接纳我，好让我稍作自夸。我所说的话，不是按主说的，而是如同愚妄人，在这自夸的凭据中说的。既然有许多人凭着血肉夸耀，我也要夸耀。因为你们很乐意容忍愚妄人，既然你们自己是明智的。因为，若有人奴役你们，若有人吞噬你们，若有人夺取你们，若有人自高自大，若有人打你们的脸，你们都能容忍。我讲这话，是出于羞辱，似乎我们软弱。然而，凡在什么事上别人胆大（我说句愚妄话），我也胆大。他们是希伯来人吗？我也是。他们是以色列人吗？我也是。他们是亚巴郎的后裔吗？我也是。他们是基督的仆役吗？（我说句愚妄话）我更是：多受劳苦，多受鞭打，多下监牢，常冒死亡。被犹太人鞭打五次，每次四十下减去一下；被棍打了三次，被石头砸了一次；遇船坏三次，在深海一昼一夜；多次行路，遭江河的危险、盗贼的危险、同族的危险、外邦人的危险、城中的危险、旷野的危险、海上的危险、假弟兄的危险；劳碌辛苦，屡次不得睡；忍饥受渴，屡次不得食；受冷受赤。除了外面的事，还有每日压在我身上的，就是对众教会的挂虑。有谁软弱，我不软弱呢？有谁跌倒，我不焦急呢？若必需夸耀，我就夸耀我软弱的事。}\BibleRef{格后 11:16}深思熟虑的人必能看出这些话中有多少智慧。连沉沉入睡的人也必注意到其中蕴含的何等雄辩之流。\par

13. 更进一步，有学识的人会注意到，那些希腊人称为\OriginalTerm{断句}{\GreekText{κόμματα}}的部分，以及我刚才提到的分句和句段，以最美丽的变换相互交织，构成了那语辞的整体形貌与特征（可以这么说），连未受过教育的人也为之欣喜感动。因为，从我开始引述的地方起，这段话由若干句段组成：第一个句段尽可能短，由两个成员组成（一个句段不能少于两个成员，但可以有更多）：\BibleQuote{我再言，不要有人以为我是愚妄的。}第二个有三个成员：\BibleQuote{即使如此，也请当作愚妄人接纳我，好让我稍作自夸。}第三个有四个成员：\BibleQuote{我所说的话，不是按主说的，而是如同愚妄人，在这自夸的凭据中说的。}第四个有两个：\BibleQuote{既然有许多人凭着血肉夸耀，我也要夸耀。}第五个有两个：\BibleQuote{因为你们很乐意容忍愚妄人，既然你们自己是明智的。}第六个又由两个成员组成：\BibleQuote{因为，若有人奴役你们。}接着是三个节段（\OriginalTerm{节段}{cæsa}）：\BibleQuote{若有人吞噬你们，若有人夺取你们，若有人自高自大。}然后是三个分句（\OriginalTerm{分句}{membra}）：\BibleQuote{有人打你们的脸。我讲这话，是出于羞辱，似乎我们软弱。}随后接上一个由三个成员组成的句段：\BibleQuote{然而，凡在什么事上别人胆大（我说句愚妄话），我也胆大。}此后，以疑问形式提出若干单独的节段，也以单独的节段给予回答，三对三：\BibleQuote{他们是希伯来人吗？我也是。他们是以色列人吗？我也是。他们是亚巴郎的后裔吗？我也是。}但第四个节段同样以疑问形式提出，回答却不是以另一个节段（\OriginalTerm{节段}{cæsum}）而是以一个分句（\OriginalTerm{分句}{membrum}）给出：\BibleQuote{他们是基督的仆役吗？（我说句愚妄话。）我更是。}接着连续给出四个节段，疑问形式被极为文雅地省略了：\BibleQuote{多受劳苦，多受鞭打，多下监牢，常冒死亡。}然后插入一个短句段；因为，通过停顿，\BibleQuote{被犹太人五次}应作为一个成员标出，第二个成员接上：\BibleQuote{鞭打，每次四十下减去一下。}然后他回到节段，列出三个：\BibleQuote{被棍打了三次，被石头砸了一次，遇船坏三次。}接着是一个分句：\BibleQuote{在深海一昼一夜。}接着十四个节段以最恰当的猛烈迸发出来：\BibleQuote{多次行路，遭江河的危险、盗贼的危险、同族的危险、外邦人的危险、城中的危险、旷野的危险、海上的危险、假弟兄的危险；劳碌辛苦，屡次不得睡；忍饥受渴，屡次不得食；受冷受赤。}之后是一个由三个成员组成的句段：\BibleQuote{除了外面的事，还有每日压在我身上的，就是对众教会的挂虑。}接着他以询问的语气加上两个分句：\BibleQuote{有谁软弱，我不软弱呢？有谁跌倒，我不焦急呢？}最后，整段话仿佛在喘息中，以一个由两个成员组成的句段结束：\BibleQuote{若必需夸耀，我就夸耀我软弱的事。}我无法充分表达，在这段迸发之后，他插入一段简短的叙述，让自己休息，也让听者休息，这是多么美好而令人愉悦。因为他接着说：\BibleQuote{我们的主耶稣基督的天主和父，是应当永远受赞美的，他知道我没有说谎。}然后他非常简短地讲述了他所遭遇的危险以及逃脱的方式。\par

14. 若继续深究，或在圣经其他段落中指出同样的事实，未免冗长。假设我费了额外的工夫，至少就我所引的宗徒著作中的段落，指出修辞学所教导的修辞格，那又如何呢？岂不是更可能让严肃的人认为我做得太过，而好学者中却没有人认为我已足够？凡由师长传授的这些事，都被视为极有价值，人们付出高价购买，贩卖者对其大肆吹嘘。我担心，在我这样讨论这类事情时，也不免沾染那种吹嘘之气。然而，有必要回答那些学养不足的人，他们轻视我们的作者，并非因为作者不具备，而是因为他们没有展示出这些人极为看重的雄辩。\par

15. 但或许有人会想，我之所以选择\Person{宗徒保禄}，是因其为我们伟大的演说家。因为当他说：\Quote{我虽然在言语上粗鲁，但在知识上却不如此}\BibleRef{格后 11:6}，他似乎是向诽谤者让步这样说，并非承认他认可其真实性。倘若他曾说：\Quote{我确实在言语上粗鲁，但在知识上却不如此}，我们便绝不能另作解释。他毫不犹豫地明确宣告自己的知识，因为没有知识，他就不能做外邦人的教师。当然，如果我们提出他的某篇著述作为雄辩的模范，我们取自那些书信，就连那些认为他身体软弱、言语可鄙的诽谤者，也承认这些信是沉重有力的。\BibleRef{格后 10:10}\par

于是，我看必须也谈谈先知们的雄辩，其中许多事物被隐藏在隐喻的风格之下，它们越是完全埋没在修辞格中，一旦显露出来，就越令人愉悦。然而在此处，我的职责是选择这样一段经文，使我无需解释内容，而只需称赞其文体。我将这样做，主要引用那位自称是牧羊人或牧人的先知的书，他被天主从那职业中召叫，并奉命向天主的人民说预言。然而，我不愿遵循七十子译本，他们在翻译中虽受圣神引导，但似乎为引导读者更专注于灵性意义的探究而改动了一些段落（因此他们译本的某些段落因更加具象征性而更加晦涩）；但我将遵循由司铎\Person{热罗尼莫}从希伯来文译为拉丁文的译本，他对两种语言都非常熟悉。\par

16. 那么，当这位乡野的，或说\Emph{quondam}乡野的先知，责骂那些不虔敬的、骄傲的、奢侈的，因而最忽视兄弟之爱的人时，他大声呼喊说：\Quote{祸哉，你们在\Place{熙雍}安逸，在\Place{撒玛黎雅}山上自恃，你们是人民的首领和元老，以盛仪进入\Place{以色列}家！你们往\Place{加耳乃}去看看吧，从那里往\Place{大哈玛特}去，然后下到\Place{培肋舍特人的加特}，以及这些地方最好的王国：难道他们的疆界比你们的疆界还大吗？你们这些被分别出来为灾祸之日，接近暴虐之座的人；你们躺在象牙床上，伸展在卧榻上，吃羊群中的羔羊，和牛群中的肥犊；你们弹琴歌唱，自以为有如\Person{达味}的乐器；用碗喝酒，用最贵重的油抹身，却不为\Person{若瑟}的苦难而悲伤。}假如那些自以为是博学善辩的人，轻视我们的先知为不学无术、言语拙劣的人，他们不得不传递这样一种信息，并且是向这样的人传递，他们是否会选择在任何方面以不同的方式表达自己？——至少那些不愿像疯子般咆哮的人。\par

17. 因为有什么是清醒的耳朵希望在这段讲话中改变的呢？首先是痛斥本身：它以何等猛烈之势扑向昏沉的感官，使它们警醒：\Quote{祸哉，你们在\Place{熙雍}安逸，在\Place{撒玛黎雅}山上自恃，你们是人民的首领和元老，以盛仪进入\Place{以色列}家！}接着，为了利用天主赐予他们广大领土的恩惠，以显示他们崇拜偶像而依靠\Place{撒玛黎雅}山的忘恩负义，他说：\Quote{你们往\Place{加耳乃}去，}他说，\Quote{看看；从那里往\Place{大哈玛特}去，然后下到\Place{培肋舍特人的加特}，以及这些地方最好的王国：难道他们的疆界比你们的疆界还大吗？}同时，在说这些话时，文体被地名如明灯般装饰，例如\Quote{熙雍，}\Quote{撒玛黎雅，}\Quote{加耳乃，}\Quote{大哈玛特，}和\Quote{培肋舍特人的加特。}然后与这些地方相连的词语也极为恰当地变化：\Quote{你们安逸，}\Quote{你们自恃，}\Quote{前往，}\Quote{去，}\Quote{下来。}\par

18. 接着预告了将来在残暴君王下的被掳，加上一句：\Quote{你们这些为邪恶之日所分别出来、又亲近欺压之座的人。}然后接续奢靡之恶：\Quote{你们躺卧在象牙床上，伸展在榻上；吃羊群中的羔羊，和牛群中的牛犊。}这六句话组成三个对句，每对两个分句。因为他没有说：\Quote{你们这些为邪恶之日所分别出来的人，你们这些亲近欺压之座的人，你们这些睡在象牙床上的人，你们这些伸展在榻上的人，你们这些吃羊群中羔羊的人和牛群中牛犊的人。}如果这样表达，也有其美感：六个独立分句连续，同一代词每次重复，每个分句由说话者一口气完成。但这样更美：分句在同一代词下成对结合，形成三个句子，一句关于被掳的预言：\Quote{你们这些为邪恶之日所分别出来、又亲近欺压之座的人}；第二句关于放荡：\Quote{你们躺卧在象牙床上，伸展在榻上}；第三句关于贪食：\Quote{你们吃羊群中的羔羊，和牛群中的牛犊。}因此，说话者可以自行判断：是每个分句单独结束，共计六个；还是在第一、第三、第五分句处停顿，将第二与第一、第四与第三、第六与第五结合，形成三个最优雅的对句，每对两个分句：一个描述迫近的灾祸；另一个描述放荡的床榻；第三个描述奢侈的餐桌。\par

19. 接着，他斥责他们在听觉上追求享乐的奢靡。在此，当他说\Quote{你们弹奏琴瑟而歌唱}时，考虑到智慧人可以智慧地运用音乐，他以极妙的演说技巧控制了斥责的流势；现在他不是对这些人说，而是论及这些人，为向我们表明必须分辨智慧之乐与纵欲之乐，他没有说：\Quote{你们弹奏琴瑟而歌唱，并自以为有如\Person{达味}的乐器}；而是先对他们说出了纵欲者应当听到的话：\Quote{你们弹奏琴瑟而歌唱}；然后，转向别人，暗示这些人连自己的技艺都未曾掌握：\Quote{他们自以为有如\Person{达味}的乐器；用碗饮酒，用最贵重的油膏抹自己。}这三个分句最好这样念：在周期前两个分句处停顿，在第三个分句处结束。\par

20. 现在来看紧随这些之后的句子：\Quote{他们不为\Person{若瑟}的苦难而忧伤。}无论这是一个分句连续念出，还是更优雅地先停顿\Quote{他们不为…忧伤}，然后加上\Quote{…\Person{若瑟}的苦难}，从而形成两个分句的周期；无论如何，这是一种奇妙的优美，没有说\Quote{他们不为他们兄弟的苦难而忧伤}，而是用\Person{若瑟}代替兄弟，以这位在弟兄中因所受伤害和所回报善行而卓著之人的专名来指代一般兄弟。事实上，我不知道这种以\Person{若瑟}代替一般兄弟的修辞格是否属于我学过并曾教授的那门修辞艺术。但它是多么优美，多么打动有智识的读者，无需告诉任何本身不感受到的人。\par

21. 在我所选为例的这段经文中，还可以找到许多其他涉及修辞法则的要点。但有智识的读者与其通过仔细分析而受教，不如通过富有精神地背诵而被点燃。这段经文并非出于人的技巧和关注，而是从天主的智慧与口才中涌流而出；智慧不追求口才，口才也不回避智慧。因为，正如某些极具口才和敏锐的人所察觉并指出的那样，口才艺术中确立的规则，若非先在演说家的天才中诞生，就不能被观察、标记并系统化，那么，这些规则出现在那位一切天才之源的使者身上，又有什么可惊奇的呢？因此，让我们承认，正典作者不仅有智慧，而且有口才——一种与其身份和地位相称的口才。\par

% structure: p0032 source=h2 target=subsection reason=relative-heading-rank; base=h2

\subsection{第8章——圣作者之隐晦虽与口才相容，但基督徒教师不应模仿。}

22. 但虽然我从那些容易理解的经文中取一些雄辩的例子，我们绝不应当认为我们有责任模仿他们在某些段落中的表现，在这些段落中，为了锻炼和训练读者的心智，打破那些愿意学习之人的厌倦并激发他们的热忱，也为给不虔敬者的心智蒙上一层帷幕，使他们要么被转化为虔诚，要么被排除在奥秘知识之外，出于这些原因之一，他们以一种有益而健康的隐晦来表达自己。事实上，他们以这样的方式表达自己，以致后来正确地理解并解释他们的那些人，在天主的教会中获得了尊敬，虽然不与他们自己所受的尊敬相等，但仅次于它。因此，这些作者的诠释者不应当以同样的方式表达自己，仿佛将他们的诠释作为同等权威提出；而应当在所有的讲论中，以被理解为首要和主要目标，尽可能使用如此清晰的言辞，以至于要么不明白他们的人非常迟钝，要么如果他们所说的不容易或迅速被理解，原因不在于他们的表达方式，而在于他们试图解释的问题的困难和微妙之处。\par

% structure: p0034 source=h2 target=subsection reason=relative-heading-rank; base=h2

\subsection{第九章 如何以及与何人讨论难解的段落}

23. 因为有些段落，无论讲者以多长的篇幅、多么清晰或多么雄辩地阐述，都无法以其适当的力量被理解，或者很难被理解；这些段落根本不应该带到公众面前，或者只在有紧急原因时才偶尔为之。然而，在那些以这样一种风格写成的书籍中，如果被理解，它们可以说吸引读者，如果不被理解，它们也不会给那些不愿阅读的人带来麻烦；在私人谈话中，我们不应回避将我们自己已达到的真理置于他人理解之内的责任，无论它多么难以理解，也无论我们在论证中付出多少劳动。只有两个条件必须坚持：我们的听众或同伴应当有真诚的学习真理的愿望，并且有智力以任何形式接受所传达的真理，教师不应当那么关心雄辩，而更关心教学的清晰性。\par

% structure: p0036 source=h2 target=subsection reason=relative-heading-rank; base=h2

\subsection{第十章 风格明晰的必要性}

24. 对清晰性的强烈渴望有时会导致忽视更优雅的言辞形式，并且不关心听起来如何，而更关心清楚地表达和传达意图的含义。因此，某位作者在处理这种言辞时，说其中有一种\Quote{一种精心的疏忽}。然而，在去除装饰的同时，它并没有带来粗俗的言辞；尽管好教师有或应该有如此大的教学热情，以至于他们会使用一个不能变成纯正拉丁文而不变得含糊或歧义的词，但按照俗语用法使用时既不歧义也不含糊，不是按照有学问的人的方式，而是按照无学问的人使用它的方式。因为如果我们的译者不忌讳说\Quote{\Emph{Non congregabo conventicula eorum de sanguinibus},} 因为他们觉得为了意义而在此处使用一个在拉丁文中只用单数的复数词是重要的，那么为什么一个有虔诚之心的教师，在向无学问的听众说话时，会忌讳使用 \Emph{ossum} 而不是 \Emph{os}，如果他担心后者可能被理解为不是 \Emph{ossa} 的单数，而是 \Emph{ora} 的单数，因为非洲人的耳朵对元音的长短没有敏锐的感知？而纯粹的言语如果不能使听者理解，又有什么好处呢？因为如果我们说话的对象不理解我们，说话就完全没有用处。因此，教导的人将避免所有不能教导的词语；如果他能找到既纯粹又易懂的词，他优先选用；如果不能，要么因为没有这样的词，要么因为当时没想到，他将使用不太纯粹的词，只要他的思想实质被完整传达和领会。\par

25. 这一点必须坚持为使我们被理解所必需，不仅在与人谈话时——无论对象是一人还是多人——而且在公开演讲时更是如此：因为在谈话中，任何人都可以提问；但当众人都静默以倾听一人发言，且所有面孔都专注地转向他时，询问自己不理解的问题既不合惯例也不得体；因此，讲者应特别留意帮助那些不能提问的人。渴望受教的群众通常会通过他们的动作表明是否理解所讲的内容；除非有此类迹象，否则所讨论的主题应反复翻转，以各种形式、样式和措辞表达——这一点是那些背诵预先准备好的、已经记熟的话语的人无法做到的。然而，一旦讲者确定他所讲的内容已被理解，他应当结束演讲或转向另一个要点。因为若一个人在人们希望受教的要点上提供光照，他给予喜乐；但若他在已经熟知的事物上长篇大论，他就变得令人厌倦，尤其是当人们的期望原本是解除经文的难点时。即便是非常熟悉的事物，如果注意力不是转向事物本身，而是转向讲述它们的方式，那么叙述它们也会因所给予的喜乐而被讲述。甚至当风格本身已为人熟知，如果它对听众有吸引力，那么讲话者是讲者还是朗读者几乎无关紧要。因为文笔优美的东西，不仅初次接触的人读来愉快，已接触并尚未忘记的人重读也愉快；甚至这两类人听别人复述也会感到愉悦。若有人遗忘了什么，在被提醒时他就受了教导。但我现在不是在处理给予喜乐的方式。我在谈论那些渴望学习的人应当被教导的方式。最好的方式就是确保听者听到真理，并且他所听到的能被理解。当这一点达到后，就无需再在真理本身上下工夫，仿佛需要进一步解释；但或许可以稍费力气将其强化，以使之深入人心。如果这样做是适当的，那么应当适度进行，以免导致厌倦和不耐烦。\par

% structure: p0039 source=h2 target=subsection reason=relative-heading-rank; base=h2

\subsection{第十一章 基督徒教师必须清楚演讲，但不可不讲文采。}

26. 因为教导当然包括真正的口才，不在于让人们喜欢他们不喜欢的，也不在于让他们做他们退缩的，而在于使晦涩变得清晰；然而如果这一点没有文采地完成，益处只限于少数热切的学生——他们渴望了解任何需学习的内容，无论其形式多么粗陋不雅；当他们成功获得目标时，他们发现朴素的真理已是足够可口的食物。善于思考的心灵的一个显著特征就是不爱词语，而爱词语中的真理。金钥匙若不能打开我们想打开的东西，有何用处？木钥匙若能——因为打开关闭的东西正是我们想要的——又有何不妥？但由于学习与饮食之间存在某种类比，没有它就无法生存的食物本身必须调味以适应多数人的口味。\par

% structure: p0041 source=h2 target=subsection reason=relative-heading-rank; base=h2

\subsection{第十二章 按西塞罗所言，演说家的目的是教导、悦乐与感动。其中，教导最为根本。}

27. 因此，一位伟大的演说家曾说：\Quote{有口才的人必须讲得能教导、悦乐与说服。}然后他补充道：\Quote{教导是必要的，悦乐是优美的，说服是胜利。}在这三者中，首先提到的，教导——这是必要的事——取决于我们说什么；其他两者取决于我们如何说。因此，以教导为目的说话的人，不应认为只要自己未被理解就已说出了他要说的；因为虽然他所讲的他自己明白，但对于不理解的人而言，他等于没说。然而，若他被理解了，那么无论他讲话的方式如何，他已经说出了他的意思。但若他还想悦乐或说服听众，他不会通过以任何形状呈现他的思想来实现这一目的，而需要讲话的风格。并且，正如听众必须被取悦以赢得其注意力，同样也必须被说服以促其行动。听众若因你甜美优雅的讲话而被取悦，那么若他因你的许诺而被吸引、因你的威胁而畏惧；若他拒绝你所谴责的，拥抱你所称赞的；若你堆积悲伤的对象时他悲伤，你指出喜乐的对象时他喜乐；若你呈现在他面前作为怜悯对象的人他怜悯，你放在他面前作为可惧可避的人他退缩——他就被说服了。我不必罗列所有其他强有力的口才可以用于感动听众心灵的事，不是告诉他们应该做什么，而是敦促他们去做已经知道应该做的事。\par

28. 然而，如果他们还不知道这一点，当然必须先受教导，然后才能被感动。也许仅仅知道自己的责任就足以产生如此效果，以至于无需用更强大的口才去感动他们。然而，当需要这样做时，就应该去做。当人们知道该做什么却不去做时，就需要这样做。因此，教导是必要的。因为人们知道的事，他们自己有权去做或不做。但谁能说，去做他们不知道的事是他们的责任呢？按照同样的原则，说服不是必要的：因为它并非总是被需要；例如，当听者仅仅因教导或取悦而表示同意时。正因如此，说服也是一种胜利，因为一个人可能被教导并被取悦，却仍然不同意。如果我们未达成第三个目标，获得前两个目标又有何用？取悦也非必要；因为在演说过程中，真理被清晰地指出（这正是教导的真正功能），事实并非如此，意图也非如此，即言语的风格应使真理令人愉悦，或风格本身应带来愉悦；而是真理本身，当以其纯朴的方式展现时，就带来愉悦，因为它是真理。因此，即使是虚假的事物，当被揭示和暴露时，也常常成为愉悦的来源。当然，不是它们的虚假带来愉悦；而是因为它们是虚假的这一事实是真实的，揭示这一事实的言语带来愉悦。\par

% structure: p0044 source=h2 target=subsection reason=relative-heading-rank; base=h2

\subsection{第13章 听者必须既受教导又被感动}

29. 但为了那些极其挑剔的人——他们只在意真理是否以令人愉悦的演说形式呈现——在口才中取悦的艺术被赋予了不小的地位。然而，即使如此，对于那些冥顽不灵的人来说，这也不够，他们既理解也喜欢教师的言论，却没有从中获得任何益处。因为一个人既承认真理又赞美口才，却不同意，这对他有何益处呢？讲者之所以在敦促真理时仔细注意他所说的，正是为了赢得他的同意。如果所教导的真理仅仅是相信或知道就足够了，那么表示同意无非就是承认它们是真实的。然而，当所教导的真理是必须付诸实践的，并且教导的目的正是为了实践，那么，如果未被学习以致实践，被说服所说的是真理是无用的，对说的方式感到满意也是无用的。因此，雄辩的讲道者，在敦促实践真理时，不仅要教导以提供指导，取悦以保持注意力，还必须打动心灵以征服意志。因为如果一个人没有被真理的力量所感动——尽管真理通过他自己的承认得到证明，并以优美的风格呈现——那么剩下的只有通过口才的力量征服他。\par

% structure: p0046 source=h2 target=subsection reason=relative-heading-rank; base=h2

\subsection{第14章 言辞之美须与内容相称}

30. 人们在此所谈到的表达之美上花费了如此多的精力，以至于我们不仅不应去做，反而应当躲避和憎恶许多邪恶和卑劣的行径——这些行径被邪恶卑劣之人以极大的口才推荐，并非为了获得同意，而仅仅为了被愉悦地阅读。但愿天主使祂的教会远离\Person{耶肋米亚}先知论及犹太人的会堂时所说的话：\BibleQuote{地上发生了一件令人惊奇而可怕的事：先知们假预言，司祭们用手鼓掌，我的百姓喜爱这样；到了结局，你们将做什么呢？}哦，口才，因其纯洁而更显可畏，因其坚实而更具摧毁力！它确实是\BibleQuote{打碎磐石的铁锤}。因为天主自己曾通过同一位先知将祂通过祂的圣先知所说的话比作此物（\BibleRef{耶 23:29}）。但愿不会，但愿我们中间没有司祭为假先知鼓掌，没有天主的子民喜爱这样。我说，但愿在我们中间没有如此可怕的疯狂！因为到结局我们将做什么呢？确实，即使所说的话不太明白、不那么令人愉悦、不那么有说服力，但讲真话，并且正义而非邪恶被喜悦地聆听，这更可取。当然，这不可能实现，除非真理和正义被优美地表达。\par

31. 此外，在严肃的集会中——如同经上所说\BibleQuote{我要在众多的人民中赞美祢}——人们不会从那种口才中获得乐趣，这种口才虽不说虚假之事，却将微小而不重要的真理淹没在浮夸的修饰词海洋中，即使这些修饰词用于装饰伟大而根本的真理，也不会显得优雅或庄重。类似的情形出现在真福\Person{西彼廉}的一封信中，我认为这封信是偶然出现的，或者是故意插入的，以便后代看到基督信仰教育的健全纪律如何治愈了他语言的冗余，并将他限制在更为庄重和谦逊的口才形式中，正如我们在他后来的信件中所见——一种无需费力就能被欣赏、被热切追求、但若不付出巨大努力就无法达到的风格。他在某处说：\Quote{让我们寻求这住所：邻近的荒野提供了一个退隐之处，在那里，葡萄树的枝条蔓延、下垂、缠绕，在支撑的芦苇间攀爬，枝叶的覆盖形成了一个葡萄藤廊。}这里的语言极为流畅和丰富；但它过于华丽，无法令严肃的心灵愉悦。然而，喜爱这种风格的人往往认为，不使用这种风格而采用更朴实风格的人，是因为他们无法达到前者，而不是因为他们判断上懂得避免。因此，这位圣人表明他既能使用那种风格（因为他曾这样做过一次），又选择不使用（因为他再也没有使用过）。\par

% structure: p0049 source=h2 target=subsection reason=relative-heading-rank; base=h2

\subsection{第15章 基督徒教师应在讲道前祈祷}

32. 因此，我们的基督徒演说家，当他说正义、神圣和良善的事（他永不该说别的）时，尽一切努力使人怀着明智、喜悦和服从之心聆听；他无需怀疑，如果他达成这一目标，并在达成之时，他将更多借着祈祷的虔诚而非口才的恩赐成功；因此，他应在开口之前先为自己和他将要讲话的人祈祷。当必须讲话的时刻来到，他应在开口之前，举起干渴的灵魂向天主，汲取他要倾注的东西，并使自己充满他将要分施的东西。因为，关于信德和爱德的每一件事，有许多可说，也有许多说法，除了洞察人心的天主，谁知道在某一时刻我们说什么或被人听到说什么是有益的呢？除了祂——我们和我们的言语都在祂手中——谁能使我们说我们该说的，并以我们该说的方式说呢？因此，渴望知晓又渴望教导的人，应学习一切该教导的事，并获得适合神职人员的言辞能力。但当讲话的时刻来到，让他想起吾主的那句话，更适合虔诚心灵的需要：\BibleQuote{不要思虑如何或说什么；因为在那时刻，必赐给你们当说的话。因为说话的不是你们，而是你们父的圣神在你们内说话。}\BibleRef{玛 10:19}圣神就这样在那些为基督的缘故被交给迫害者的人内说话；为何不也在那些将基督的讯息交付给愿意学习的人内说话呢？\par

% structure: p0051 source=h2 target=subsection reason=relative-heading-rank; base=h2

\subsection{第十六章——人的指引不应被轻视，尽管天主造就真正的导师。}

33. 若有人说，既然圣神使人成为导师，我们便无需指导人如何或教什么，那他同样可以说，既然吾主说\BibleQuote{你们的父在你们求祂以前，已知道你们需要什么}\BibleRef{玛 6:8}，我们便无需祈祷；或说宗徒保禄不应指示弟茂德和弟铎如何或教别人什么。这三封宗徒书信应常摆在教会中每位获得导师职分的人眼前。在《弟茂德前书》中，我们不是读到：\BibleQuote{你要命令和教导这些事？}\BibleRef{弟前 4:11}这些事是什么，先前已经说了。那里不是也写着：\BibleQuote{不要责斥老年人，但要劝他如父？}\BibleRef{弟前 5:1}在《弟茂德后书》中不是说道：\BibleQuote{你要持守你从我这里听到的健全道理的模式？}\BibleRef{弟后 1:13}他岂不该惭愧，正确地讲解真理之言？\BibleRef{弟后 2:15}同处又写道：\BibleQuote{你要宣讲圣言，不论顺境逆境，务必坚持；以极大的忍耐和教训，斥责、警戒、劝勉。}\BibleRef{弟后 4:2}同样，在《弟铎书》中，他不是说主教应当\BibleQuote{持守那合于教训的忠信之言，以便能以健全的道理劝勉并驳斥抗辩的人？}\BibleRef{铎 1:9}那儿也说：\BibleQuote{你该讲说合乎健全道理的事：使老年人节制}等等。\BibleRef{铎 2:1}那儿还说：\BibleQuote{你要讲这些事，以全权劝勉和斥责。不可让人轻视你。提醒他们要服从执政者和掌权者}等等。那么，我们该作何想？宗徒说了人是由圣神的运作而成为导师，但自己又指示他们如何教导和教什么，这是自相矛盾吗？或者，我们该理解为：虽然圣神已赐下，但人教导甚至教导导师的职责并未止息，然而栽种的和浇灌的都算不得什么，只在那使之生长的天主。\BibleRef{格前 3:7}因此，即使有圣人为我们助力，或有圣天使相助，除非天主使人预备好从祂自己学习，没有人能正确领悟关乎与天主同在的生命之事；那位在圣咏中如此呼求的天主：\BibleQuote{教导我承行你的旨意，因为你是我的天主。}因此，同一宗徒对弟茂德自己说——当然是作为导师对门徒说：\BibleQuote{至于你，应持守你所学的和所确信的，因为你知道是从谁学来的。}\BibleRef{弟后 3:14}因为，正如人彼此用在身体上的药物，除非天主赐予效力（祂可以不用药物而治愈，药物却不能没有祂而治愈），否则毫无用处；但人还是施用药物；若出于义务，便被视为慈悲或仁爱的工作；同样，通过人作为工具而施行的教导之助，只有当天主使其产生益处时，才对灵魂有益——祂甚至无需人的帮助或中介就能将福音赐予人。\par

% structure: p0053 source=h2 target=subsection reason=relative-heading-rank; base=h2

\subsection{第十七章——各种演说风格的三种划分。}

34. 因此，凡在讲论中致力于劝善的人，不应轻视那三个目的——教导、悦乐和感动——而应如上所述，祈祷并努力，使人怀着明智、喜悦和顺服之心聆听。当他优雅而恰当地做到这一点时，纵然未能赢得听众的赞同，也完全可以被称为雄辩。因为正是这三个目的——教导、悦乐和感动——那位伟大的罗马修辞学大师似乎旨在以下三条指引来服务：\Quote{因此，他将是这样一位雄辩的人：能以平实的风格说小事，以温和的风格说中等的事，以崇高的风格说伟大的事。}仿佛他也包含了上述三个目的，并用一句话概括了一切：\Quote{因此，他将是这样一位雄辩的人：能以平实的风格说小事，以教导；以温和的风格说中等的事，以悦乐；以崇高的风格说伟大的事，以感动心灵。}\par

% structure: p0055 source=h2 target=subsection reason=relative-heading-rank; base=h2

\subsection{第十八章——基督徒演说家经常处理伟大的事。}

35. 我所引用的作者本可用他本身所定的这三种规则来举例说明有关法律的问题，但他却不能以此来说明有关教会的问题——这正是我希望能加以阐述的此类讲道所唯一关心的问题。因为法律问题中，那些涉及金钱交易的问题被称为小问题；那些涉及人的生命或自由的问题称为大问题。此外，那些既不涉及前者也不涉及后者，其目的不是要让听众去做某事或对某事做出判断，而只是为了愉悦他们的，则介于前两者之间，因此而被称为中等或适中问题。因为适中事物得名于\OriginalTerm{度量}{modus}（即度量）；而把\OriginalTerm{适中}{moderate}当作\Quote{微小}来用，是对这个词的滥用，而非正确的用法。然而，在我们这样的问题中，所有的事情，特别是那些从权威位置对民众讲论的事情，都应关乎人的得救，而且不是他们暂时的得救，而是他们永恒的得救，并且需要防范的是永恒的灭亡，所以我们所说的一切都是重要的；如此重要，以至于传道者关于金钱问题的谈论，无论是关于损失或收益，金额是大是小，都不应被视为不重要。因为正义绝非不重要，即使在最小金钱事务上也必须遵循正义，正如我们的主所说：\BibleQuote{在小事上忠信的，在大事上也忠信}\BibleRef{路 16:10}。那么，最小的事确实很小；但在最小的事上忠信却是伟大的。正如圆的本质，即从圆心到圆周的所有线段都相等，无论在一个大圆盘上还是在一个最小的硬币上都是一样；同样，正义的伟大绝不因所应用的事项微小而有丝毫减弱。\par

36. 当宗徒论及有关世俗事务的审判时（这些不过是金钱问题吗？），他说：\BibleQuote{你们中间有彼此相争的事，怎敢在不义的人面前，而不在圣徒面前求审呢？岂不知圣徒要审判世界吗？若世界为你们所审判，难道你们不配审判这最小的事吗？岂不知我们要审判天使吗？何况今生的事呢？既是这样，你们若有今生的事当审判，是派教会中所轻看的人审判吗？我说这话是要叫你们羞愧。难道你们中间没有一个智慧人能审断弟兄们的事吗？你们竟是弟兄与弟兄告状，而且告在不信主的人面前。你们彼此告状，这已经是你们的大错了。为什么不情愿受欺呢？为什么不情愿吃亏呢？你们倒是欺压人、亏负人，况且所欺压、所亏负的就是弟兄。你们岂不知不义的人不能承受天主的国吗？}\BibleRef{格前 6:1}。为什么宗徒如此愤怒，如此控告、责备、斥责、威胁？为什么他的语调如此频繁而突然的变化，证明他情感的深厚？最后，为什么他对如此琐碎的事情用如此崇高的语调说话？难道世俗的事务值得他如此关注吗？绝对不。不，这一切都是为了正义、爱德和虔诚而做的，这些在每一个清醒的人看来都是伟大的，即使应用于最小的事情。\par

37. 当然，如果我们是在给人建议，告诉他们在教会法庭上应当如何处理世俗案件，不论是为自己还是为关系人，我们会正确地建议他们像处理小事一样心平气和地进行。但我们正在讨论的是那要成为真理的教师、救我们脱离永苦并导入永福的人所当用的说话方式；每当谈论这些真理时，无论是在公开场合还是私下，对一个人还是对许多人，对朋友还是对敌人，在连续的讲道中还是在交谈中，在短论中还是在书籍中，在长信或短信中，都是非常重要的。除非我们准备说，因为一杯凉水是极其微小普通的东西，所以我们的主所说\BibleQuote{凡因门徒的名，只把一杯凉水给这小子里的一个喝的，我实在告诉你们，这人不能不得赏赐}\BibleRef{玛 10:42}是非常微不足道和不重要的。或者，当一位传道者以此话语为经文时，他应该认为他的主题非常不重要，因此说话既无口才也无能力，而是用平淡谦卑的风格。难道不是这样吗？当我们恰巧对民众谈论这个主题，而天主的同在临到我们，使我们所说的话不完全与主题不相称时，从那一杯凉水中会冒出火舌，甚至点燃人们冰冷的心，使他们热切地盼望永恒的赏报而行慈悲之事？\par

% structure: p0059 source=h2 target=subsection reason=relative-heading-rank; base=h2

\subsection{第十九章——基督徒教师必须在不同场合使用不同风格}

38. 然而，我们的教师固然应当谈论重大的事，但他不应总是以庄严的语调谈论它们，而在教导时应以平和的语调，在赞扬或责备时应以温和的语调。但当作某件事时，我们是对那些应当做却不愿做的人说话，那么重大的事就必须以有力的方式、以能影响心灵的方式来说。有时，同一重要的主题在不同时间以所有这些方式来处理：在教导时平静，在强调其重要性时温和，在迫使背离真理的心灵转向并拥抱它时则有力。因为有什么比天主本身更大呢？那么，关于祂就没有什么可学习的吗？教导合一的天主圣三的人，难道应该以冷静讨论之外的方式来谈论它，好使我们关于这个不易理解的题目，能理解我们被赐予的理解吗？在这种情况下，我们要寻求装饰而非证据吗？或者听众应当被推动去做某事，而不是受教导以便学习某事吗？但当我们赞美天主时，无论是赞美祂本身还是祂的作品，那言辞的优美与华丽的领域在人面前是何等开阔！人可以竭尽全力赞美祂，无人能充分赞美祂，然而没有人在某种程度上不赞美祂！但如果天主不被敬拜，或者偶像——无论是魔鬼还是任何受造物——与祂一同或比祂更受敬拜，那么我们就应当以力量和感动说话，表明这是何等大的邪恶，并劝人逃避它。\par

% structure: p0061 source=h2 target=subsection reason=relative-heading-rank; base=h2

\subsection{二十、从圣经中取用各类风格的例子}

39. 但现在让我们谈些更具体的事。我们有圣保禄宗徒的一个平和、低调风格的例子，他说：\Quote{请告诉我，你们这愿意在法律之下的人，不听听法律吗？因为经上记载：亚巴郎有两个儿子：一个由婢女所生，一个由自由妇人所生。那由婢女所生的，是按血肉生的；那由自由妇人所生的，却是藉恩许生的。这都是寓意：它们是两个盟约，一个出自西乃山，生子为奴，即哈加尔。这哈加尔是指亚拉伯的西乃山，与现在的耶路撒冷相符，因为她和她的儿女同为奴隶。但那是天上的耶路撒冷，是自由的，她是我们的母亲：} \BibleRef{迦 4:21}。等等。同样，他这样推理：\Quote{弟兄们，我按人的常理说：即使是人的盟约，如果一经确立，谁也不得废除或增删。许给亚巴郎和他的后裔的恩许就是这样说的。他没有说\InnerQuote{后裔们}，好像指多数，而是指一个，即\InnerQuote{你的后裔}，就是基督。我的意思是：天主先前所立定的盟约，那四百三十年以后才有的法律，不能加以废除，以致取消恩许。因为如果继承出自法律，就不出自恩许；但天主藉恩许赐给了亚巴郎。} \BibleRef{迦 3:15}。由于听者可能会问：如果没有由法律而来的继承，那么法律是为什么而赐下的？他自己预见到这个异议，问道：\Quote{那么法律是为什么而有的？} 回答说：\Quote{它是因为过犯而增加的，直到那蒙恩许的后裔来到；它是由天使藉中保之手所颁布的。但中保不是为一方的中保，天主却是一位。} 这里又出现他本人提出的一个异议：\Quote{那么法律与天主的恩许相悖吗？} 他回答：\Quote{绝对不可！} 他并说明理由：\Quote{因为如果有一种能赐予生命的法律，那么正义就确实出自法律。但圣经把所有的人都圈在罪过之下，为使那藉对耶稣基督的信德而应许的恩赐，赐给相信的人。} \BibleRef{迦 3:19}。因此，教师的责任不但是解释晦涩之处、解开难题，而且同时要回应可能出现的其他问题，以免它们对我们的论述造成怀疑或不信任。然而，如果这些问题一出现就能得到解决，那么扰动那些我们无法消除的东西是无益的。此外，当一个问题引出其他问题，这些又引出更多，如果逐一讨论和解决，推理就会变得极为冗长，以至于除非记忆极其强大活跃，否则推理者无法回到最初的问题。但是，非常可取的是：任何心中可能出现且可能被提出的异议，都应陈述并驳斥，以免将来无人回应时它再次出现，或者以免它在场的人心中想起却默不作声，而从未被彻底消除。\par

40. 在宗徒的以下话语中，我们看到了温和风格：\BibleQuote{不可责斥老年人，但要劝他如劝父亲；对青年人如对弟兄；对老年妇女如对母亲；对青年妇女如对姐妹。} \BibleRef{弟前 5:1} 还有这些：\BibleQuote{弟兄们！我以天主的仁慈请求你们，献上你们的身体当作生活、圣洁和悦乐天主的祭品：这才是你们合理的敬礼。} \BibleRef{罗 12:1} 几乎整篇劝诫段落都属于温和的雄辩风格；其中最美妙的部分是，仿佛相称地支付，那些彼此相关的事物被优雅地结合在一起。例如：\BibleQuote{按照赐给我们的恩宠，我们得到了不同的恩赐：如果是说预言，就要应用与信德相称；如果是服务，就要用在服务上；如果是教导，就要用在教导上；如果是劝勉，就要用在劝勉上；施与的，要大方；监督的，要殷勤；行慈善的，要喜乐。爱不可是虚伪的；恶要厌恶，善要持守。兄弟相爱要彼此亲爱；恭敬人要彼此争先。论殷勤，不可懒惰；论心神，要热切；论盼望，要喜乐；论患难，要忍耐；论祈祷，要恒心；论圣徒的急需，要分担；论款待，要热诚。迫害你们的，要祝福；只可祝福，不可诅咒。与喜乐的一同喜乐，与哭泣的一同哭泣。彼此要同心。} \BibleRef{罗 12:6} 这一切多么优雅地以双成员句结尾：\BibleQuote{不可心高，却要俯就卑微的！} 稍后又说：\BibleQuote{凡你们所当付的，要付清；如粮税，就纳税；如关税，就付关税；如敬畏，就敬畏；如尊敬，就尊敬。} \BibleRef{罗 13:7} 还有一些，虽以单句表达，却以双成员句结束：\BibleQuote{除了彼此相爱外，你们不可再欠人什么。} 再往后：\BibleQuote{黑夜深了，白日已近，所以我们该脱去黑暗的行为，佩戴光明的武器。行动要端庄，好似在白天一样，不可狂宴豪饮，不可淫乱放荡，不可争斗嫉妒；但该穿上主耶稣基督，不要只挂念肉性的事，去满足私欲。} \BibleRef{罗 13:12} 如果这段经文这样翻译：\Quote{\Emph{et carnis providentiam ne in concupiscentiis feceritis}}，无疑耳朵会因更和谐的结尾而愉悦；但我们的译者更严谨，宁可保留词序。这段话在宗徒所说的希腊语中听起来如何，那些更精通希腊语的人可以判断。然而，我的意见是，以相同词序翻译给我们的这段话，即使在原文中也不十分和谐。\par

41. 确实，我必须承认，我们的作者在由和谐结尾构成的言谈优美方面非常欠缺。这究竟是译者的过失，还是如我更愿意相信的，作者有意避免了这种装饰，我不敢断言；因为我承认我不知道。然而，我知道这一点：如果有人精通这种和谐，若要取这些作者的结尾句按和谐法则编排（他很容易通过改换一些同义词，或保留原词但改变排列来实现），他将发现这些天主所感发的人在文法家和修辞学家的学校中被视为重要的那些方面毫无欠缺；他将在他们中发现许多极美的言谈形式——即使在我们的语言中也很美，尤其在原文中更美——那些他们自夸的著作中找不到任何这种形式。但必须注意，在增加和谐的同时，不要削弱这些神圣而权威言语的任何分量。我们的先知们离这种和谐所由充分习得的音乐训练相距如此之远，以至于热罗尼莫，一位非常博学的人，至少描述了其中一些人在希伯来语中使用的韵律；尽管为了准确翻译词语，他在译文中没有保留这些。然而，我（谈论自己的感受，这对我来说比对他人的感受更清楚，也比他人对我的感受更清楚），虽然在我自己的言谈中，无论我自认多么谦逊，我并未忽略这些和谐结尾，但我同样乐意在神圣作者中极少见到它们。\par

42. 豪迈风格的演讲与刚才所说的温和风格之不同，主要在于它并非以言辞的装饰来点缀，而是藉心灵的激情升华为激昂。它确实使用了近乎所有温和风格所用的装饰；但如果这些装饰恰好不在手边，它也不会刻意寻找。因为它被自身的激情所推动；思想的力道，而非对装饰的渴望，使它抓住任何出现在途中的表达之美。对其目标而言，热忱的情感足以暗示合适的词语；它们无需经过精心雕琢的言辞来选择。如果一位勇敢的人以黄金和宝石装饰的武器武装自己，他在战斗的热潮中用这些武器完成英勇事迹，并非因为它们昂贵，而是因为它们是武器；然而，即使愤怒为他提供了一件从地上挖出的武器，他也能造成巨大伤害。宗徒在以下段落中敦促我们，为了福音的职务，并因天主恩宠的安慰所支撑，应当耐心忍受今生的一切苦难。这是一个伟大的主题，并以力量来处理，而且言辞的装饰也不缺少：\Quote{看，}他说，\Quote{如今正是悦纳的时候；看，如今正是救恩的日子。我们在任何事上都不给人绊倒，免得这职务受毁谤；反而处处表明自己是天主的仆役，以极大的忍耐、在患难、匮乏、困苦、鞭打、监禁、暴乱、劳苦、不眠、禁食；以纯洁、以知识、以坚忍、以慈惠、以圣神、以无伪的爱德、以真理之言、以天主的能力、以左右两手的正义武器；藉着荣耀和羞辱、藉着恶名和美名；被视为迷惑人的，却是真诚的；似乎不为人知，却是人所共知的；似乎正在死，看，我们却活着；似乎受惩罚，却没有被杀；似乎忧愁，却常常喜乐；似乎贫穷，却使许多人富足；似乎一无所有，却拥有万有。}\BibleRef{格后 6:2}且看他依然燃烧着：\Quote{啊，格林多人哪！我们的口向你们张开了，我们的心放宽了，}等等；全部引用会过于冗长。\par

43. 同样，在写给罗马人的书信中，他敦促应藉着爱德，在对天主帮助的确信依靠下，战胜这世界的迫害。他以力量与优美处理这个主题：\Quote{我们知道，}他说，\Quote{一切事都协助那些爱天主的人，就是那些按祂旨意蒙召的人，获得益处。因为祂所预知的人，也预定他们与祂儿子的肖像相同，使祂在许多弟兄中作长子。而且祂所预定的人，又召叫了他们；祂所召叫的人，又使他们成义；祂所成义的人，又使他们得光荣。对这一切我们还有什么可说的？如果天主为我们，谁能反对我们呢？祂既然没有怜惜自己的儿子，反而为我们众人把祂交出，怎能不也把万有与祂一同赐给我们呢？谁能控告天主所拣选的人？是那使人成义的天主吗？谁能定他们的罪？是那已死、且更已复活、现今在天主右边、也为我们转求的基督耶稣吗？谁能使我们与基督的爱隔绝？是患难、困苦、迫害、饥饿、赤身、危险、刀剑吗？（正如所记载的：\InnerQuote{为了祢的缘故，我们终日被处死；我们被看作待宰的羊。}）然而，靠着那爱我们的主，我们在这一切事上大获全胜。因为我深信：无论是死、是生、是天使、是率领者、是现时的、是将来的、是有权能者、是崇高、是深遂、或是任何其他受造之物，都不能使我们与天主的爱隔绝，这爱是在我们的主基督耶稣内的。}\BibleRef{罗 8:28}\par

44. 此外，在写给迦拉达人的书信中，虽然整封信是以平缓风格写成，除了结尾处升华为温和的雄辩，但他插入了一段感情如此丰富的话，以致尽管没有上述引文中的任何装饰，却不能不说它是强有力的：\Quote{你们谨守日子、月份、节期、年份。我为你们害怕，惟恐我在你们身上白费了劳力。弟兄们，我恳求你们，要像我一样，因为我也像你们一样；你们丝毫没有亏负过我。你们知道，我最初是因肉体的软弱而向你们传福音。那时你们并没有轻视或厌弃我肉身上的考验，反而接待我如同天主的使者，如同基督耶稣。那么，你们当时的福气在哪里？我给你们作证：如果可能，你们甚至会把你们的眼睛挖出来给我。那么，我因为对你们说实话，就成了你们的仇人吗？那些人热心待你们，但并非好意；他们是要你们疏远我，好使你们热心待他们。但善事上常怀热心，不但在我与你们同在时，就是现在不在你们那里，也是好的。我的孩子们，我为你们再受产痛，直到基督在你们内形成；我渴望现在就在你们那里，并改变我的声音，因为我为你们担心。}\BibleRef{迦 4:10}这里有任何对照词作对偶排列，或词语逐渐升向高潮，或响亮的分句、段落和周期吗？然而，尽管如此，有一股强烈情感的火焰，使我们感受到雄辩的热忱。\par

% structure: p0068 source=h2 target=subsection reason=relative-heading-rank; base=h2

\subsection{第21章 各种风格的范例，取自教会教师，特别是\Person{安波罗削}和\Person{西彼廉}。}

45. 然而宗徒们的这些著作虽则明晰，却仍深奥，且其文笔如此，以致不满足于浅尝辄止、而欲透彻认识之人，不但须阅读聆听，还须借助诠释者。那么，就让我们研究这些不同的言谈方式，一如它们在那些因研读圣经而获得神圣救恩真理、并传授予教会之人著作中的范例。真福西彼廉在其论述圣爵圣事的论文中，采用平实文体。在此书中，他解决了主之圣爵应仅盛水、或水酒相混的问题。但我们必须引述一段文字以资说明。在惯常的引言之后，他进行了对所述问题的讨论。他说：\Quote{请注意，我们受训导，在奉献圣爵时，应持守由主传予我们的惯例，不做任何主未曾先为我们行过之事：因此，为纪念祂而奉献的圣爵应与酒相混。因为，正如基督所说：\BibleQuote{我是真葡萄树}（\BibleRef{若 15:1}），由此可知，基督的血是酒而非水；若缺少酒，圣爵便不能彰显其盛有祂的血——那赎了我们并赐予我们生命的血；因为酒预表基督的血，那在圣经所有预象和宣告中预示并宣告的血。因为我们在《创世纪》中发现，与圣事相关的同一情形已被预示，主的苦难被典型地表达在诺厄身上：当他饮酒而醉，在帐幕内赤身露体，其裸体被次子揭露，而被长子和幼子小心遮盖（\BibleRef{创 9:20}）。无需详述其余情形，只需注意此点：诺厄预表未来的现实，饮的不是水而是酒，从而显示主的苦难。同样，我们看到主晚餐的圣事在默基瑟德司祭身上得到预表，根据圣经的见证，其中说：\BibleQuote{撒冷王默基瑟德带来了饼和酒：他是至高天主的司祭。他祝福了亚巴郎。}（\BibleRef{创 14:18}）现在，默基瑟德是基督的预象，圣神在圣咏中予以宣告，其中圣父对圣子说：\BibleQuote{你按照默基瑟德的品位，永为司祭。}} 在这段文字以及随后整个书信中，平实文体一直保持，读者可以轻易自行证实。\par

46. 圣盎博罗削虽在处理极为重大的问题——圣神与圣父圣子的平等——却也采用平实文体，因为他的目的并不在于辞藻之美，亦非以情感激荡来打动心智，而在于事实与明证。因此，在其著作的引言中，我们看到如下段落，此外还有其他：\Quote{当基德红听到来自天主的消息而惊愕——尽管百姓数千人失败，但天主将藉一人拯救祂的子民脱离仇敌——他牵来一只山羊羔，按天使的指示，连同无酵饼放在磐石上，并将肉汤浇在上面；当天主的杖一端触及磐石时，火从磐石中升起，焚尽了祭品（\BibleRef{民 6:14}）。这迹象似乎表明，磐石是基督身体的预象，因为经上记载：\BibleQuote{他们饮了那跟着他们的神迹之磐石，那磐石就是基督}（\BibleRef{格前 10:4}）。这当然不是指基督的神性，而是指祂的肉躯，其不断涌流的血泉常使渴慕的百姓心满意足。同样，当时以奥秘宣告，主耶稣在被钉十字架时，应在祂的肉躯内消除全世界的罪，不仅包括外在行为的罪责，也包括内心的邪恶欲念。因为羔羊之肉指代外在行为的罪责，肉汤指代内在的私欲诱惑，正如经上记载：\BibleQuote{在他们中间混杂的群众起了贪欲，以色列子民又再次哭泣说：谁给我们肉吃？}（\BibleRef{户 11:4}）。于是，当天使伸出杖触及磐石，火从中升起时，这便是征兆：我们的主之肉躯，充满天主之神，将焚尽人类的一切罪恶。因此，主也说：\BibleQuote{我来是为把火投在地上。}（\BibleRef{路 12:49}）} 并且以同一文体他继续论述主题，主要致力于证明和强调他的论点。\par

47. \Emph{温和的}风格的典范，是\Person{西彼廉}那篇颂赞童贞的著名讲词：\Quote{现在我们的讲论转向童贞者们，她们既获享更高的尊荣，也受到更大的关注。这些是教会之树上的花朵，是属灵恩宠的光荣与装饰，是尊荣与赞美的喜乐，是完美无瑕的工程，是天主回应主的圣德的肖像，是基督羊群中最光辉的部分。她们的母亲教会因她们的丰饶果实而欢欣，因她们而更加繁茂；圣洁的童贞女增加得越多，母亲的喜乐也越增长。在信函的另一处他又说：\BibleQuote{我们既带了那属于土的肖像，也要带那属于天的肖像。}\BibleRef{格前 15:49} 童贞带着这肖像，纯洁带着它，圣洁与真理带着它；那些铭记主的管教、遵守正义与虔诚、在信德上坚强、在敬畏中谦卑、在忍受痛苦中坚忍、在忍受伤害中温良、乐于怜悯、在兄弟般的平安中同心合意的人，都带着它。圣洁的童贞者们，你们应当遵守、珍爱并实践这一切，你们的心为天主和基督而空闲，选择了那更好的一份，引领并指示通向主的道路，你们已向主许下誓愿。年长的，要管理年幼的；年幼的，要服事年长的，并鼓励同辈；以彼此的劝勉互相激励；以美德的榜样激发彼此走向光荣；勇敢地忍受，在属灵上进步，喜乐地完成你们的赛程；只请在我们开始收获尊荣的赏报时记得我们。}\par

48. \Person{盎博罗削}在向已发愿的童贞女提出一个模仿的榜样时，也使用了温和而华丽的风格，他说：\Quote{她是童贞女，不仅身体上，而且心灵上；没有将情感的纯洁与任何伪善的渣滓混在一起；言语庄重，性情明达，寡言少语，喜爱学习；不依靠不可靠的财富，而依靠穷人的祈祷；勤劳工作，言语恭敬；习惯以天主，而非人，作为良心的指引；不伤害任何人，愿所有人好；对长者孝顺，对同辈不嫉妒；避免夸耀，遵循理性，热爱德行。她何时曾以眼神伤害过父母？何时与邻舍争吵？何时鄙视卑微者，嘲笑软弱者，或回避贫困者？她只习惯走访那些怜悯不会脸红、端庄不会经过的场所。她眼中没有傲慢，言语没有鲁莽，姿态没有放荡；她的举止不纵欲，步态不轻浮，声音不蛮横；因此她的外在形象是她心灵的形象，是纯洁的图画。因为一所好房子在门槛处就应让人认出，并在入口处就显示出里面没有黑暗的角落，正如里面点着的灯的光照亮外面。我何需详述她在饮食上的节制，在责任上的丰盈——前者低于本性的需求，后者超越本性的能力？后者没有间断的间隙，前者以禁食加倍日子；当恢复精力的渴望兴起时，她以满足生命的食物，而不助长食欲。} 现在我已引用了后面这些段落作为温和风格的例子，因为它们的意图不是诱使那些尚未献身的人发守贞愿，而是表明那些发了愿的人应当具有怎样的品格。要说服任何人迈出如此本性且如此重要的一步，需要以庄严的风格激动和点燃心灵。然而，殉道者\Person{西彼廉}并没有写关于承担守贞职分的问题，而是关于童贞女的服饰和举止。但那位伟大的主教甚至在这些方面也以庄严的雄辩之力敦促她们尽职。\par

49. 但我要从他们两人都触及的一个主题中选取庄严风格的例子。两人都谴责那些用脂粉涂饰，或者更确切地说，毁坏自己面容的女子。第一位在论述这一主题时说：\Quote{假设一位画家用与自然媲美的色彩描绘了某个人的面貌、身形和肤色，当肖像以精湛的艺术完成后，另一位画家将手覆于其上，好像要以他的更高超技艺改进已完成的画作；那么第一位艺术家必定会深感侮辱，他的愤怒也会正当地被激起。那么，你以为你竟能如此胆大妄为的恶行而无罪脱身吗？如此侮辱了伟大的艺术家天主？因为，即使你承认自己在人前并非不端正，你的心也没有被这些娼妓般的骗术所玷污，然而，你在败坏和侵犯天主的所有物时，你证明自己比淫妇更坏。你自以为这种技艺使自己得到装饰和美化，这本身就是对天主作品的指责，对真理的违背。请听宗徒的警告之声：\BibleQuote{你们应把旧酵母除净，好使你们成为新和的面团，正如你们原是无酵的一样。因为我们的逾越节羔羊基督，已被祭杀，作了牺牲。所以，我们过节，不可用旧酵母，也不可用奸诈和邪恶的酵母，而只用纯洁和真诚的无酵饼。}\BibleRef{格前 5:7} 如今，真诚和真理如何能继续存在呢？当真诚被污染，真理被娼妓般的着色和骗术的谎言所改变时？你的主说：\BibleQuote{你不能使一根头发变白或变黑。}\BibleRef{玛 5:36} 而你竟想拥有更大的权能，以致废弃你主的话语？你以鲁莽和亵渎的手改变你头发的颜色；但愿你能以预言的目光展望未来，把你的头发染成火焰的颜色。} 若要将后面的话都引出来，那就太长了。\par

50. 盎博罗削再次谴责此类行为时说道：\Quote{由此产生了这些趋向邪恶的诱因：妇女因害怕不能吸引男子，便用经过精心挑选的颜色涂饰面容，于是从面容上的污点发展到贞洁上的污点。改变自然的面貌为绘画的面貌是多么愚蠢，为了担心招致丈夫的不悦，竟公开宣告她们已经不悦于自己！因为那想要改变自己天然容貌的妇女，是在宣判自己的罪过；她急切地取悦别人，恰恰证明她首先已不悦于自己。喔，妇人，还有什么比你自己更确凿地证明你的丑陋呢？当你害怕显露自己时，你若美丽，为何隐藏你的美丽？你若平凡，为何假扮美丽？既然你在自己的意识中或在他人的意识中都享受不到谎言的乐趣。因为他爱的是另一个女人，你想要取悦另一个男人；你因他爱别人而生气，而他却是从你这里学会了通奸。你是自身伤害的邪恶教唆者。就连一个曾被拉皮条者坑害的妇女，也羞于充当拉皮条的角色，尽管她卑劣，但她作的恶是针对自己，而非他人。通奸之罪几乎比你的罪更可容忍；因为通奸败坏的是廉耻，而你败坏的是自然。} 这已足够清楚，我认为，此雄辩之辞强烈地呼吁妇女避免用欺诈手段修饰外貌，并培养谦逊与敬畏。因此我们注意到，其风格既非平实，亦非温和，而是通篇庄严。现在，在这两位我选作其他作者代表的作家中，以及在其他阐述真理且言辞优美的教会作家——即他们用明智、犀利、富有美感和力量的方式表述真理——，在其各类著作和讲道中，都可以找到这三种风格的许多范例；勤勉的学生通过孜孜不倦的阅读，并辅以自身实践，便能充分领会所有这些风格。\par

% structure: p0075 source=h2 target=subsection reason=relative-heading-rank; base=h2

\subsection{第22章——风格多样化的必要性。}

51. 但我们不能认为混合这些不同的风格是违反规则的；相反，只要符合审美趣味，就应当引入各种风格的变化。因为当我们单调地保持一种风格时，便不能保持听众的注意力；但从一种风格过渡到另一种风格时，讲演会更优雅地展开，即使它延伸到更长的篇幅。每一种风格自身也有其变化，可以防止听众的注意力冷却或变得倦怠。不过，与庄严风格相比，平实风格可以更长久地坚持而不需变化。因为为了打动听众的情感而必须激起的心理激动，一旦被充分激发，被提升到的音调越高，能够维持的时间就越短。因此，我们必须警惕，在试图将已激起的情感推向更高点时，反而失去我们已经取得的效果。但在插入了需要用较平静风格处理的内容之后，我们可以有效地回归到必须用强力处理的内容，从而使雄辩的潮汐如同大海般涨落。由此可知，庄严风格若想持续长久，就不应单调不变，而应间歇地与其他风格交替出现；然而整个讲演或作品应参照那种占主导地位的风格。\par

% structure: p0077 source=h2 target=subsection reason=relative-heading-rank; base=h2

\subsection{第23章——如何混合各种风格。}

52. 现在，确定哪些风格应相互交替，以及在哪些地方有必要使用某种特定风格，是一件重要的事。例如，在庄严风格中，几乎总是或经常希望开头是温和的。演讲者可以自行决定，在允许使用庄严风格的地方，甚至使用平实风格，以便庄严风格通过对比显得更为庄重，仿佛从暗色背景中闪耀出更明亮的光芒。再者，无论讲演或作品的风格如何，当遇到棘手的难题需要解决时，就需要精确的区分，而这自然要求平实风格。因此，每当这类问题出现时，就必须用平实风格与其他两种风格交替使用；同样，无论论述的总体基调如何，当需要赞美或责备，而不涉及对任何人的谴责或开释，或不需获取任何人对某个行动路线的同意时，我们就必须使用温和风格。因此，在庄严风格中，也在平实风格中，偶尔会用到其他两种风格。另一方面，温和风格虽不总是，但偶尔也需要平实风格；例如，当出现一个棘手的问题需要解决，或当某些可加修饰的点被保留而不加修饰地以平实风格表达时，以便给某些（可称之为）修饰的繁复之处带来更大的效果。但温和风格从不需要庄严风格的帮助；因为它的目标是愉悦，而非激动心灵。\par

% structure: p0079 source=h2 target=subsection reason=relative-heading-rank; base=h2

\subsection{第24章——庄严风格所产生的效果。}

53. 若随后有频繁而热烈的掌声，我们不应因此以为讲者是在采用雄辩风格；因为这种效果常由平和风格的精准分辨，以及温文风格的优美而产生。雄辩风格则反之，常以其感染力使听众静默，却引出他们的泪水。例如，当我在\Place{毛里塔尼亚的凯撒勒雅}劝阻民众脱离那种他们称之为\OriginalTerm{卡泰尔瓦}{Caterva}的内战（甚或比内战更糟）时（因为不仅是同胞，而且是邻居、兄弟、甚至父与子，分成两派，持石武装，每年在特定季节连续数日互相厮杀，每人都杀死自己能杀之人），我竭尽全力以雄辩言辞，试图根除并驱离如此残忍且根深蒂固的邪恶；然而，并非当我听到他们的掌声，而是当我看见他们的泪水时，我才认为产生了效果。因为掌声表明他们得到教导和喜悦，而泪水表明他们被降服。当我看见他们的泪水时，甚至在结果证实之前，我就确信这可怕而野蛮的习俗（该习俗由他们的父辈和久远的祖先流传下来，像仇敌一样围攻着他们的心，更确切地说，完全占据了他们的心）已被推翻；立刻在讲道结束后，我以心与声呼吁他们赞美感谢天主。看，靠基督的降福，至今八年或更久，当地再未尝试类似事情。除此之外，在许多其他情形中，我观察到人们显示智者雄辩对他们的影响，不是通过喧闹的掌声，而是通过呻吟，有时甚至通过泪水，最终通过生活的改变。\par

54. 平和风格也在许多人身上造成改变；但这改变是为了教导他们未知之事，或说服他们相信原以为不可信之事，而非促使他们去做他们明知该做却不愿做的事。要打破这类刚硬，言辞需要强劲。同样，当赞美与谴责被雄辩地表达时，即使在温文风格中，也会对一些人有如此影响：他们不仅对赞美与谴责的雄辩感到喜悦，而且被引导过一种配得赞美的生活，避免那种招致责备的生活。但没有人会说，所有如此喜悦的人因此改变了习性，然而所有被雄辩风格所感动的人都会付诸行动，所有被平和风格教导的人都会知道或相信一个他们先前所不知道的真理。\par

% structure: p0082 source=h2 target=subsection reason=relative-heading-rank; base=h2

\subsection{第25章 温文风格当如何使用}

55. 由以上一切我们可以得出结论：最后两种风格所达到的目的，正是那些渴望以智慧与雄辩讲论的人最必须确保的。另一方面，温文风格正当地追求的目标，即通过表达之美使人喜悦，本身并非充分的目的；但当我们要说的是善而有益之理，且听众已熟悉并乐意接受，以至于无需教导或说服他们时，文体之美可以发挥作用，促使他们更快地顺从，或更坚定地持守。因为一切雄辩的功能，无论采取这三种形式中的哪一种，都是要说服性地说话，其目的是说服；一个雄辩的人无论采用何种风格，都会说服性地说话；但除非他成功说服，否则他的雄辩便未达到目的。在平和风格中，他说服听众相信他所言为真；在雄辩风格中，他说服他们去做他们明知该做却未做的事；在温文风格中，他说服他们相信他的言辞优雅华美。但达到类似最后一种的目的又有何用？那些自诩雄辩、以颂词及类似表演为荣的人或许渴望如此；这些表演的目的不是教导听众，也不是说服他们采取任何行动，而仅仅是为了愉悦他们。然而，我们应当使这个目的从属于另一个目的，即通过这种雄辩风格实现我们在使用雄辩风格时所追求的效果。因为我们可以通过使用这种风格说服人们养成好习惯、放弃坏习惯，如果他们尚未顽固到需要强劲风格的程度；或者，如果他们已开始步入善途，我们可以促使他们更热切地追求，并持之以恒。因此，即使在温文风格中，我们也必须使用表达之美，不是为了炫耀，而是为了智慧的目的；不满足于仅仅取悦听众，而是旨在帮助他追求我们在他面前提出的善的目标。\par

% structure: p0084 source=h2 target=subsection reason=relative-heading-rank; base=h2

\subsection{第26章 在任何风格中，演说者都应追求明晰、优美与说服力}

55. 现在关于我刚才提出的三个条件，任何希望有智慧且雄辩地说话的人都必须满足，即：清晰、风格优美和说服力，我们不应理解为这三个品质分别对应三种演说风格，一种风格一个品质，以致清晰是\OriginalTerm{平实风格}{genus submissum}的独特优点，优美是\OriginalTerm{温和风格}{genus temperatum}，而说服力是\OriginalTerm{庄严风格}{genus grande}。相反，所有演说，无论其风格如何，都应当始终追求并尽可能展示所有这三个优点。因为即使是平实风格，我们也不希望让听者感到乏味；因此我们希望听众不仅理解我们，而且也能愉悦地聆听。再者，为什么我们用神圣的证词来强化我们所教导的，除非我们希望带着听众一起走，也就是说，通过呼求那位的帮助来迫使听众同意，关于那位曾说：\Quote{你的证词是极其可靠的}？当有人用平实风格叙述一个故事时，他除了希望被相信，还希望什么？但如果他没有通过某种风格的优美来吸引注意力，谁会听他呢？如果他不清晰，显然既不能带来愉悦也不能强制说服。此外，平实风格以其赤裸的简洁，当它解开非常困难的问题并投下意想不到的光明；当它从毫无预料的地方挖掘并带来一些非常敏锐的观察；当它抓住并揭露一个反对意见的虚假性，这个意见最初看起来似乎坚不可摧；尤其是当这一切伴随着自然的、不刻意追求的表达之美，以及不显眼地强制性的节奏和风格的平衡，似乎是由主题的性质所召唤的：这种风格这样使用时，常常会引起如此大的掌声，以至于人们几乎不能相信它是平实风格。因为它毫无修饰或防御地出现，并以赤裸的简洁进行战斗，这并不妨碍它用神经和肌肉的重量粉碎对手，并仅凭自身右臂的力量压倒并摧毁反对它的谎言。如何解释那些如此说话的人经常获得的热情和强烈的掌声，若不是因为真理如此不可抗拒地确立并如此胜利地辩护所带来的自然愉悦？因此，基督教教师和演说家在使用平实风格时，不仅应努力做到清晰和易懂，还应带来愉悦并让听众信服。\par

57. 温和风格的雄辩，在基督教演说家手中，也必须既不完全缺乏装饰，也不适合装饰，也不以愉悦为唯一目的（这在他人手中是它所宣称的全部成就）；而是在其赞美和谴责中，应旨在引导听者追求、避免或弃绝它所谴责的。另一方面，没有清晰，这种风格就无法带来愉悦。因此，这三个品质——清晰、优美和说服力——也应在这种风格中寻求；当然，优美是其主要目标。\par

58. 再者，当必须用庄严风格来激动和感动听众的心灵时（当听众承认你所说的是真实和愉快的，却不愿据此行动时，总是需要这样），你当然必须用庄严风格说话。但如果一个人不理解所说的话，怎么会被感动呢？如果他不得到愉悦，谁会留下听呢？因此，在这种风格中，当要说服一颗顽固的心服从时，你必须说话清晰而愉悦，如果你希望被顺从的心听到。\par

% structure: p0088 source=h2 target=subsection reason=relative-heading-rank; base=h2

\subsection{第27章——生活与教导和谐一致的人将更有效地教导。}

59. 但无论风格的庄严如何，说话者的生活对于确保听众的服从更为重要。那说话明智而雄辩，但生活邪恶的人，确实可以教导许多渴望学习的人；尽管，如经上所记，他\BibleQuote{对自己没有益处。}\BibleRef{德 37:19}因此，宗徒也说：\BibleQuote{或假意或真诚，无论怎样，基督终被宣讲。}\BibleRef{斐 1:18}基督就是真理；然而我们看到真理可以被宣讲，尽管不是以真理的方式——也就是说，一个内心乖谬、欺骗的人可以宣讲正确且真实的东西。因此，耶稣基督被那些寻求自己而非耶稣基督之事的人所宣讲。但由于真正的信徒服从的并非任何人的声音，而是主自己的声音，祂说：\BibleQuote{凡他们所吩咐你们的，你们都要谨守遵行；但不要效法他们的行为；因为他们能说不能行。}\BibleRef{玛 23:3}因此，那些自己过着无益生活的人，却被他人在有益中听到。因为虽然他们寻求自己的目标，但他们不敢教导自己的教义，因为他们坐在教会权威的高位上，这权威是建立在纯正教义之上的。因此，我们主自己在说刚才关于这类人的话之前，说了这样一句话：\BibleQuote{经师和法利塞人坐在梅瑟的座位上。}\BibleRef{玛 23:2}那么，他们所坐的座位，不是他们自己的而是梅瑟的，迫使他们说出好的话，尽管他们行恶。因此，他们在生活中遵循自己的道路，但被属于他人的座位所阻止，不敢宣讲自己的教义。\par

60. 如今，这些人宣讲自己所不实践的事，却造福了许多人；但如果他们言行一致，将会造福更多人。因为有许多人拿导师的教训与其行为相比较，为自己邪恶的生活寻找借口，他们心中甚至口头说：为什么你命令我做的，你自己却不去做？因此，他们不再顺从地聆听那不听从他自身的人，在轻视宣讲者的同时，也学会了轻视所宣讲的话语。为此，宗徒写信给弟茂德，在告诉他\Quote{不要让人轻视你的青年时代}之后，立即加上他避免轻视的方法：\Quote{却要在言语、行为、爱德、心神、信德、洁德上，做信徒的榜样。} \BibleRef{弟前 4:12}\par

% structure: p0091 source=h2 target=subsection reason=relative-heading-rank; base=h2

\subsection{第二十八章——真理比表达更重要；何谓关于言辞的争论。}

61. 如上所述的导师，为了使人服从，不仅可以安静温和地说话，甚至可以激昂慷慨地说话，而不会失于谦逊，因为他的生活保护他不受轻视。因为当他追求正直生活时，他也留意保持良好的名声，在天主和众人眼前成就光明的事，\BibleRef{格后 8:21}敬畏天主，关心众人。甚至在言谈中，他也宁愿以内容而不以言辞取悦于人；他认为，事物在事实上越真实，就说得越好，并且导师应该驾驭言辞，而不应让言辞驾驭他。这正是宗徒所说：\Quote{并不是用智慧的言辞，免得基督的十字架失去效力。} \BibleRef{格前 1:17}同样，他对弟茂德所说的话也有同样的意义：\Quote{你在天主面前嘱咐他们，不可在言辞上争辩，因为那毫无益处，只能败坏听众。} \BibleRef{弟后 2:14}这并不意味着，当仇敌反对真理时，我们不应为真理辩护。因为，宗徒描述主教应有的品格时所说的话又在哪里：\Quote{能以健全的道理，劝勉并驳斥反对的人？} \BibleRef{铎 1:9}在言辞上争辩，并不是注意用真理克服错误的方法，而是关心自己的表达方式应该比别人更受偏爱。那不在言辞上争辩的人，无论他说话安静、温和，还是激昂，使用言辞的唯一目的就是使真理变得明白、悦人、有效；因为甚至连爱德本身——它是诫命的目的和法律的满全——也不能正确地实行，除非爱的对象是真实的而非虚假的。因为正如一个人身体俊美而心灵丑恶，比他的身体也畸形更令人痛苦，所以那些教导谎言的人，如果碰巧口才流利，就更可悲。因此，既雄辩又智慧地说话，就是用恰当合适的言辞来表达应当教导的真理——这些言辞在平实的风格中是充分的，在温和的风格中是优雅的，在庄严的风格中是有力的。但是，一个人如果既不能雄辩又智慧地说话，就应该智慧而不雄辩地说话，胜过雄辩而不智慧。\par

% structure: p0093 source=h2 target=subsection reason=relative-heading-rank; base=h2

\subsection{第二十九章——宣讲者可以将比自己更有口才之人所写的作品向民众宣讲。}

但如果他连这一点也做不到，那就让他的生活不仅为自己赢得赏报，也为他人提供榜样；让他生活方式本身成为一篇雄辩的讲道。\par

63. 确实，有些人口才很好，却写不出任何讲稿。这样的人若将别人用智慧和修辞写成的讲稿记在心里，向百姓宣讲，只要他们不存欺骗之心，就不能责备他们。因为这样，许多人就成了真理的宣讲者（这当然是可取的），而并非许多都是教师；因为所有人宣讲的都是同一位真正教师所撰写的讲辞，他们之间并无分裂。这样的人也不应被耶肋米亚先知的话吓倒，天主藉先知斥责那些\BibleQuote{各邻舍偷窃我言语}的人。\BibleRef{耶 23:30} 因为偷窃的人拿走的不是属于自己的东西；但天主的言语属于所有服从它的人。而那讲得好却生活败坏的人，才是真正拿走了属于别人的言语。因为他所说的善言似乎出于自己的思想，却与他的生活毫无共同之处。因此，天主说那些偷窃他言语的人，是那些藉着说天主的言语而显得好，实际上却是坏的，因为他们随从自己的道路。如果仔细审视这事，其实并不是他们自己在说那些善言。因为他们怎能用言语说出他们用行为否认的事呢？宗徒论及这样的人时并非无故说：\BibleQuote{他们声称认识天主，但在行为上却否认他。}\BibleRef{铎 1:16} 从某种意义上说，他们确实说那些事；从另一种意义上说，他们并不说那些事；因为这两句话都必须是真的，都是出自那位真理者之口。论及这样的人，他有一处说：\BibleQuote{凡他们所吩咐你们的，你们都要遵守实行；但不要照他们的行为去做}——这就是说，你们从他们口里听到的，你们去做；你们在他们生活中看到的，你们不要做——\BibleQuote{因为他们能说不能行。}\BibleRef{玛 23:3} 因此，他们虽然不行，却仍说。但在另一处，他斥责这样的人说：\BibleQuote{毒蛇的种类，你们既是恶的，怎能说出善来呢？}\BibleRef{玛 12:34} 由此看来，即使他们所说的善言，也不是他们自己说的，因为他们在意志和行为上否认了他们所说的。因此，一个邪恶而能言善辩的人可以写一篇宣讲真理的讲辞，由一位不善言辞的善人来讲；当这事发生时，前者从自己身上取出不属于自己的东西，后者从别人那里接受实际上属于自己的东西。但当真正的信徒为真正的信徒提供这种服务时，双方所说的都是自己的，因为天主是他们的，他们所讲的一切都属于天主；甚至那些不能撰写自己所讲内容的人，通过使自己的生活与那讲辞相协调，而使之成为自己的。\par

% structure: p0096 source=h2 target=subsection reason=relative-heading-rank; base=h2

\subsection{第三十章 讲道者应以祈祷向天主开始其宣讲。}

63. 但无论是将要向百姓说话，还是将要口述让别人向百姓宣读或宣讲，他都应祈求天主赐给他合宜的言语。因为如果艾斯德尔王后为着民族现世的幸福而向国王说话时，曾祈求天主将合宜的话放在她口中，那么，为着人类的永恒幸福而从事言语和教理工作的人，岂不更应祈求同样的恩惠？那些将要宣讲别人为他们撰写的讲辞的人，在接受讲辞之前，应为撰写的人祈祷；在接受之后，应祈求自己讲得好，也祈求听众能留心听；当讲辞顺利结束时，应感谢天主，因为他们知道这样的恩惠是从他而来的，\BibleQuote{在他的手中，有我们和我们的言语。}\BibleRef{智 7:16}\par

% structure: p0098 source=h2 target=subsection reason=relative-heading-rank; base=h2

\subsection{第三十一章 为本书长度致歉。}

64. 本书比我预期或希望的篇幅要长。但乐于阅读或聆听的人不会觉得它长。觉得它长却急于了解其内容的人可以部分阅读。不愿了解的人也不必抱怨它的长度。我感谢天主，以我仅有的一点能力，在这四卷书中努力描绘的，不是我自己是怎样的人（因为我的缺点很多），而是那愿意在健全的道理（即基督徒教义）中劳苦的人应当是怎样的人，不仅为自身教导，也为教导他人。\par

\SourceRecord{Translated by James Shaw.；Nicene and Post-Nicene Fathers, First Series；Vol. 2.；Edited by Philip Schaff.；Buffalo, NY: Christian Literature Publishing Co.,；1887.；<http://www.newadvent.org/fathers/12024.htm>.}
